% Dissertation produire le plan détaillé d’une dissertation — Français Tle Tech — Cours signé OnBuch+
% Programme officiel MINESEC. Compiler avec : tectonic N07-dissertation-produire-plan-detaille-dissertation.tex
\newcommand{\DOCMATIERE}{Français}
\newcommand{\DOCNIVEAU}{Tle Tech}
\newcommand{\DOCMODULE}{Français – classe de Terminale technique}
\newcommand{\DOCLECON}{N07}
\newcommand{\DOCTITRE}{Dissertation produire le plan détaillé d’une dissertation}
% OnBuch+ — préambule « cours signés » ENRICHI (compilé avec tectonic / XeLaTeX)
% Système visuel « Pop » d'OnBuch+ : fond crème (jamais blanc), bordures encre
% épaisses, ombres « dures » décalées, titres Archivo Black en capitales.
% Version MATHÉMATIQUES : graphes (pgfplots/tikz), boîtes pédagogiques
% (définition, propriété, méthode, attention, le savais-tu, exercice résolu,
% exercice, TP, bilan), page de garde et signature OnBuch+ ancrée en bas.
%
% Macros attendues dans main.tex AVANT \input{preamble} :
%   \DOCMATIERE  \DOCNIVEAU  \DOCMODULE  \DOCLECON (numéro)  \DOCTITRE
\documentclass[11pt]{article}

\usepackage{fontspec}
\usepackage[french]{babel}
\usepackage[a4paper,margin=2cm,top=3.4cm,bottom=2.8cm,headheight=1.4cm,footskip=1.3cm]{geometry}
\usepackage{amsmath,amssymb,amsfonts}
\usepackage{mathtools} % \xmapsto, \coloneqq... souvent utilisés par les modèles
\usepackage{cancel} % simplifications barrées (bilans de forces)
% \abs{z} (module, valeur absolue) et \norm{u} (norme), utilisés par les modèles
\providecommand{\abs}{}\let\abs\relax\DeclarePairedDelimiter{\abs}{\lvert}{\rvert}
\providecommand{\norm}{}\let\norm\relax\DeclarePairedDelimiter{\norm}{\lVert}{\rVert}
\usepackage[table]{xcolor} % [table] : fournit aussi \rowcolors (alternance de lignes)
\usepackage{pagecolor}
\usepackage{tikz}
\usetikzlibrary{arrows.meta,positioning,calc,shapes.geometric,shapes.misc,shapes.arrows,decorations.pathmorphing,decorations.pathreplacing,decorations.markings,patterns,shadows,mindmap,trees,fit,backgrounds,matrix,intersections,angles,quotes}
\usepackage{pgfplots}
\usepackage[version=4]{mhchem} % équations nucléaires \ce{^{235}_{92}U}
\usepackage[european,siunitx]{circuitikz} % schémas électriques
\pgfplotsset{compat=1.18}
\usepgfplotslibrary{fillbetween,groupplots} % aires entre courbes (calcul intégral)
\usepackage{enumitem}
\usepackage{fancyhdr}
% [most] charge toutes les bibliothèques tcolorbox (listings, minted, raster,
% documentation...) et gonfle inutilement la mémoire TeX sur un document long
% (~300 boîtes) : on ne charge que ce qui est réellement utilisé.
\usepackage[skins,breakable]{tcolorbox}
\usepackage{graphicx}
\usepackage{colortbl}
\usepackage{booktabs}
\usepackage{tabularx}
\usepackage{multirow}
\usepackage{diagbox} % case d'en-tête diagonale des tableaux à double entrée
% Colonnes extensibles centrée / gauche / droite (C, L, R), souvent utilisées par les modèles
\newcolumntype{C}{>{\centering\arraybackslash}X}
\newcolumntype{L}{>{\raggedright\arraybackslash}X}
\newcolumntype{R}{>{\raggedleft\arraybackslash}X}
\usepackage{array}
\usepackage{multicol}
\usepackage{siunitx}
% parse-numbers=false : accepte aussi \SI{2,5 \times 10^{-3}}{...}
\sisetup{locale=FR,output-decimal-marker={,},group-minimum-digits=4,parse-numbers=false}
\DeclareSIUnit\ohm{\ensuremath{\symup{\Omega}}} % U+2126 absent de la police
\DeclareSIUnit\division{div} % réglages d'oscilloscope (ms/div, V/div)
\DeclareSIUnit\tr{tr} % tours (vitesse de rotation, ex. \SI{1500}{\tr\per\minute})
\DeclareSIUnit\molar{M} % concentration molaire (ex. \SI{3}{\molar}, \SI{1}{\micro\molar})
\DeclareSIUnit\an{an} % année (le modèle confond parfois avec \year, un compteur TeX) — ex. \SI{5}{\kilo\meter\per\an}
\DeclareSIUnit\FCFA{FCFA} % franc CFA (ex. \SI{500}{\FCFA\per\kilo\gram})
\DeclareSIUnit\degreeBrix{\textdegree Brix} % teneur en sucre d'un sirop/jus (réfractométrie)
\DeclareSIUnit\month{mois} % le modèle utilise parfois \month, un compteur TeX, comme unité de durée
\DeclareSIUnit\microliter{\micro\liter} % le modèle écrit parfois \microliter en un seul token
\DeclareSIUnit\million{millions} % dénombrement (ex. \SI{15}{\million\per\mL} spermatozoïdes)
\usepackage{qrcode}
% \degree : commande fréquemment utilisée par les modèles pour un angle ou une
% température en dehors de \SI{}{} (non fournie par les paquets chargés).
\providecommand{\degree}{^{\circ}}
% \qed (fin de démonstration), utilisé par les modèles sans amsthm
\providecommand{\qed}{\hfill$\square$}
% \barre : double barre des valeurs interdites dans un tableau de variations (tkz-tab)
\providecommand{\barre}{\ensuremath{\|}}
% environnement proof (amsthm non chargé) utilisé par les modèles
\ifdefined\proof\else\newenvironment{proof}[1][Démonstration]{\par\noindent\textit{#1.}\ }{\qed\par}\fi
\usepackage{lastpage}
\usepackage{hyperref}
\hypersetup{hidelinks}

% ============================================================
% Palette « Pop » OnBuch+ — mêmes hex que la classe P (plus_ui.dart)
% ============================================================
\definecolor{poporange}{HTML}{F59321}
\definecolor{popdark}{HTML}{E07A0C}
\definecolor{popcream}{HTML}{FFF6E8}
\definecolor{popink}{HTML}{1C1714}
\definecolor{poppurple}{HTML}{7A5AE0}
\definecolor{popgreen}{HTML}{2E9E6A}
\definecolor{poppink}{HTML}{E0457B}
\definecolor{popblue}{HTML}{2F7DE1}
\definecolor{popgold}{HTML}{C98A12}
\definecolor{popmuted}{HTML}{8A7F76}
\definecolor{popline}{HTML}{EADFD2}
\definecolor{popgray}{HTML}{8A7F76}
\colorlet{popgrey}{popgray}
% Alias tolérants (le modèle nomme parfois les couleurs différemment)
\colorlet{popwhite}{white}
\colorlet{popblack}{popink}
\colorlet{popred}{poppink}
\colorlet{popyellow}{popgold}
% Teintes claires (fonds de tableaux/encadrés) — le modèle nomme parfois
% les couleurs avec un suffixe L (ex. popblueL) sans les avoir définies.
\colorlet{poporangeL}{poporange!20}
\colorlet{popdarkL}{popdark!20}
\colorlet{poppurpleL}{poppurple!20}
\colorlet{popgreenL}{popgreen!20}
\colorlet{poppinkL}{poppink!20}
\colorlet{popblueL}{popblue!20}
\colorlet{popgoldL}{popgold!20}
\colorlet{popredL}{popred!20}
\colorlet{popgrayL}{popgray!20}
% Teintes claires pour fonds de tableaux / zones
\colorlet{popblueL}{popblue!10!white}
\colorlet{poporangeL}{poporange!14!white}
\colorlet{popgreenL}{popgreen!12!white}
\colorlet{poppurpleL}{poppurple!12!white}
\colorlet{poppinkL}{poppink!12!white}
\colorlet{popgoldL}{popgold!14!white}

\pagecolor{popcream}

% ============================================================
% Typographie
% ============================================================
\defaultfontfeatures{Ligatures=TeX}
\setmainfont{PlusJakartaSans-Medium.ttf}[Path=../../../fonts/,BoldFont=PlusJakartaSans-Bold.ttf]
% Maths en sans-serif (Fira Math) avec chiffres et lettres droites de Plus
% Jakarta, pour que les formules se fondent dans le texte.
\usepackage[math-style=ISO,bold-style=ISO]{unicode-math}
\setmathfont{FiraMath-Regular.otf}[Path=../../../fonts/]
\setmathfont{PlusJakartaSans-Medium.ttf}[Path=../../../fonts/,range={up/num,up/{Latin,latin},\mathup}]
\usepackage{icomma} % virgule décimale sans espace en mode math
% Caractères mathématiques Unicode absents des polices de texte (Plus Jakarta,
% Archivo Black) : rendus par leur équivalent mathématique, en texte comme en
% formule — sinon ils s'affichent en carré vide (ex. « Arithmétique dans ℤ »).
\usepackage{newunicodechar}
\newunicodechar{ℝ}{\ensuremath{\mathbb{R}}}
\newunicodechar{ℤ}{\ensuremath{\mathbb{Z}}}
\newunicodechar{ℂ}{\ensuremath{\mathbb{C}}}
\newunicodechar{ℕ}{\ensuremath{\mathbb{N}}}
\newunicodechar{ℚ}{\ensuremath{\mathbb{Q}}}
\newunicodechar{𝔻}{\ensuremath{\mathbb{D}}}
\newunicodechar{α}{\ensuremath{\alpha}}
\newunicodechar{β}{\ensuremath{\beta}}
\newunicodechar{γ}{\ensuremath{\gamma}}
\newunicodechar{δ}{\ensuremath{\delta}}
\newunicodechar{ε}{\ensuremath{\varepsilon}}
\newunicodechar{θ}{\ensuremath{\theta}}
\newunicodechar{λ}{\ensuremath{\lambda}}
\newunicodechar{μ}{\ensuremath{\mu}}
\newunicodechar{σ}{\ensuremath{\sigma}}
\newunicodechar{τ}{\ensuremath{\tau}}
\newunicodechar{φ}{\ensuremath{\varphi}}
\newunicodechar{ψ}{\ensuremath{\psi}}
\newunicodechar{ω}{\ensuremath{\omega}}
\newunicodechar{Σ}{\ensuremath{\Sigma}}
% majuscules produites par \MakeUppercase dans les titres (α-aminés -> Α)
\newunicodechar{Α}{\ensuremath{\alpha}}
\newunicodechar{Β}{\ensuremath{\beta}}
\newunicodechar{Γ}{\ensuremath{\Gamma}}
\newunicodechar{Δ}{\ensuremath{\Delta}}
\newunicodechar{Ε}{\ensuremath{\varepsilon}}
\newunicodechar{Θ}{\ensuremath{\Theta}}
\newunicodechar{Λ}{\ensuremath{\Lambda}}
\newunicodechar{Μ}{\ensuremath{\mu}}
\newunicodechar{Τ}{\ensuremath{\tau}}
\newunicodechar{Φ}{\ensuremath{\Phi}}
\newunicodechar{Ψ}{\ensuremath{\Psi}}
\newunicodechar{Ω}{\ensuremath{\Omega}}
\newunicodechar{↦}{\ensuremath{\mapsto}}
\newunicodechar{∈}{\ensuremath{\in}}
\newunicodechar{∉}{\ensuremath{\notin}}
\newunicodechar{∧}{\ensuremath{\wedge}}
\newunicodechar{⇔}{\ensuremath{\Leftrightarrow}}
\newunicodechar{⟺}{\ensuremath{\Longleftrightarrow}}
\newunicodechar{⇒}{\ensuremath{\Rightarrow}}
\newunicodechar{⟹}{\ensuremath{\Longrightarrow}}
\newunicodechar{∘}{\ensuremath{\circ}}
\newunicodechar{∪}{\ensuremath{\cup}}
\newunicodechar{∩}{\ensuremath{\cap}}
\newunicodechar{≡}{\ensuremath{\equiv}}
\newunicodechar{⊂}{\ensuremath{\subset}}
\newunicodechar{⊥}{\ensuremath{\perp}}
\newunicodechar{∃}{\ensuremath{\exists}}
\newunicodechar{∀}{\ensuremath{\forall}}
\newunicodechar{∛}{\ensuremath{\sqrt[3]{\,}}}
\newunicodechar{∜}{\ensuremath{\sqrt[4]{\,}}}
\newunicodechar{⃗}{\ensuremath{{}^{\to}}}
\newunicodechar{ⁿ}{\ifmmode^{n}\else\textsuperscript{\ensuremath{n}}\fi}
\newunicodechar{ˣ}{\ifmmode^{x}\else\textsuperscript{\ensuremath{x}}\fi}
\newunicodechar{ᵇ}{\ifmmode^{b}\else\textsuperscript{\ensuremath{b}}\fi}
\newunicodechar{ʳ}{\ifmmode^{r}\else\textsuperscript{\ensuremath{r}}\fi}
\newunicodechar{ᵘ}{\ifmmode^{u}\else\textsuperscript{\ensuremath{u}}\fi}
\newunicodechar{ʸ}{\ifmmode^{y}\else\textsuperscript{\ensuremath{y}}\fi}
\newunicodechar{ᵃ}{\ifmmode^{a}\else\textsuperscript{\ensuremath{a}}\fi}
\newunicodechar{ᵉ}{\ifmmode^{e}\else\textsuperscript{\ensuremath{e}}\fi}
\newunicodechar{ᵏ}{\ifmmode^{k}\else\textsuperscript{\ensuremath{k}}\fi}
\newunicodechar{ᵖ}{\ifmmode^{p}\else\textsuperscript{\ensuremath{p}}\fi}
\newunicodechar{ᴺ}{\ifmmode^{N}\else\textsuperscript{\ensuremath{N}}\fi}
\newunicodechar{ᶜ}{\ifmmode^{c}\else\textsuperscript{\ensuremath{c}}\fi}
\newunicodechar{ʰ}{\ifmmode^{h}\else\textsuperscript{\ensuremath{h}}\fi}
\newunicodechar{ᵠ}{\ifmmode^{q}\else\textsuperscript{\ensuremath{q}}\fi}
\newunicodechar{ₙ}{\ifmmode_{n}\else\textsubscript{\ensuremath{n}}\fi}
\newunicodechar{ₐ}{\ifmmode_{a}\else\textsubscript{\ensuremath{a}}\fi}
\newunicodechar{ᵢ}{\ifmmode_{i}\else\textsubscript{\ensuremath{i}}\fi}
\newunicodechar{ᵣ}{\ifmmode_{r}\else\textsubscript{\ensuremath{r}}\fi}
\newunicodechar{ₖ}{\ifmmode_{k}\else\textsubscript{\ensuremath{k}}\fi}
\newunicodechar{ᵦ}{\ifmmode_{\beta}\else\textsubscript{\ensuremath{\beta}}\fi}
\newunicodechar{ₓ}{\ifmmode_{x}\else\textsubscript{\ensuremath{x}}\fi}
\newfontfamily\popdisplay{ArchivoBlack-Regular.ttf}[Path=../../../fonts/]
\newfontfamily\poplogo{SpaceGrotesk-Bold.ttf}[Path=../../../fonts/]
\newfontfamily\popsemi{PlusJakartaSans-SemiBold.ttf}[Path=../../../fonts/]
\newfontfamily\popxbold{PlusJakartaSans-ExtraBold.ttf}[Path=../../../fonts/]
\newfontfamily\popxboldls{PlusJakartaSans-ExtraBold.ttf}[Path=../../../fonts/,LetterSpace=3.0]

\setlist[enumerate]{leftmargin=1.7em,itemsep=5pt,topsep=4pt}
\setlist[itemize]{leftmargin=1.7em,itemsep=4pt,topsep=4pt,label={\color{poporange}\textbullet}}
\setlength{\parindent}{0pt}
\setlength{\parskip}{5pt}
\renewcommand{\arraystretch}{1.25}

% pgfplots : style Pop par défaut
\pgfplotsset{
  every axis/.append style={axis line style={popink,line width=1pt},tick style={popink},
    grid style={popline,line width=0.5pt},label style={font=\small},tick label style={font=\footnotesize}},
  popcurve/.style={poporange,line width=1.8pt,no markers,smooth},
}
% Même style disponible dans un \draw TikZ simple (hors pgfplots)
\tikzset{popcurve/.style={poporange,line width=1.8pt}}
\tikzset{popgrid/.style={popline,very thin}}
\colorlet{popcurve}{poporange} % le nom du style est parfois utilisé comme couleur

% ============================================================
% Pill / squiggle
% ============================================================
\newcommand{\poppill}[3]{%
  \tikz[baseline=-0.55ex]\node[draw=popink,line width=1.2pt,rounded corners=2.6pt,%
    fill=#2,inner xsep=6.5pt,inner ysep=3pt]{%
    \popxboldls\fontsize{7}{7}\selectfont\color{#3}\MakeUppercase{#1}};%
}
\newcommand{\popsquiggle}[1]{%
  \tikz[baseline=-0.4ex]{
    \draw[#1,line width=2.6pt,line cap=round]
      plot [smooth] coordinates {(0,0.09) (0.34,0.01) (0.68,0.21) (1.02,0.01) (1.36,0.10)};
  }%
}
% Mot-clé surligné (marqueur orange sous le texte)
\newcommand{\cle}[1]{{\popxbold\color{popdark}#1}}

% ============================================================
% En-tête / pied
% ============================================================
\pagestyle{fancy}
\fancyhf{}
\renewcommand{\headrulewidth}{0pt}
\renewcommand{\footrulewidth}{0pt}
\lhead{\raisebox{-0.15ex}{\poplogo\fontsize{13}{13}\selectfont\color{popink}OnBuch\color{poporange}+}}
\rhead{\poppill{\DOCMATIERE\ \textbullet\ \DOCNIVEAU\ \textbullet\ Leçon \DOCLECON}{white}{popink}}
\lfoot{\popxbold\fontsize{7.6}{7.6}\selectfont\color{popmuted}\MakeUppercase{Cours signé OnBuch+}\ \textbullet\ {\popsemi\DOCTITRE}}
\rfoot{\tikz[baseline=-0.6ex]\node[rounded corners=8.5pt,fill=popink,minimum height=17pt,minimum width=17pt,inner xsep=5pt,inner sep=0]{\popxbold\fontsize{8}{8}\selectfont\color{popcream}\thepage\,/\,\pageref*{LastPage}};}

% ============================================================
% Titres
% ============================================================
\newcommand{\chaptitre}[2]{%
  \begin{tcolorbox}[enhanced,colback=poporange,colframe=popink,boxrule=2.6pt,arc=11pt,
      drop shadow=popink,left=15pt,right=15pt,top=12pt,bottom=11pt,before skip=3pt,after skip=18pt]
    \poppill{#2}{popink}{popcream}\par\vspace{9pt}
    {\raggedright\hyphenpenalty=10000\exhyphenpenalty=10000\popdisplay\fontsize{22}{24}\selectfont\color{popcream}\MakeUppercase{#1}\par}
  \end{tcolorbox}
}
% Section numérotée : pastille orange avec le numéro + titre Archivo
\newcounter{popsec}
\newcommand{\coursec}[1]{%
  \stepcounter{popsec}\setcounter{popsub}{0}%
  \par\vspace{16pt}\noindent
  \begin{minipage}{\linewidth}
  \tikz[baseline=-0.7ex]\node[draw=popink,line width=1.6pt,fill=poporange,rounded corners=4pt,minimum size=22pt,inner sep=0]{\popdisplay\fontsize{12}{12}\selectfont\color{popcream}\thepopsec};%
  \hspace{8pt}{\popdisplay\fontsize{15}{17}\selectfont\color{popink}\MakeUppercase{#1}}\par
  \vspace{4pt}\noindent{\color{poporange}\rule{36pt}{3.6pt}}
  \end{minipage}\par\nopagebreak\vspace{8pt}
}
\newcounter{popsub}
\newcommand{\courssub}[1]{%
  \stepcounter{popsub}%
  \par\vspace{9pt}\noindent{\popxbold\fontsize{11.5}{13}\selectfont\color{popink}{\color{poporange}\thepopsec.\thepopsub}\hspace{6pt}#1}\par\nopagebreak\vspace{4pt}
}

% ============================================================
% Boîtes pédagogiques — même recette : fond blanc, bordure épaisse colorée,
% ombre dure encre, badge plein. Titre optionnel [#1].
% ============================================================
\tcbset{popbase/.style={breakable,enhanced,colback=white,boxrule=2.3pt,arc=9pt,
  shadow={3pt}{-3pt}{0pt}{popink},left=14pt,right=14pt,top=11pt,bottom=11pt,before skip=11pt,after skip=15pt,
  fontupper=\popsemi\fontsize{10.6}{14.5}\selectfont\color{popink},
  fontlower=\popsemi\fontsize{10.6}{14.5}\selectfont\color{popink}}}
% Coche dessinée (le glyphe ✓ n'existe pas dans les polices du document)
\newcommand{\popcheck}[1]{\tikz[baseline=-0.5ex]\draw[#1,line width=1.6pt,line cap=round,line join=round](-0.13,0.0)--(-0.04,-0.09)--(0.13,0.1);}
\providecommand{\coche}{\popcheck{popgreen}}
\providecommand{\resultat}[1]{\par\smallskip\noindent\fcolorbox{popgreen}{popgreenL}{\parbox{\dimexpr\linewidth-2\fboxsep-2\fboxrule\relax}{\textbf{Résultat.} #1}}\par\smallskip}
\newcommand{\popboxhead}[3]{\poppill{#1}{#2}{white}\ifx\relax#3\relax\else\hspace{6pt}{\popxbold\fontsize{10.6}{12}\selectfont\color{popink}#3}\fi\par\vspace{7pt}}

\newtcolorbox{aretenir}[1][]{popbase,colframe=popgreen,before upper={\popboxhead{À retenir}{popgreen}{#1}}}
\newtcolorbox{exemplebox}[1][]{popbase,colframe=poporange,before upper={\popboxhead{Exemple}{poporange}{#1}}}
\newtcolorbox{definition}[1][]{popbase,colframe=popblue,before upper={\popboxhead{Définition}{popblue}{#1}}}
\newtcolorbox{propriete}[1][]{popbase,colframe=poppurple,before upper={\popboxhead{Propriété}{poppurple}{#1}}}
\newtcolorbox{methode}[1][]{popbase,colframe=popgold,before upper={\popboxhead{Méthode}{popgold}{#1}}}
\newtcolorbox{attention}[1][]{popbase,colframe=poppink,before upper={\popboxhead{Attention}{poppink}{#1}}}
\newtcolorbox{experience}[1][]{popbase,colframe=popblue,colback=popblueL,before upper={\popboxhead{Expérience}{popblue}{#1}}}
% « Le savais-tu ? » : carte violette pleine (contexte, culture, Cameroun)
\newtcolorbox{savaistu}[1][]{popbase,colback=poppurple,colframe=popink,
  fontupper=\popsemi\fontsize{10.4}{14.2}\selectfont\color{white},
  before upper={\poppill{Le savais-tu\,?}{white}{poppurple}\ifx\relax#1\relax\else\hspace{6pt}{\popxbold\color{white}#1}\fi\par\vspace{7pt}}}
% Exercice résolu : énoncé en haut, \tcblower puis la solution sur fond crème
\newcounter{popexo}\newcounter{popexr}
\newtcolorbox{exoresolu}[1][]{popbase,colframe=poporange,colbacklower=popcream,
  segmentation style={popink,line width=1.2pt,dashed},
  before upper={\stepcounter{popexr}\popboxhead{Exercice résolu \thepopexr}{poporange}{#1}},
  before lower={\poppill{Solution}{popink}{popcream}\par\vspace{6pt}}}
% Exercice (non résolu en place) — la correction est dans la partie Corrigés
\newtcolorbox{exercice}[1][]{popbase,colframe=popink,
  before upper={\stepcounter{popexo}\popboxhead{Exercice \thepopexo}{popink}{#1}}}
% Corrigé
\newtcolorbox{corrige}[1][]{popbase,colframe=popgreen,colback=popgreenL,
  before upper={\popboxhead{Corrigé}{popgreen}{#1}}}
% Objectifs / prérequis / bilan
\newtcolorbox{objectifs}{popbase,colframe=popink,colback=poporangeL,
  before upper={\popboxhead{Objectifs de la leçon}{popink}{}}}
\newtcolorbox{prerequis}{popbase,colframe=popink,colback=popblueL,
  before upper={\popboxhead{Prérequis}{popblue}{}}}
\newtcolorbox{bilan}{popbase,colframe=popink,colback=white,boxrule=2.8pt,
  before upper={\popboxhead{Fiche bilan}{poporange}{L'essentiel à connaître pour l'examen}}}
% Figure encadrée avec légende
\newtcolorbox{popfigure}[1][]{enhanced,breakable=false,colback=white,colframe=popline,boxrule=1.4pt,arc=8pt,
  left=10pt,right=10pt,top=10pt,bottom=8pt,before skip=10pt,after skip=12pt,halign=center,
  lower separated=false,halign lower=center,
  fontlower=\popsemi\fontsize{9}{11.5}\selectfont\color{popmuted},
  #1}
\newcommand{\legende}[1]{\par\vspace{6pt}{\popsemi\fontsize{9}{11.5}\selectfont\color{popmuted}#1}}

% ============================================================
% Signature OnBuch+
% ============================================================
\newcommand{\poplogoinline}[1]{{\poplogo\fontsize{#1}{#1}\selectfont\color{popink}OnBuch\color{poporange}+}}
\newcommand{\signatureonbuch}{%
  \begin{tcolorbox}[enhanced,colback=popink,colframe=popink,boxrule=2.2pt,arc=10pt,
      drop shadow=poporange,left=14pt,right=14pt,top=10pt,bottom=10pt,before skip=0pt,after skip=0pt]
    \tikz[baseline=-0.7ex]\node[circle,fill=popgreen,draw=popcream,line width=1.3pt,minimum size=22pt,inner sep=0]{\popcheck{white}};%
    \hspace{9pt}\begin{minipage}[c]{0.62\linewidth}
      {\popxbold\fontsize{12}{13}\selectfont\color{popcream}Signé\ \textbullet\ {\poplogo OnBuch\color{poporange}+}}\par\vspace{2pt}
      {\raggedright\popsemi\fontsize{8}{10.5}\selectfont\color{popline}\MakeUppercase{Équipe pédagogique OnBuch+}\\\MakeUppercase{Conforme au programme officiel MINESEC}\par}
    \end{minipage}\hfill
    {\poplogo\fontsize{22}{22}\selectfont\color{popcream}OnBuch\color{poporange}+}
  \end{tcolorbox}%
}

% ============================================================
% Page de garde
% #1 titre · #2 matière · #3 niveau · #4 module · #5 n° leçon · #6 durée ·
% #7 description / objectifs (texte ou itemize)
% ============================================================
\newcommand{\pagegarde}[7]{%
  \thispagestyle{empty}
  \newgeometry{a4paper,margin=1.7cm,top=1.6cm,bottom=1.4cm}%
  \begin{tikzpicture}[remember picture,overlay]
    \fill[poporange!18] ([xshift=-3.2cm,yshift=-3.6cm]current page.north east) circle (4.6cm);
    \fill[poppurple!14] ([xshift=2.4cm,yshift=7.6cm]current page.south west) circle (3.8cm);
  \end{tikzpicture}%
  {\poplogo\fontsize{30}{30}\selectfont\color{popink}OnBuch\color{poporange}+}\hfill
  \raisebox{0.6ex}{\poppill{Cours signé \textbullet\ Premium}{popink}{popcream}}\par
  \vspace{1pt}\popsquiggle{poppurple}\par
  \vspace{8pt}
  \begin{tcolorbox}[enhanced,colback=poporange,colframe=popink,boxrule=2.8pt,arc=13pt,
      drop shadow=popink,left=18pt,right=18pt,top=12pt,bottom=13pt,after skip=10pt]
    \poppill{#2 \textbullet\ #3}{popink}{popcream}\hspace{5pt}\poppill{#4}{white}{popink}\par\vspace{8pt}
    {\popdisplay\fontsize{14}{14}\selectfont\color{popink}LEÇON #5}\par\vspace{4pt}
    {\raggedright\hyphenpenalty=10000\exhyphenpenalty=10000\popdisplay\fontsize{25}{27}\selectfont\color{popcream}\MakeUppercase{#1}\par}
    \vspace{8pt}
    {\popxbold\fontsize{9.5}{9.5}\selectfont\color{popink}\MakeUppercase{Durée conseillée : #6}}
  \end{tcolorbox}
  \begin{tcolorbox}[enhanced,colback=white,colframe=popink,boxrule=2.2pt,arc=10pt,
      drop shadow=popink,left=16pt,right=16pt,top=10pt,bottom=10pt,after skip=8pt]
    \poppill{Dans cette leçon}{poppurple}{white}\par\vspace{5pt}
    \popsemi\fontsize{9.3}{12.6}\selectfont\color{popink}#7
  \end{tcolorbox}
  \noindent\begin{minipage}[t]{0.487\linewidth}
    \begin{tcolorbox}[enhanced,colback=popgreen,colframe=popink,boxrule=2.2pt,arc=10pt,
        drop shadow=popink,left=11pt,right=11pt,top=10pt,bottom=10pt,height=3.1cm,valign=center]
      \tcbox[colback=white,colframe=popink,boxrule=1.3pt,arc=5pt,boxsep=0pt,nobeforeafter,
        left=3pt,right=3pt,top=3pt,bottom=3pt,valign=center]{\qrcode[height=1.9cm]{https://play.google.com/store/apps/details?id=com.luvvix.onbuch&hl=fr}}%
      \hspace{8pt}\begin{minipage}[c]{0.5\linewidth}
        {\popdisplay\fontsize{11}{12.5}\selectfont\color{white}\MakeUppercase{Télécharge OnBuch}}\par\vspace{4pt}
        {\raggedright\popsemi\fontsize{8.4}{11}\selectfont\color{white}Gratuit \textbullet\ Google Play\\Tuteur IA, annales, concours, résultats\par}
      \end{minipage}
    \end{tcolorbox}
  \end{minipage}\hfill
  \begin{minipage}[t]{0.487\linewidth}
    \begin{tcolorbox}[enhanced,colback=poppurple,colframe=popink,boxrule=2.2pt,arc=10pt,
        drop shadow=popink,left=11pt,right=11pt,top=10pt,bottom=10pt,height=3.1cm,valign=center]
      \tikz[baseline=-0.7ex]\node[circle,fill=white,draw=popink,line width=1.3pt,minimum size=1.6cm,inner sep=0]{\popdisplay\fontsize{24}{24}\selectfont\color{poppurple}+};%
      \hspace{8pt}\begin{minipage}[c]{0.6\linewidth}
        {\popdisplay\fontsize{11}{12.5}\selectfont\color{white}\MakeUppercase{Passe à OnBuch+}}\par\vspace{4pt}
        {\raggedright\popsemi\fontsize{8.4}{11}\selectfont\color{white}Tous les cours signés, Tuteur IA illimité, fiches PDF — onglet OnBuch+ de l'app\par}
      \end{minipage}
    \end{tcolorbox}
  \end{minipage}
  \vspace{8pt}
  \popcontenu
  \vfill
  \signatureonbuch
  \restoregeometry
  \newpage
}

% Rangée « Ce cours contient » (page de garde)
\newcommand{\poptile}[3]{% #1 couleur #2 gros texte #3 légende
  \begin{tcolorbox}[enhanced,colback=#1,colframe=popink,boxrule=2pt,arc=9pt,shadow={2.5pt}{-2.5pt}{0pt}{popink},
      left=6pt,right=6pt,top=9pt,bottom=9pt,halign=center,valign=center,height=1.75cm,width=\linewidth,nobeforeafter]
    {\popdisplay\fontsize{12.5}{13}\selectfont\color{white}\MakeUppercase{#2}}\par\vspace{3pt}
    {\centering\popsemi\fontsize{7.8}{9.5}\selectfont\color{white}#3\par}
  \end{tcolorbox}}
\newcommand{\popcontenu}{%
  \par\noindent{\popxboldls\fontsize{7.5}{7.5}\selectfont\color{popmuted}CE COURS SIGNÉ CONTIENT}\par\vspace{7pt}
  \noindent\begin{minipage}[t]{0.235\linewidth}\poptile{popblue}{Cours}{complet, illustré, pas à pas}\end{minipage}\hfill
  \begin{minipage}[t]{0.235\linewidth}\poptile{poporange}{Méthodes}{et exercices résolus}\end{minipage}\hfill
  \begin{minipage}[t]{0.235\linewidth}\poptile{poppink}{Exercices}{type examen + corrigés}\end{minipage}\hfill
  \begin{minipage}[t]{0.235\linewidth}\poptile{popgreen}{Fiche bilan}{l'essentiel à retenir}\end{minipage}\par}

% Sommaire
\newtcolorbox{sommaire}{enhanced,colback=white,colframe=popink,boxrule=2.2pt,arc=10pt,
  drop shadow=popink,left=18pt,right=18pt,top=13pt,bottom=13pt,before skip=6pt,after skip=20pt,
  before upper={\poppill{Sommaire}{poppurple}{white}\par\vspace{9pt}\popsemi\fontsize{10.6}{14.5}\selectfont\color{popink}}}

% Signature de fin de cours — poussée tout en bas de la dernière page
\newcommand{\findecours}{%
  \par\vfill\nopagebreak
  \begin{center}\popsquiggle{poporange}\end{center}\vspace{4pt}
  \signatureonbuch
}
\AtBeginDocument{\def\llbracket{\mathopen{[\mkern-3.5mu[}}\def\rrbracket{\mathclose{]\mkern-3.5mu]}}} % intervalles d'entiers (glyphe ⟦ absent de la police)
% Glyphes absents de Fira Math : redéfinis à partir de glyphes disponibles
\AtBeginDocument{%
  \def\setminus{\mathbin{\backslash}}\let\smallsetminus\setminus
  \def\vdots{\mathord{\vbox{\baselineskip3.2pt\lineskiplimit0pt\kern4pt\hbox{.}\hbox{.}\hbox{.}}}}%
  \def\ddots{\mathinner{\mkern1mu\raise6.5pt\hbox{.}\mkern2mu\raise3.6pt\hbox{.}\mkern2mu\raise0.7pt\hbox{.}\mkern1mu}}%
  \def\triangle{\mathord{\scalebox{0.9}{$\symup{\Delta}$}}}%
  \def\nabla{\mathord{\raisebox{\depth}{\scalebox{1}[-1]{$\symup{\Delta}$}}}}%
  \def\checkmark{\popcheck{popdark}}%
  \def\star{\mathbin{\ast}}\def\bigstar{\mathord{\ast}}%
  \def\diamond{\mathbin{\scalebox{0.6}{$\Diamond$}}}%
  \def\bigcup{\mathop{\mathchoice{\raisebox{-0.3em}{\scalebox{1.7}{$\cup$}}}{\scalebox{1.25}{$\cup$}}{\cup}{\cup}}\displaylimits}%
  \def\bigcap{\mathop{\mathchoice{\raisebox{-0.3em}{\scalebox{1.7}{$\cap$}}}{\scalebox{1.25}{$\cap$}}{\cap}{\cap}}\displaylimits}%
  \def\lesseqqgtr{\mathrel{\lessgtr}}\def\bigoplus{\mathop{\oplus}\displaylimits}%
}
\providecommand{\ssi}{si et seulement si} % abréviation fréquente
\AtBeginDocument{\providecommand{\wideparen}{\overparen}} % arc de cercle

%% FIGFIX-BEGIN v5 — correctifs globaux des figures (docs/onbuch-generation/tools/figfix)
\usepackage{adjustbox}\usepackage{etoolbox}\tcbuselibrary{raster}
% toute figure TikZ/pgfplots plus large que la ligne est réduite à la largeur disponible
\BeforeBeginEnvironment{tikzpicture}{\begin{adjustbox}{max width=\linewidth}}
\AfterEndEnvironment{tikzpicture}{\end{adjustbox}}
% légendes pgfplots lisibles par-dessus les courbes
\pgfplotsset{every axis/.append style={legend style={fill=white,fill opacity=0.92,text opacity=1,draw=black!15,inner sep=3pt}}}
% étiquettes d'arêtes (syntaxe edge["..."]) avec fond blanc
\tikzset{every edge quotes/.append style={fill=white,inner sep=1.2pt}}
%% FIGFIX-END
\begin{document}

\pagegarde{Dissertation produire le plan détaillé d’une dissertation}{Français}{Tle Tech}{Français – classe de Terminale technique}{N07}{16 séances (fiche de progression harmonisée)}{Cette leçon outille l'élève de Terminale technique pour analyser un sujet de dissertation, formuler une problématique pertinente, mobiliser ses ressources lexicales (synonymie, antonymie, hypéronymie) et construire un plan détaillé rigoureux. Elle s'appuie sur la lecture des œuvres au programme (Le Vieux nègre et la médaille, Le Lion et la perle) et sur des sujets authentiques du Baccalauréat technique camerounais.
\vspace{4pt}
\begin{multicols}{2}\small
\begin{itemize}
\item À la fin de cette leçon, je sais identifier les paratextes d'une œuvre et en tirer des hypothèses de lecture
\item À la fin de cette leçon, je sais analyser un sujet de dissertation (mots-clés, type de sujet, consigne implicite)
\item À la fin de cette leçon, je sais formuler une problématique claire à partir d'un sujet
\item À la fin de cette leçon, je sais distinguer le référent de ses substituts dans un texte pour assurer la cohérence de mon écriture
\item À la fin de cette leçon, je sais utiliser la synonymie, l'antonymie, l'hypéronymie et l'hyponymie pour enrichir et nuancer mon argumentation
\item À la fin de cette leçon, je sais repérer les caractéristiques d'un texte argumentatif et les contenus latents d'un énoncé
\end{itemize}\end{multicols}}

\begin{sommaire}
\begin{multicols}{2}
\begin{enumerate}
\item Réception des textes : paratextes et projet d'étude
\item Le référent et ses substituts
\item Lexique de l'argumentation : synonymie, antonymie, hypéronymie
\item Le texte argumentatif et l'analyse du sujet
\item Recherche des idées et élaboration du plan détaillé
\item Compte rendu, remédiation et entraînement au Bac
\item Activité d'intégration
\item Méthodes et exercices résolus
\item Exercices
\item Corrigés des exercices
\item Fiche bilan
\end{enumerate}
\end{multicols}
\end{sommaire}

\begin{objectifs}
\begin{itemize}
\item À la fin de cette leçon, je sais identifier les paratextes d'une œuvre et en tirer des hypothèses de lecture
\item À la fin de cette leçon, je sais analyser un sujet de dissertation (mots-clés, type de sujet, consigne implicite)
\item À la fin de cette leçon, je sais formuler une problématique claire à partir d'un sujet
\item À la fin de cette leçon, je sais distinguer le référent de ses substituts dans un texte pour assurer la cohérence de mon écriture
\item À la fin de cette leçon, je sais utiliser la synonymie, l'antonymie, l'hypéronymie et l'hyponymie pour enrichir et nuancer mon argumentation
\item À la fin de cette leçon, je sais repérer les caractéristiques d'un texte argumentatif et les contenus latents d'un énoncé
\item À la fin de cette leçon, je sais rechercher, trier et organiser des idées en plan détaillé (thèse, antithèse, synthèse ou plan thématique)
\item À la fin de cette leçon, je sais confronter, évaluer et améliorer ma production et celle de mes camarades
\item À la fin de cette leçon, je sais rédiger un compte rendu de lecture et justifier ma démarche
\end{itemize}
\end{objectifs}

\begin{prerequis}
\begin{itemize}
\item Notions de paragraphe argumenté et de texte argumentatif vues en Première
\item Connecteurs logiques et organisation d'un paragraphe (classe de Seconde)
\item Lecture méthodique d'un texte littéraire (incipit, narrateur, personnages)
\item Types de phrases et expansion du nom (révision utile)
\end{itemize}
\end{prerequis}

\newpage

\coursec{Réception des textes : paratextes et projet d'étude}

\textbf{Situation d'accroche.} À l'occasion du concours d'entrée à l'ENAM, un candidat de Bafoussam échoue non par manque d'idées, mais parce que son plan ne répond pas au sujet posé : il a tout dit sur la corruption alors qu'on lui demandait si l'État seul peut l'éradiquer. Au Cameroun, réussir le Bac technique ou un concours administratif exige de savoir transformer un sujet en problématique claire, puis en plan détaillé cohérent. Mais avant même de produire, il faut savoir \emph{recevoir} un texte : c'est en lisant intelligemment des œuvres argumentatives que l'on apprend à construire sa propre argumentation. Comment passer méthodiquement d'un sujet de dissertation à un plan détaillé convaincant ? C'est tout l'enjeu de cette leçon, dont la première étape consiste à apprendre à aborder une œuvre littéraire par ses \cle{paratextes} et à organiser sa lecture en \cle{projet d'étude}.

\courssub{Le paratexte : définition et composantes}

Avant d'ouvrir un roman, le lecteur rencontre déjà beaucoup de choses : un titre, un nom d'auteur, une image de couverture, quelques lignes au dos du livre. Ces éléments ne font pas partie du récit lui-même, et pourtant ils orientent puissamment notre lecture. C'est l'intuition fondatrice : \emph{on ne lit jamais un texte nu, on lit toujours un texte habillé}.

\begin{definition}[Le paratexte]
Le \cle{paratexte} est l'ensemble des éléments qui entourent le texte et orientent sa réception : le \textbf{titre} et le \textbf{sous-titre}, le \textbf{nom de l'auteur}, la \textbf{couverture} (image, maquette, nom de l'éditeur, collection), la \textbf{quatrième de couverture} (résumé, extraits de presse), la \textbf{préface}, la \textbf{dédicace}, ainsi que le sommaire ou la table des matières. Le terme, forgé par le critique Gérard Genette (\emph{Seuils}, 1987), signifie littéralement « ce qui est à côté du texte ».
\end{definition}

Distinguons deux espaces du paratexte :
\begin{itemize}
\item le \cle{péritexte}, situé \emph{dans} le livre lui-même : titre, préface, dédicace, notes, table des matières ;
\item l'\cle{épitexte}, situé \emph{hors} du livre : interviews de l'auteur, articles de presse, affiches publicitaires.
\end{itemize}

Chaque élément paratextuel remplit une fonction précise, qu'il faut savoir identifier à l'examen :

\begin{center}
\begin{tabularx}{\linewidth}{|l|X|X|}
\hline
\rowcolor{popblueL} \textbf{Élément} & \textbf{Description} & \textbf{Fonction principale} \\
\hline
Titre & Premier mot lu ; souvent énigmatique & Attirer, annoncer le thème ou le héros \\
\hline
Nom de l'auteur & Identité, nationalité, réputation & Situer l'œuvre (littérature africaine, époque) \\
\hline
Couverture & Image, couleurs, graphisme & Suggérer l'ambiance, le cadre, le ton \\
\hline
Quatrième de couverture & Résumé partiel, citations élogieuses & Donner envie sans tout dévoiler \\
\hline
Préface & Texte liminaire de l'auteur ou d'un tiers & Éclairer les intentions, le contexte \\
\hline
Dédicace & Hommage bref à une personne & Révéler l'engagement de l'auteur \\
\hline
\end{tabularx}
\end{center}

\begin{popfigure}
\begin{tikzpicture}
\node[draw=popdark, fill=popcream, rounded corners, minimum width=3.2cm, minimum height=4.2cm, line width=1pt] (livre) at (0,0) {};
\node[font=\bfseries\small, text=popdark] at (0,0.6) {LE VIEUX NÈGRE};
\node[font=\bfseries\small, text=popdark] at (0,0.2) {ET LA MÉDAILLE};
\node[font=\small\itshape, text=popmuted] at (0,-0.6) {Ferdinand OYONO};
\node[draw=popblue, fill=popblueL, rounded corners, font=\small, align=center] (titre) at (-4.2,2.2) {Titre et sous-titre\\(annoncent le héros\\et le symbole)};
\node[draw=poporange, fill=poporangeL, rounded corners, font=\small, align=center] (auteur) at (4.2,2.2) {Nom de l'auteur\\(écrivain camerounais,\\génération 1956)};
\node[draw=popgreen, fill=popgreenL, rounded corners, font=\small, align=center] (image) at (-4.2,-2.2) {Image de couverture\\(vieil homme, décor\\villageois, uniforme)};
\node[draw=poppurple, fill=poppurpleL, rounded corners, font=\small, align=center] (resume) at (4.2,-2.2) {Quatrième de couverture\\(résumé partiel, ton\\satirique annoncé)};
\node[draw=popgold, fill=popgoldL, rounded corners, font=\small, align=center] (editeur) at (0,-3.4) {Éditeur et collection (10/18, Julliard)};
\draw[-{Stealth}, popblue, thick] (titre) -- (livre.north west);
\draw[-{Stealth}, poporange, thick] (auteur) -- (livre.north east);
\draw[-{Stealth}, popgreen, thick] (image) -- (livre.south west);
\draw[-{Stealth}, poppurple, thick] (resume) -- (livre.south east);
\draw[-{Stealth}, popgold, thick] (editeur) -- (livre.south);
\end{tikzpicture}
\legende{Carte d'identité paratextuelle d'un roman : chaque élément qui entoure le texte (1. titre, 2. auteur, 3. image, 4. quatrième de couverture, 5. éditeur) fournit un indice de lecture avant même l'ouverture du livre.}
\end{popfigure}

\courssub{Étude des paratextes du \emph{Vieux nègre et la médaille} : hypothèses de lecture}

Appliquons la démarche au roman de Ferdinand Oyono, œuvre au programme de Terminale technique. L'exercice consiste à formuler des \cle{hypothèses de lecture} : des suppositions raisonnées que la lecture du texte devra confirmer ou infirmer.

\textbf{À partir du titre.} \emph{Le Vieux nègre et la médaille} associe deux termes : un personnage (« le vieux nègre », expression qui porte la marque du regard colonial) et un objet symbolique (« la médaille », récompense de l'administration). On peut émettre les hypothèses suivantes :
\begin{itemize}
\item le héros est un vieil homme africain, sans doute soumis à l'autorité coloniale ;
\item la médaille sera l'enjeu du récit : récompense pour services rendus à l'administration ;
\item la coordination « et » suggère un rapport problématique : la médaille honorera-t-elle ou humiliera-t-elle le vieil homme ?
\end{itemize}

\textbf{À partir de l'auteur.} Ferdinand Oyono (1929-2010) est un écrivain camerounais, auteur d'\emph{Une vie de boy} (1956). Son nom évoque la génération des romans de la dénonciation coloniale. Hypothèse : le roman critiquera vraisemblablement l'aliénation des Africains séduits par les valeurs coloniales.

\textbf{À partir de la quatrième de couverture.} Elle présente Meka, vieux paysan qui reçoit une médaille pour avoir donné ses terres à la mission et ses fils à la guerre, avant d'être humilié. Hypothèse confirmée et précisée : le récit sera une satire tragique de la « reconnaissance » coloniale.

\begin{methode}[Formuler des hypothèses de lecture à partir du titre et de la quatrième de couverture]
\begin{enumerate}
\item \textbf{Relever} chaque mot du titre et s'interroger : qui ? quoi ? quel rapport entre les termes ?
\item \textbf{Identifier} le genre suggéré (récit, satire, témoignage) grâce au titre, à l'image et à la collection.
\item \textbf{Situer} l'auteur : époque, nationalité, autres œuvres connues, engagement.
\item \textbf{Croiser} ces indices avec le résumé de la quatrième de couverture, sans le lire intégralement pour garder le plaisir de la découverte.
\item \textbf{Formuler} deux à quatre hypothèses au conditionnel (« le roman pourrait dénoncer... », « le héros sera peut-être... »).
\item \textbf{Vérifier} ces hypothèses au cours de la lecture et les consigner dans le cahier de lecture.
\end{enumerate}
\end{methode}

\begin{savaistu}[Oyono, écrivain de la dénonciation coloniale]
Ferdinand Oyono écrit \emph{Le Vieux nègre et la médaille} en 1956, avant l'indépendance du Cameroun (1\textsuperscript{er} janvier 1960), pour dénoncer l'aliénation coloniale : son héros Meka a tout sacrifié à la France — ses terres, ses fils morts à la guerre — et ne reçoit en échange qu'une médaille dérisoire, symbole d'une reconnaissance hypocrite. Diplomate devenu ministre, Oyono reste l'une des grandes voix de la littérature camerounaise de langue française.
\end{savaistu}

\courssub{Négocier le projet d'étude de l'œuvre}

La lecture d'une œuvre intégrale en classe de Terminale ne s'improvise pas : elle se \emph{négocie} entre le professeur et les élèves. Cette négociation du \cle{projet d'étude} donne à chacun des responsabilités et un calendrier clair.

\begin{definition}[Projet d'étude]
Le \cle{projet d'étude} est un contrat de lecture fixé collectivement : il précise les \textbf{objectifs} (comprendre l'œuvre, enrichir sa culture argumentative, préparer la dissertation), le \textbf{calendrier} (répartition des chapitres par semaine), les \textbf{modalités de lecture} (lecture personnelle, lecture cursive, fiches de lecture) et les \textbf{moments d'évaluation} (contrôles de lecture, compte rendu, remédiation).
\end{definition}

En pratique, la classe fixe par exemple : lecture des chapitres 1 à 3 en semaine 1, contrôle de lecture en semaine 2, débat interprétatif en semaine 3, compte rendu écrit en semaine 4. Chaque élève tient un \cle{cahier de lecture} où il note les personnages, les lieux, les passages marquants et ses hypothèses.

\begin{popfigure}
\begin{tikzpicture}[x=1cm]
\draw[-{Stealth}, thick, popdark] (0,0) -- (12.6,0) node[below left=2pt and 0pt, font=\small]{temps};
\foreach \x/\txt/\col in {0.6/{Négociation\\du projet}/popblue, 3.0/{Lecture\\de l'œuvre}/popgreen, 5.4/{Contrôles\\de lecture}/poporange, 7.8/{Compte\\rendu}/poppurple, 10.2/{Remédiation\\et bilan}/popgold}{
  \draw[\col, line width=2.5pt] (\x,0) -- (\x+1.6,0);
  \node[draw=\col, fill=white, rounded corners, font=\scriptsize, align=center, above=6pt] at (\x+0.8,0) {\txt};
  \fill[\col] (\x,0) circle (2.5pt);
}
\fill[popgold] (12.2,0) circle (2.5pt);
\node[font=\scriptsize\itshape, text=popmuted, below=8pt] at (6.3,0) {Semaine 1 \hspace{1.2cm} Semaine 2 \hspace{1.2cm} Semaine 3 \hspace{1.2cm} Semaine 4};
\end{tikzpicture}
\legende{Frise des étapes du projet d'étude négocié : 1. négociation des objectifs et du calendrier, 2. lecture régulière de l'œuvre, 3. contrôles de lecture intermédiaires, 4. compte rendu écrit, 5. remédiation et bilan collectif.}
\end{popfigure}

\courssub{Lecture du texte ouvroir : première lecture globale}

Le \cle{texte ouvroir} est le texte qui ouvre l'unité d'apprentissage : c'est à partir de lui que seront construits les savoirs de langue et de production. La première lecture du \emph{Vieux nègre et la médaille} (ou de l'extrait proposé) est une lecture \emph{globale} : on lit pour comprendre l'ensemble, sans s'arrêter sur chaque mot difficile.

Le premier travail consiste à repérer le \cle{cadre spatio-temporel} :
\begin{itemize}
\item \textbf{Cadre spatial} : le village camerounais de Doum, la mission catholique, le commandement (le « Mbengui », siège de l'administration) ; l'opposition entre l'espace africain et l'espace colonial structure le récit.
\item \textbf{Cadre temporel} : l'époque coloniale (avant 1960), marquée par l'administration française, le commandant, le père Vandermayer, les gardes.
\item \textbf{Personnages} : Meka (le vieux nègre), sa femme Kelara, Engamba, le commandant, le père Vandermayer, le chef de la région.
\end{itemize}

\begin{exemplebox}[Repérage du cadre dans l'incipit]
Dès les premières pages, Oyono plante le décor : le village s'agite parce que le commandant a convoqué Meka. Les cases, les pistes, la mission au sommet de la colline : tout indique un village camerounais sous domination coloniale. Le lecteur comprend immédiatement que le pouvoir est du côté de l'administration, et l'attente anxieuse des villageois crée le suspense : pourquoi convoque-t-on le vieil homme ?
\end{exemplebox}

\courssub{Contrôle de lecture et compte rendu : \emph{Le Lion et la perle} de Wole Soyinka}

Pour vérifier que la lecture cursive est réellement faite, le professeur organise un \cle{contrôle de lecture} : questions de compréhension (qui, où, quand, quoi), restitution des événements, sens des passages clés. \emph{Le Lion et la perle} (1963), pièce du Nigérian Wole Soyinka, prix Nobel de littérature 1986, met en scène Sidi, la « perle » du village, courtisée par le vieux chef Baroka, le « lion », et par l'instituteur moderniste Lakunle : un affrontement entre tradition et modernité.

Le contrôle débouche sur un \cle{compte rendu de lecture}, exercice à part entière au programme.

\begin{definition}[Compte rendu de lecture]
Le \cle{compte rendu de lecture} est un texte bref qui présente une œuvre de manière \textbf{fidèle et neutre} : il en donne les références (auteur, titre, date, genre), un \textbf{résumé exact} de l'intrigue et des thèmes, puis une \textbf{appréciation argumentée} clairement distinguée des faits. Il respecte la neutralité : on rapporte l'œuvre avant de la juger.
\end{definition}

\begin{attention}[Ne pas confondre résumé et compte rendu critique]
Erreur fréquente au contrôle de lecture : transformer le compte rendu en simple résumé, ou inversement y donner seulement son avis personnel. Le \textbf{résumé} restitue l'intrigue sans jugement ; le \textbf{compte rendu} comporte trois parties distinctes : présentation de l'œuvre, résumé fidèle, appréciation \emph{argumentée} (et non « j'ai aimé / je n'ai pas aimé »). Autre piège : raconter la fin sans l'annoncer, ou trahir l'œuvre par des erreurs de faits (noms, lieux, chronologie).
\end{attention}

\begin{exoresolu}[Compte rendu de lecture : \emph{Le Lion et la perle}]
Rédige en dix à quinze lignes un compte rendu de lecture de la pièce de Wole Soyinka.
\tcblower
\textbf{Présentation.} \emph{Le Lion et la perle} (1963) est une pièce de théâtre du dramaturge nigérian Wole Soyinka. \textbf{Résumé fidèle.} Dans le village d'Ilujinle, la belle Sidi, rendue fière par la publication de ses photos dans un magazine, est courtisée par Lakunle, jeune instituteur qui rêve de modernité et refuse de payer la dot, et par Baroka, le vieux chef polygame, le « lion », qui la séduit par la ruse. Sidi finit par choisir Baroka. \textbf{Appréciation argumentée.} La pièce interroge avec humour le conflit entre tradition et modernité : la « modernité » de Lakunle apparaît superficielle (il imite l'Occident sans le comprendre), tandis que Baroka, malin, sait aussi moderniser son village à sa manière. L'écriture mêle danse, chant et satire, ce qui rend la lecture vivante. On peut toutefois regretter que Sidi soit davantage un enjeu qu'un personnage libre de ses choix.
\end{exoresolu}

\begin{exercice}[Hypothèses de lecture]
À partir du seul titre \emph{Le Vieux nègre et la médaille} et de ce que tu sais de Ferdinand Oyono, formule trois hypothèses de lecture au conditionnel, puis indique quel élément paratextuel te permettrait de les vérifier avant la lecture.
\end{exercice}

\begin{corrige}[Exercice : hypothèses de lecture]
Exemples de réponses attendues : (1) « Le héros pourrait être un vieil homme africain récompensé par l'administration coloniale » — hypothèse fondée sur les deux termes du titre ; (2) « Le roman pourrait dénoncer l'injustice coloniale » — hypothèse fondée sur l'auteur, Oyono, écrivain de la dénonciation ; (3) « La médaille pourrait être un symbole ironique, source d'humiliation plutôt que d'honneur » — hypothèse fondée sur le rapport étrange entre un vieil homme méprisé (« nègre ») et une récompense prestigieuse. La quatrième de couverture et la préface permettraient de vérifier ces hypothèses avant la lecture intégrale.
\end{corrige}

\begin{aretenir}[L'essentiel de la section]
Le \cle{paratexte} (titre, auteur, couverture, quatrième de couverture, préface, dédicace) entoure le texte et oriente sa réception : il permet de formuler des \cle{hypothèses de lecture} vérifiables. L'étude d'une œuvre intégrale s'organise en \cle{projet d'étude négocié} (objectifs, calendrier, modalités, évaluations). La première lecture repère le \cle{cadre spatio-temporel} et les personnages. Le \cle{compte rendu de lecture} articule présentation, résumé fidèle et appréciation argumentée, dans le respect de la neutralité.
\end{aretenir}

\coursec{Le référent et ses substituts}

Dans la séquence précédente, nous avons appris à aborder un texte par ses \cle{paratextes} (titre, nom de l'auteur, couverture, quatrième de couverture) et à négocier un projet d'étude autour du \emph{Vieux nègre et la médaille} de Ferdinand Oyono. Nous savons donc désormais \emph{de quoi} parle le roman : un vieil homme, Meka, reçoit une médaille de l'administration coloniale. Mais une question de langue reste en suspens : comment l'auteur fait-il pour parler de Meka pendant des pages entières sans jamais répéter lourdement « Meka... Meka... Meka... » ? C'est précisément l'objet de cette leçon de grammaire du texte : le \cle{référent} et ses \cle{substituts}. Cette compétence est décisive pour la dissertation au Baccalauréat technique : un devoir clair est un devoir où le correcteur sait toujours \emph{de qui} et \emph{de quoi} l'élève parle.

\courssub{Le référent : la réalité désignée par les mots}

Commençons par une intuition simple. Quand je dis : « Le directeur est entré dans la salle », le mot « directeur » n'est pas une réalité en soi : c'est un signe linguistique qui renvoie à une réalité extérieure à la phrase, une personne bien précise que je pourrais montrer du doigt. Cette réalité désignée, c'est le référent.

\begin{definition}[Référent]
Le \cle{référent} est la réalité extralinguistique — personne, animal, objet, lieu, notion, idée — à laquelle renvoie un mot ou un groupe de mots dans un énoncé. Le référent n'est pas dans le texte : il est ce dont le texte parle.
\end{definition}

Il faut bien distinguer trois choses :
\begin{itemize}
\item le \textbf{mot} (ou signifiant) : la suite de sons ou de lettres, par exemple \emph{M-e-k-a} ;
\item le \textbf{sens} (ou signifié) : l'idée associée au mot, par exemple « vieil homme africain décoré » ;
\item le \textbf{référent} : la réalité désignée, ici le personnage de Meka dans le roman d'Oyono.
\end{itemize}

\begin{exemplebox}[Un même référent, plusieurs mots]
Dans la phrase : « \emph{Meka} reçut sa médaille. \emph{Le vieillard} la fixa sur sa poitrine. \emph{Il} sourit. », les expressions « Meka », « le vieillard » et « il » sont trois mots différents, mais ils désignent \textbf{un seul et même référent} : le personnage de Meka.
\end{exemplebox}

\begin{attention}[Ne pas confondre référent et mot]
Le référent n'est pas un mot du texte. Dire « le référent de \emph{il} est le mot \emph{Meka} » est inexact : le référent de \emph{il} est le \textbf{personnage} Meka ; le mot \emph{Meka} est simplement la première expression qui l'a désigné dans le texte (on l'appelle l'\cle{antécédent}).
\end{attention}

\courssub{Les substituts du référent : les pronoms}

Pourquoi ne répète-t-on pas indéfiniment le même nom ? Parce que la langue dispose d'outils de remplacement : les \cle{substituts}. Leur fonction est double : assurer la \textbf{cohérence} du texte (on comprend qu'on parle toujours de la même chose) et sa \textbf{fluidité} (on évite les répétitions lourdes).

\begin{definition}[Substitut]
Un \cle{substitut} est un mot qui remplace, dans un énoncé, une expression déjà présente dans le texte (ou évidente dans la situation de communication) et qui renvoie au même référent qu'elle.
\end{definition}

\courssub{Les pronoms substituts à connaître}

\textbf{1. Les pronoms personnels.} Ce sont les substituts les plus fréquents : \emph{il, elle, ils, elles} (3\textsuperscript{e} personne, reprise d'un élément déjà nommé), mais aussi les formes complétives \emph{le, la, les, lui, leur, en, y}.

\begin{exemplebox}[Pronoms personnels]
« Meka reçut \emph{sa médaille}. \emph{Il} \emph{la} regarda longtemps. »
\begin{itemize}
\item \emph{il} = substitut de « Meka » (référent : le personnage) ;
\item \emph{la} = substitut de « sa médaille » (référent : l'objet).
\end{itemize}
\end{exemplebox}

\textbf{2. Les pronoms démonstratifs} : \emph{celui, celle, ceux, celles, ceci, cela, ça, ce}. Ils reprennent un nom ou toute une idée.

« Les Blancs décidèrent de décorer Meka. \emph{Cela} surprit tout le village. » — \emph{cela} reprend toute la proposition précédente.

\textbf{3. Les pronoms possessifs} : \emph{le mien, la tienne, le sien, les nôtres...} Ils remplacent un nom et indiquent en même temps la possession.

« Tous les anciens avaient une médaille ; Meka montra \emph{la sienne} avec fierté. » — \emph{la sienne} = « sa médaille à lui ».

\textbf{4. Les adjectifs possessifs} : \emph{mon, ton, son, ma, ta, sa, mes, tes, ses, notre, votre, leur...} Attention : ce ne sont pas des substituts à proprement parler, car ils accompagnent un nom (« \emph{sa} médaille »), mais ils participent à la chaîne référentielle en rattachant un objet à un possesseur déjà identifié.

\courssub{La reprise nominale : répétition, synonyme, hyperonyme, périphrase}

Les pronoms ne sont pas les seuls substituts. Un texte bien écrit utilise aussi la \cle{reprise nominale}, c'est-à-dire un \textbf{groupe nominal} qui désigne à nouveau le même référent. On distingue quatre procédés :

\begin{center}
\begin{tabularx}{\linewidth}{|l|X|X|}
\hline
\rowcolor{popblueL} \textbf{Procédé} & \textbf{Principe} & \textbf{Exemple (référent : Meka)} \\
\hline
Répétition du nom & On reprend le même mot & « Meka... Meka se leva. » \\
\hline
Synonyme & On reprend un mot de sens proche & « Meka... le vieillard sourit. » \\
\hline
Hyperonyme & On reprend un terme plus général & « Meka... le vieux nègre... l'homme... » \\
\hline
Périphrase & On décrit le référent en quelques mots & « Meka... ce brave homme... le père de famille décoré... » \\
\hline
\end{tabularx}
\end{center}

La reprise nominale est précieuse dans une dissertation : elle permet de varier les désignations tout en apportant une nuance (la périphrase « ce brave homme » ajoute un jugement de valeur que le pronom « il » ne peut pas porter).

\courssub{La chaîne de référence}

\begin{definition}[Chaîne de référence]
Une \cle{chaîne de référence} est l'ensemble des expressions d'un texte (nom propre, groupes nominaux, pronoms) qui renvoient à un même référent. Le premier maillon, qui introduit le référent, s'appelle l'\cle{antécédent} ; les maillons suivants sont les \cle{reprises} (pronominales ou nominales).
\end{definition}

Suivre une chaîne de référence, c'est donc suivre un personnage ou une notion à travers tout un texte. C'est un exercice fondamental de lecture (pour le commentaire et le résumé) et d'écriture (pour la dissertation).

\begin{popfigure}
\begin{center}
\begin{tikzpicture}[font=\small]
\node[draw=popdark, fill=popblueL, rounded corners, thick, minimum width=2.6cm, minimum height=1cm, align=center] (meka) at (0,0){\textbf{MEKA}\\ \footnotesize (référent : le personnage)};
\node[draw=popblue, rounded corners, fill=white, align=center] (a) at (-4.6,2.2){« Meka »\\ \footnotesize antécédent};
\node[draw=popblue, rounded corners, fill=white, align=center] (b) at (4.6,2.2){« le vieillard »\\ \footnotesize reprise nominale (synonyme)};
\node[draw=popblue, rounded corners, fill=white, align=center] (c) at (-4.6,-2.2){« il »\\ \footnotesize reprise pronominale};
\node[draw=popblue, rounded corners, fill=white, align=center] (d) at (0,-2.8){« le vieux nègre »\\ \footnotesize reprise nominale (hyperonyme)};
\node[draw=popblue, rounded corners, fill=white, align=center] (e) at (4.6,-2.2){« ce brave homme »\\ \footnotesize reprise nominale (périphrase)};
\draw[-{Stealth}, thick, poporange] (a) -- (meka);
\draw[-{Stealth}, thick, poporange] (b) -- (meka);
\draw[-{Stealth}, thick, poporange] (c) -- (meka);
\draw[-{Stealth}, thick, poporange] (d) -- (meka);
\draw[-{Stealth}, thick, poporange] (e) -- (meka);
\end{tikzpicture}
\end{center}
\legende{Chaîne de référence autour du personnage de Meka dans \emph{Le Vieux nègre et la médaille} : toutes les expressions (flèches) renvoient à un seul et même référent.}
\end{popfigure}

\begin{exoresolu}[Identifier les référents des pronoms]
Énoncé. Dans la phrase : « Meka reçut sa médaille. Il la regarda longtemps », identifiez le référent de « il » et de « la », puis précisez la nature de chaque substitut.
\tcblower
Solution détaillée.
\begin{enumerate}
\item « \emph{Il} » est un \textbf{pronom personnel} sujet, 3\textsuperscript{e} personne du singulier. Son antécédent est « Meka » ; son référent est donc le personnage de Meka.
\item « \emph{La} » est un \textbf{pronom personnel} complément d'objet direct, 3\textsuperscript{e} personne du singulier féminin. Son antécédent est « sa médaille » ; son référent est l'objet, la médaille remise par l'administration.
\item « \emph{sa} » (dans « sa médaille ») est un \textbf{adjectif possessif} qui rattache la médaille à Meka : il participe lui aussi à la chaîne de référence du personnage.
\end{enumerate}
Conclusion : deux substituts différents peuvent coexister dans la même phrase parce qu'ils renvoient à deux référents différents (le personnage et l'objet). C'est la clarté du contexte qui permet de les distinguer sans effort.
\end{exoresolu}

\courssub{Le rôle des substituts dans la dissertation}

Dans une dissertation, les chaînes de référence remplissent trois fonctions majeures :

\begin{itemize}
\item \textbf{Cohérence} : le correcteur suit sans peine le fil du raisonnement parce que chaque pronom renvoie clairement à son antécédent (« Ferdinand Oyono dénonce l'absurdité coloniale. \emph{Il} montre que... »).
\item \textbf{Fluidité} : les substituts évitent les répétitions maladroites qui alourdissent le style. Comparer : « Meka reçut la médaille, Meka fixa la médaille, Meka sourit » (lourd) et « Meka reçut la médaille, la fixa sur sa poitrine et sourit » (fluide).
\item \textbf{Précision} : la reprise nominale permet de désigner le référent sous un angle utile à l'argumentation : « le vieillard » insiste sur l'âge, « ce brave homme » sur la sympathie, « le vieux nègre » sur le regard colonial.
\end{itemize}

\begin{aretenir}[L'essentiel]
Un texte clair est un texte où chaque substitut renvoie sans hésitation possible à un référent unique. Avant d'utiliser un pronom, le rédacteur doit pouvoir répondre instantanément à la question : « de qui ou de quoi est-ce que je parle ici ? »
\end{aretenir}

\begin{methode}[Vérifier la clarté référentielle de son propre texte]
Pour contrôler votre devoir avant de le rendre, appliquez le \textbf{test de la question} :
\begin{enumerate}
\item Relisez votre paragraphe phrase par phrase.
\item Pour chaque pronom (\emph{il, elle, ils, le, la, leur, celui, cela...}) et chaque possessif (\emph{son, sa, leur...}), posez-vous la question : « \emph{de qui ? de quoi ?} ».
\item Cherchez l'antécédent : il doit être \textbf{unique}, \textbf{récent} (dans la phrase ou la phrase précédente) et \textbf{accordé} en genre et en nombre avec le pronom.
\item Si deux référents sont possibles (deux hommes, deux idées), remplacez le pronom par une reprise nominale : « le romancier », « cette seconde idée », « le personnage ».
\item Si l'antécédent est trop loin (plus de deux phrases), répétez le nom ou utilisez un synonyme.
\end{enumerate}
\end{methode}

\begin{attention}[Erreur fréquente : l'ambiguïté référentielle]
Le piège classique de la copie d'examen consiste à accumuler les pronoms « il » renvoyant à \textbf{deux personnages différents} dans la même phrase. Exemple fautif : « Meka rencontra le commandant ; \emph{il} était furieux. » — Qui est furieux ? Meka ou le commandant ? Le lecteur ne peut pas trancher : c'est une \cle{ambiguïté référentielle}. Corrections possibles : « Meka rencontra le commandant ; \emph{ce dernier} était furieux » ou « Meka rencontra le commandant, \emph{qui} était furieux ». Au Baccalauréat, chaque ambiguïté de ce type fait perdre des points de cohérence textuelle.
\end{attention}

\begin{exercice}[Repérage d'une chaîne de référence]
Relevez toutes les expressions qui renvoient au référent « Meka » dans le passage suivant, et classez-les (antécédent, reprise pronominale, reprise nominale) : « Meka attendait depuis l'aube. Le vieillard s'était habillé avec soin. Il tenait sa canne d'une main tremblante. Ce brave homme, que tout le village respectait, allait enfin recevoir sa récompense. On la lui remit devant la foule assemblée. »
\end{exercice}

\begin{corrige}[Exercice 1]
Chaîne de référence de « Meka » :
\begin{itemize}
\item Antécédent : « Meka » (nom propre, premier maillon).
\item Reprises nominales : « Le vieillard » (synonyme) ; « Ce brave homme » (périphrase avec jugement de valeur).
\item Reprises pronominales : « Il » (pronom personnel sujet) ; « lui » (dans « la lui remit », pronom personnel COI).
\item Adjectifs possessifs rattachés à la chaîne : « sa canne », « sa récompense ».
\end{itemize}
Le passage est clair car chaque reprise ne peut renvoyer qu'à Meka : aucun autre personnage masculin singulier n'interfère.
\end{corrige}

En résumé, maîtriser le référent et ses substituts, c'est tenir le fil invisible qui relie toutes les phrases d'un texte. Cette compétence de lecture nous servira immédiatement pour analyser le vocabulaire de l'argumentation — synonymie, antonymie, hypéronymie — dans la suite du cours, puis pour construire des paragraphes de dissertation parfaitement cohérents.

\coursec{Lexique de l'argumentation : synonymie, antonymie, hypéronymie}

Dans la section précédente, nous avons vu que le référent d'un texte peut être repris par des substituts : pronoms, groupes nominaux équivalents, périphrases. Mais ces reprises ne sont jamais de simples copies : l'auteur choisit ses mots avec soin, en jouant sur la proximité de sens, l'opposition ou le degré de généralité. C'est précisément ce jeu sur les mots que nous allons étudier maintenant. Pour argumenter efficacement dans une dissertation, il ne suffit pas d'avoir des idées : il faut disposer d'un lexique riche et précis, capable de dire la même chose autrement, d'opposer deux notions, de classer des réalités du plus général au plus particulier. Trois relations de sens fondamentales nous y aideront : la \cle{synonymie}, l'\cle{antonymie} et le couple \cle{hypéronymie}/\cle{hyponymie}. Nous verrons ensuite comment ces outils permettent de déceler le contenu latent d'un sujet et de construire un champ lexical complet avant de rédiger.

\courssub{La synonymie : dire presque la même chose autrement}

\textbf{Intuition.} Observez ces trois phrases : « Le candidat a produit un \emph{devoir} remarquable », « Le candidat a produit une \emph{copie} remarquable », « Le candidat a produit un \emph{travail} remarquable ». Les mots \emph{devoir}, \emph{copie} et \emph{travail} désignent la même réalité, mais ils ne se ressemblent pas et ne s'emploient pas dans les mêmes situations. C'est le phénomène de synonymie.

\begin{definition}[La synonymie]
La \cle{synonymie} est la relation de sens qui unit deux ou plusieurs mots de signification très proche, appelés \cle{synonymes}. On distingue :
\begin{itemize}
\item les \textbf{synonymes parfaits} (ou totaux), interchangeables dans presque tous les contextes : \emph{couloir}/\emph{corridor}, \emph{commencer}/\emph{débuter} ;
\item les \textbf{synonymes partiels}, qui ne sont interchangeables que dans certains contextes : on dit \emph{un homme grand} et \emph{un homme de haute taille}, mais on ne dit pas \emph{une montagne grande} pour \emph{une haute montagne}.
\end{itemize}
\end{definition}

\textbf{Les nuances de sens et de registre.} Deux synonymes ne sont presque jamais strictement équivalents. Ils diffèrent par trois types de nuances :

\begin{itemize}
\item la \textbf{nuance d'intensité} : \emph{aimer}, \emph{adorer}, \emph{idolâtrer} forment une gradation croissante ;
\item la \textbf{nuance de sens} : \emph{regarder}, \emph{observer}, \emph{contempler}, \emph{scruter} n'impliquent pas la même manière de voir ;
\item la \textbf{nuance de registre} : \emph{travail} (courant), \emph{besogne} (familier, parfois péjoratif), \emph{labeur} (soutenu, littéraire), \emph{boulot} (très familier, à bannir d'une dissertation).
\end{itemize}

\begin{exemplebox}[Gradation de synonymes]
Pour le verbe \emph{critiquer} : \emph{discuter} $<$ \emph{contester} $<$ \emph{condamner} $<$ \emph{fustiger}. Dans une dissertation sur \emph{Le Vieux nègre et la médaille}, on pourra écrire que Ferdinand Oyono \emph{fustige} l'administration coloniale : le choix du synonyme le plus fort traduit la violence satirique du roman.
\end{exemplebox}

\begin{attention}[Le piège du registre]
L'erreur la plus fréquente consiste à employer un synonyme d'un registre inadapté à la situation d'écriture. La dissertation exige un \textbf{registre soutenu ou courant}. Écrire « Meka a trimé au boulot toute sa vie » au lieu de « Meka a accompli un dur labeur toute sa vie » dévalorise la copie, même si le sens est correct. Avant d'employer un synonyme, demandez-vous toujours : ce mot convient-il à un examen officiel ?
\end{attention}

\courssub{L'antonymie : opposer pour mieux délimiter}

\textbf{Intuition.} Pour définir la justice, rien de plus efficace que de la confronter à son contraire : l'injustice. L'argumentation procède constamment par opposition (thèse/antithèse, avantage/inconvénient, pour/contre). Maîtriser l'antonymie, c'est donc maîtriser l'outil de base du raisonnement dialectique.

\begin{definition}[L'antonymie]
L'\cle{antonymie} est la relation de sens qui unit deux mots de signification opposée, appelés \cle{antonymes}. On distingue :
\begin{itemize}
\item les \textbf{antonymes par préfixation}, formés en ajoutant un préfixe négatif (\emph{in-}, \emph{im-}, \emph{dé(s)-}, \emph{a-}, \emph{anti-}) : \emph{juste}/\emph{injuste}, \emph{possible}/\emph{impossible}, \emph{accord}/\emph{désaccord}, \emph{typique}/\emph{atypique} ;
\item les \textbf{antonymes lexicaux}, mots différents sans rapport de forme : \emph{guerre}/\emph{paix}, \emph{richesse}/\emph{pauvreté}, \emph{liberté}/\emph{esclavage}.
\end{itemize}
On distingue enfin deux types d'opposition :
\begin{itemize}
\item les \textbf{contraires} (opposition graduable, avec des degrés intermédiaires) : entre \emph{chaud} et \emph{froid}, il existe \emph{tiède} ; entre \emph{riche} et \emph{pauvre}, il existe toute une échelle sociale ;
\item les \textbf{contradictoires} (opposition absolue, sans intermédiaire) : \emph{vivant}/\emph{mort}, \emph{présent}/\emph{absent} ; nier l'un, c'est affirmer l'autre.
\end{itemize}
\end{definition}

\begin{exemplebox}[Antonymes et argumentation]
Sujet : « Le développement est-il toujours bénéfique ? » Les antonymes permettent de construire l'axe opposé du plan : \emph{développement}/\emph{sous-développement} (préfixation), \emph{progrès}/\emph{régression}, \emph{enrichissement}/\emph{appauvrissement}. La partie antithétique de la dissertation s'appuiera naturellement sur ces termes opposés.
\end{exemplebox}

\courssub{Hypéronymie et hyponymie : du général au particulier}

\textbf{Intuition.} Si l'on vous demande de citer des moyens de transport au Cameroun, vous répondrez spontanément : la moto-taxi, le bus, le train, le benskin (moto-taxi urbain). Tous ces mots sont des cas particuliers du mot général \emph{transport}. Cette relation d'inclusion est essentielle pour organiser les idées d'une dissertation : elle permet de passer de la notion abstraite du sujet à des exemples concrets.

\begin{definition}[Hypéronymie et hyponymie]
L'\cle{hypéronymie} est la relation qui unit un terme générique, l'\cle{hyperonyme}, à des termes spécifiques, les \cle{hyponymes}, dont le sens est inclus dans le sien. Ainsi :
\begin{itemize}
\item \emph{injustice} est l'hyperonyme de \emph{corruption}, \emph{tribalisme}, \emph{népotisme}, \emph{favoritisme} ;
\item \emph{transport} est l'hyperonyme de \emph{moto-taxi}, \emph{bus}, \emph{train}, \emph{benskin}.
\end{itemize}
Inversement, on dit que \emph{corruption} est un hyponyme d'\emph{injustice}. La relation est hiérarchique et peut comporter plusieurs niveaux : \emph{fait social} $>$ \emph{injustice} $>$ \emph{corruption} $>$ \emph{corruption policière}.
\end{definition}

\begin{popfigure}
\begin{center}
\begin{tikzpicture}[
  level distance=1.5cm,
  level 1/.style={sibling distance=3.4cm},
  every node/.style={draw=popblue, fill=popblueL, rounded corners=2pt, font=\small, inner sep=4pt},
  edge from parent/.style={draw=popmuted, thick}
]
\node[fill=poporangeL, draw=poporange, font=\small\bfseries] {TRANSPORT (hyperonyme)}
  child { node {moto-taxi} }
  child { node {bus} }
  child { node {benskin} }
  child { node {train} };
\end{tikzpicture}
\end{center}
\legende{Arbre hypéronymique : le terme générique « transport » domine quatre hyponymes adaptés au contexte camerounais. Chaque hyponyme peut à son tour devenir hyperonyme (le bus : car interurbain, bus urbain, coaster...).}
\end{popfigure}

\textbf{Intérêt de la hiérarchie pour la dissertation.} Lorsque le sujet porte sur un hyperonyme (par exemple \emph{l'injustice}), les hyponymes fournissent les axes du plan : une partie sur la corruption, une partie sur le tribalisme, une partie sur le népotisme. Réciproquement, lorsque le sujet porte sur un hyponyme, l'hyperonyme permet de situer la notion dans un cadre plus large en introduction.

\courssub{Contenus manifestes et contenus latents : dénotation et connotation}

\textbf{Intuition.} Un personnage du \emph{Lion et la perle} déclare : « Nos jeunes filles ont désormais le goût des choses de la ville. » En apparence, la phrase parle de goûts vestimentaires. En réalité, elle suggère une critique de la modernité, de la perte des traditions, peut-être de la corruption des mœurs. Ce que la phrase dit explicitement, c'est son contenu manifeste ; ce qu'elle laisse entendre, c'est son contenu latent.

\begin{definition}[Contenu manifeste et contenu latent]
\begin{itemize}
\item Le \cle{contenu manifeste} d'un énoncé est ce qu'il dit explicitement, son sens littéral et directement observable. Il relève de la \cle{dénotation} : le sens objectif, stable, dictionnairique des mots.
\item Le \cle{contenu latent} est ce qu'un énoncé suggère sans le dire : l'\cle{implicite}, les sous-entendus, les présupposés. Il relève de la \cle{connotation} : les valeurs affectives, culturelles ou idéologiques associées aux mots.
\end{itemize}
Ainsi, le mot \emph{foyer} dénote le lieu où l'on fait le feu, mais connote la chaleur familiale, l'intimité, la sécurité.
\end{definition}

\begin{exemplebox}[Dénotation et connotation dans un sujet de dissertation]
Sujet : « L'école est-elle la seule voie de la réussite ? » Contenu manifeste : une question sur le rapport entre scolarisation et succès sociale. Contenus latents : le mot \emph{seule} présuppose que d'autres voies existent (commerce, artisanat, agriculture, sport) ; le mot \emph{réussite} connote la richesse et le prestige social, mais peut aussi signifier l'épanouissement personnel. Repérer ces implicites, c'est déjà trouver la problématique.
\end{exemplebox}

\courssub{Intérêt argumentatif de ces relations lexicales}

Ces quatre outils lexicaux ne sont pas de simples curiosités grammaticales : ils structurent directement le travail du dissertateur.

\begin{center}
\begin{tabularx}{\linewidth}{|l|X|X|}
\hline
\rowcolor{popblueL} \textbf{Relation} & \textbf{Fonction argumentative} & \textbf{Exemple d'emploi} \\
\hline
Synonymie & Éviter les répétitions ; graduer les idées & \emph{travail} $\rightarrow$ \emph{labeur}, \emph{activité}, \emph{métier} \\
\hline
Antonymie & Construire l'opposition thèse/antithèse & \emph{tradition}/\emph{modernité} \\
\hline
Hypéronymie & Catégoriser, annoncer le plan & \emph{injustice} : corruption, tribalisme, népotisme \\
\hline
Connotation & Décoder les sous-entendus du sujet & \emph{« seule voie »} = présupposé à discuter \\
\hline
\end{tabularx}
\end{center}

\begin{popfigure}
\begin{center}
\begin{tikzpicture}
\draw[thick, popmuted] (0,0) rectangle (10,4.4);
\draw[thick, popmuted] (0,2.2) -- (10,2.2);
\draw[thick, popmuted] (5,0) -- (5,4.4);
\node[fill=popgreenL, draw=popgreen, rounded corners=2pt, font=\small\bfseries] at (2.5,4.85) {SYNONYMES de « développement »};
\node[fill=poppinkL, draw=poppink, rounded corners=2pt, font=\small\bfseries] at (7.5,4.85) {ANTONYMES de « développement »};
\node[font=\small, align=left, anchor=north west] at (0.25,2.0) {\textbf{Registre courant :} progrès,\\ croissance, essor};
\node[font=\small, align=left, anchor=north west] at (0.25,4.0) {\textbf{Registre soutenu :} épanouissement,\\ avènement, expansion};
\node[font=\small, align=left, anchor=north west] at (5.25,2.0) {\textbf{Par préfixation :}\\ sous-développement};
\node[font=\small, align=left, anchor=north west] at (5.25,4.0) {\textbf{Lexicaux :} régression,\\ stagnation, déclin, recul};
\end{tikzpicture}
\end{center}
\legende{Tableau à double entrée : synonymes et antonymes de « développement » classés selon la nuance de registre ou le mode de formation.}
\end{popfigure}

\courssub{Méthode : constituer le champ lexical d'un sujet avant de rédiger}

\begin{methode}[Construire un champ lexical en quatre étapes]
Avant toute recherche d'idées, appliquez systématiquement la démarche suivante au mot-clé du sujet :
\begin{enumerate}
\item \textbf{Identifier l'hyperonyme} : à quelle catégorie générale appartient la notion ? (\emph{travail} $\rightarrow$ \emph{activité humaine}, \emph{fait social}).
\item \textbf{Lister les hyponymes} : quelles réalités précises recouvre la notion ? (\emph{travail} $\rightarrow$ agriculture, artisanat, fonction publique, commerce, petits métiers).
\item \textbf{Lister les synonymes avec leurs registres} : \emph{travail} $\rightarrow$ labeur (soutenu), métier, emploi, occupation (courants), besogne, boulot (familiers, à éviter).
\item \textbf{Lister les antonymes} : \emph{travail} $\rightarrow$ oisiveté, chômage, paresse, repos, loisir. Ils préparent la partie antithétique du plan.
\end{enumerate}
Ce tableau lexical, rédigé au brouillon en cinq minutes, garantit une copie sans répétitions et un plan déjà esquissé.
\end{methode}

\begin{exoresolu}[Champ lexical du sujet « le travail »]
Énoncé. Sujet de dissertation : « Le travail est-il une contrainte ou une source d'épanouissement ? » Constituez le champ lexical du mot-clé \emph{travail} selon la méthode en quatre étapes, puis dégagez le contenu latent du sujet.
\tcblower
Solution détaillée.
\begin{itemize}
\item \textbf{Hyperonymes} : activité humaine, fait social, obligation sociale.
\item \textbf{Hyponymes} : agriculture, élevage, artisanat, enseignement, commerce, fonction publique, travail domestique, petits métiers (benskin, call-box, vente ambulante).
\item \textbf{Synonymes} : labeur (soutenu, connote l'effort), métier (connote la qualification), emploi (connote le salariat), occupation (neutre), besogne et boulot (familiers, à proscrire).
\item \textbf{Antonymes} : chômage, oisiveté, paresse (contraires péjoratifs) ; repos, loisir (contraires valorisés).
\item \textbf{Contenu latent} : l'alternative « contrainte \emph{ou} épanouissement » présuppose que les deux sont incompatibles ; or le correcteur attend qu'on discute ce présupposé : le travail peut être \emph{à la fois} contrainte (nécessité de survivre) et épanouissement (réalisation de soi). Le mot \emph{épanouissement} connote le bonheur et l'accomplissement personnel, au-delà du simple salaire.
\end{itemize}
Ce travail lexical fournit déjà les deux axes du plan dialectique : I. Le travail comme contrainte (labeur, fatigue, exploitation) ; II. Le travail comme épanouissement (métier, créativité, dignité).
\end{exoresolu}

\begin{exoresolu}[Repérer le contenu latent dans \emph{Le Lion et la perle}]
Énoncé. Dans \emph{Le Lion et la perle} de Wole Soyinka, le chef Baroka est surnommé « le Lion » et affirme défendre la tradition contre l'école et le progrès incarnés par Lakunle. Dégager le contenu manifeste et le contenu latent de cette affirmation.
\tcblower
Solution détaillée. Contenu manifeste : Baroka s'oppose publiquement à la modernité (école, routes, chemises occidentales) et se présente comme le gardien des coutumes du village d'Ilujinle. Contenu latent : le surnom « le Lion » connote la force, la ruse et la domination ; en réalité, Baroka instrumentalise la tradition pour conserver son pouvoir et satisfaire ses désirs (séduire Sidi, la « perle »). Sa défense des coutumes cache donc une stratégie de domination. L'élève qui repère cet implicite comprend que la pièce ne oppose pas simplement tradition et modernité, mais dénonce les manipulations que les deux discours peuvent servir : Lakunle lui-même use du « progrès » pour imposer ses vues à Sidi. Ce décodage nourrit directement la problématique d'une dissertation sur le conflit des générations ou sur tradition et modernité.
\end{exoresolu}

\begin{savaistu}[Le lexique, critère officiel de correction]
Au Probatoire et au Baccalauréat technique, la grille de correction de la dissertation comporte des critères explicites portant sur la langue : richesse et précision du vocabulaire, correction de l'expression, absence de répétitions. Deux copies aux idées identiques peuvent obtenir des notes très différentes selon la qualité du lexique employé. Constituer un champ lexical au brouillon avant de rédiger n'est donc pas un luxe : c'est un investissement directement payé en points.
\end{savaistu}

\begin{exercice}[Entraînement]
1. Classez les mots suivants en synonymes et antonymes de \emph{liberté}, en précisant le registre : \emph{affranchissement}, \emph{captivité}, \emph{indépendance}, \emph{servitude}, \emph{autonomie}, \emph{esclavage}, \emph{émancipation}.

2. Construisez l'arbre hypéronymique du mot \emph{enseignement} avec au moins quatre hyponymes adaptés au système éducatif camerounais.

3. Sujet : « Faut-il préserver les traditions ? » Dégagez le contenu manifeste et deux contenus latents de ce sujet, puis constituez le champ lexical de \emph{tradition} en quatre étapes.
\end{exercice}

\begin{corrige}[Exercice]
1. Synonymes : \emph{affranchissement} et \emph{émancipation} (soutenus), \emph{indépendance} et \emph{autonomie} (courants). Antonymes : \emph{captivité}, \emph{servitude} (soutenu), \emph{esclavage} (courant). Nuance : \emph{indépendance} s'applique plutôt aux nations, \emph{autonomie} aux individus ou aux institutions.

2. Hyperonyme : \emph{enseignement}. Hyponymes : enseignement primaire, enseignement secondaire général, enseignement secondaire technique, enseignement supérieur, formation professionnelle. Chaque hyponyme peut être subdivisé (secondaire technique : lycée technique, CETIC...).

3. Contenu manifeste : question sur la valeur de la conservation des coutumes. Contenus latents : le verbe \emph{préserver} connote que les traditions sont menacées (par la modernité, la mondialisation) ; l'interrogation présuppose que certaines traditions pourraient être abandonnées (excision, mariages forcés), ce qui ouvre la voie à une réponse nuancée. Champ lexical de \emph{tradition} : hyperonyme \emph{culture} ; hyponymes : rites, danses, langues, cultes, artisanat ; synonymes : coutume, usage, héritage (soutenu), folklore (parfois péjoratif) ; antonymes : modernité, innovation, progrès, rupture.
\end{corrige}

\begin{aretenir}[L'essentiel de la section]
\begin{itemize}
\item La \cle{synonymie} unit des mots de sens proche, rarement interchangeables : attention aux nuances d'intensité et de registre.
\item L'\cle{antonymie} unit des mots opposés (par préfixation ou lexicaux), contraires ou contradictoires : c'est l'outil de l'opposition argumentative.
\item L'\cle{hyperonyme} (terme générique) domine ses \cle{hyponymes} (termes spécifiques) : c'est l'outil du classement des idées et de l'annonce de plan.
\item Tout énoncé possède un \cle{contenu manifeste} (dénotation) et un \cle{contenu latent} (connotation, implicite, présupposés) que le candidat doit apprendre à décoder dans le sujet.
\item Avant de rédiger : hyperonyme $\rightarrow$ hyponymes $\rightarrow$ synonymes $\rightarrow$ antonymes. Cinq minutes de brouillon lexical pour une copie plus riche et mieux structurée.
\end{itemize}
\end{aretenir}

\coursec{Le texte argumentatif et l'analyse du sujet}

Dans la section précédente, nous avons enrichi notre vocabulaire de l'argumentation grâce à la synonymie, à l'antonymie et à l'hypéronymie : nous savons désormais varier les mots pour éviter les répétitions et exprimer des nuances fines. Mais un beau vocabulaire ne suffit pas à réussir une dissertation. Encore faut-il comprendre \emph{ce qu'on nous demande exactement} et \emph{comment construire un raisonnement convaincant}. C'est l'objet de cette section : d'abord reconnaître les caractéristiques du texte argumentatif, puis apprendre à analyser un sujet de dissertation pas à pas, jusqu'à la formulation de la problématique.

\courssub{Les caractéristiques du texte argumentatif}

\textbf{Intuition.} Quand un élève essaie de convaincre ses parents de lui acheter un téléphone, il ne se contente pas de dire « je veux un téléphone ». Il avance des raisons (« tous mes camarades en ont un », « je pourrai réviser avec les groupes WhatsApp de la classe »), il illustre (« regarde, Marlyse a progressé en maths grâce aux leçons en ligne ») et il conclut. Sans le savoir, il produit un \cle{texte argumentatif}.

\begin{definition}[Le texte argumentatif]
Un \cle{texte argumentatif} est un texte qui défend une \cle{thèse} (une opinion, un point de vue) en s'appuyant sur des \cle{arguments} (des raisons) illustrés par des \cle{exemples}, dans le but de convaincre ou de persuader un destinataire. Il est organisé par des \cle{connecteurs logiques} qui articulent le raisonnement.
\end{definition}

\textbf{Les cinq piliers du texte argumentatif.}
\begin{itemize}
\item La \cle{thèse} : c'est l'idée que l'on défend. Exemple : « L'école reste la voie royale vers la réussite. »
\item Les \cle{arguments} : ce sont les raisons qui justifient la thèse. Exemple : « L'école transmet des savoirs et des diplômes exigés par le marché du travail. »
\item Les \cle{exemples} : ce sont des cas précis qui rendent l'argument concret et crédible. Exemple : « La plupart des cadres camerounais recrutés dans la fonction publique sont titulaires d'au moins une licence. »
\item Les \cle{connecteurs logiques} : ils marquent les relations entre les idées : \emph{d'abord, ensuite, en effet, car, cependant, néanmoins, par conséquent, enfin}. Ils sont les « charnières » du devoir.
\item La \cle{visée persuasive} : le texte cherche à faire adhérer le lecteur à la thèse défendue, par la raison (convaincre) ou par les émotions (persuader).
\end{itemize}

\begin{propriete}[Structure d'un bon argument]
Un bon argument se compose de \textbf{trois éléments indissociables} :
\[
\text{Argument complet} = \text{une idée} + \text{une justification} + \text{un exemple précis}.
\]
Un argument sans exemple reste abstrait ; un exemple sans idée reste une anecdote. Au Baccalauréat technique, chaque paragraphe du développement doit contenir ces trois éléments.
\end{propriete}

\textbf{Les types d'arguments.} Tous les arguments n'ont pas la même force ni la même nature. On en distingue classiquement quatre :

\begin{center}
\begin{tabularx}{\linewidth}{|l|X|X|}
\hline
\rowcolor{popblueL} \textbf{Type d'argument} & \textbf{Principe} & \textbf{Exemple} \\
\hline
Argument d'\cle{autorité} & On s'appuie sur la parole d'une personne compétente, d'un spécialiste, d'une institution. & « Selon l'OMS, le tabac tue chaque année des millions de personnes. » \\
\hline
Argument d'\cle{expérience} & On invoque un fait vécu, observé, une réalité concrète. & « Dans mon quartier, les jeunes diplômés trouvent plus facilement un emploi. » \\
\hline
Argument de \cle{logique} & On raisonne par enchaînement de causes et de conséquences. & « Sans formation, pas de qualification ; sans qualification, pas d'emploi stable. » \\
\hline
Argument de \cle{valeur} & On fait appel à un principe moral, culturel ou social admis (justice, honneur, travail). & « Le travail honnête est la fierté de l'homme : mieux vaut réussir par le mérite. » \\
\hline
\end{tabularx}
\end{center}

\begin{attention}[Piège fréquent]
L'argument d'autorité n'est solide que si l'autorité citée est \emph{compétente dans le domaine} et si la citation est exacte. Inventer une citation ou attribuer à un savant une parole qu'il n'a jamais tenue est sanctionné par le correcteur. Si tu n'es pas sûr de la source, préfère une formule prudente : « Comme le rappellent de nombreux éducateurs... »
\end{attention}

\courssub{L'analyse du sujet : une étape décisive}

\textbf{Intuition.} Un menuisier qui se trompe dans la lecture du plan du client fabriquera une belle porte... qui ne correspond pas à la commande. De même, un candidat qui lit mal le sujet rédigera une belle copie... hors sujet. L'analyse du sujet est la phase où l'on « décode la commande » du correcteur.

\begin{methode}[Analyser un sujet de dissertation en 5 étapes]
\begin{enumerate}
\item \textbf{Lire} le sujet au moins trois fois, lentement, sans rien écrire, puis en le reformulant mentalement avec ses propres mots.
\item \textbf{Souligner} les mots-clés : les noms importants, les verbes, les adverbes restrictifs (\emph{seule, toujours, vraiment, nécessairement}) et la consigne.
\item \textbf{Définir} chaque mot-clé au brouillon : donner son sens, ses synonymes, ses contraires, son champ d'application.
\item \textbf{Relier} les termes entre eux : quel rapport le sujet établit-il entre eux ? Y a-t-il opposition, équivalence, condition ?
\item \textbf{Questionner} : transformer cette tension en une ou deux questions directrices — c'est la problématique.
\end{enumerate}
\end{methode}

\textbf{Repérer la consigne.} La formulation du sujet n'est jamais neutre : elle indique le type de travail attendu.
\begin{itemize}
\item « \emph{Pensez-vous que... ?} » : le sujet invite à prendre position clairement (plan dialectique ou thématique possible).
\item « \emph{Discutez cette opinion...} » : il faut confronter plusieurs points de vue, le pour et le contre (plan dialectique).
\item « \emph{Commentez cette affirmation...} » : il faut expliquer, apprécier, nuancer une pensée donnée, souvent une citation.
\end{itemize}

\courssub{Les trois types de sujets}

\begin{definition}[Types de sujets de dissertation]
On distingue trois grands types de sujets : le \cle{sujet de réflexion} (une question directe à laquelle on répond en organisant ses idées), le \cle{sujet de commentaire} (une affirmation ou une citation à expliquer et à apprécier) et le \cle{sujet de discussion} (une opinion à confronter à d'autres points de vue).
\end{definition}

\begin{popfigure}
\begin{tikzpicture}[font=\small]
\node[draw=popblue, fill=popblueL, rounded corners, text width=3.1cm, align=center, minimum height=2.6cm] (A) at (0,0)
{\textbf{Sujet de réflexion}\\[2pt] Consigne : \emph{Pensez-vous que...}\\ Plan : thématique ou dialectique};
\node[draw=poporange, fill=poporangeL, rounded corners, text width=3.1cm, align=center, minimum height=2.6cm] (B) at (5.4,0)
{\textbf{Sujet de commentaire}\\[2pt] Consigne : \emph{Commentez...}\\ Plan : explication puis appréciation};
\node[draw=poppurple, fill=poppurpleL, rounded corners, text width=3.1cm, align=center, minimum height=2.6cm] (C) at (10.8,0)
{\textbf{Sujet de discussion}\\[2pt] Consigne : \emph{Discutez...}\\ Plan : dialectique (thèse / antithèse / synthèse)};
\draw[-{Stealth}, thick, popmuted] (A) -- (B);
\draw[-{Stealth}, thick, popmuted] (B) -- (C);
\end{tikzpicture}
\legende{Les trois types de sujets de dissertation au Baccalauréat technique : consignes typiques et plans attendus.}
\end{popfigure}

\courssub{Présupposés et contenu latent du sujet}

Un sujet dit souvent plus qu'il ne semble dire. Le \cle{contenu manifeste} est ce qui est écrit explicitement ; le \cle{contenu latent} est ce qui est sous-entendu. Parmi ces sous-entendus, les \cle{présupposés} sont les idées admises sans être démontrées.

\begin{exemplebox}[Débusquer un présupposé]
Sujet : « L'école est-elle la \emph{seule} voie de la réussite ? »
\begin{itemize}
\item Contenu manifeste : on interroge le lien entre école et réussite.
\item Présupposé 1 : l'école est \emph{une} voie de la réussite (personne ne le nie).
\item Présupposé 2 : il existerait d'autres voies possibles (sinon la question n'aurait pas de sens).
\item Contenu latent : la réussite sociale au Cameroun se mesure souvent au diplôme — est-ce légitime ?
\end{itemize}
Repérer ces présupposés permet de construire les deux axes du devoir : oui, l'école est une voie majeure ; mais non, elle n'est pas la seule.
\end{exemplebox}

\courssub{Formuler la problématique}

\begin{definition}[La problématique]
La \cle{problématique} est la question centrale qui \emph{problématise} le sujet, c'est-à-dire qui met en évidence la difficulté, la tension entre les termes du sujet. Elle guide tout le devoir : chaque partie du plan doit y répondre partiellement, et la conclusion y répond globalement.
\end{definition}

Une bonne problématique naît de la \textbf{tension} entre les mots-clés. Dans « L'école est-elle la seule voie de la réussite ? », la tension oppose \emph{école} (voie institutionnelle) et \emph{seule} (exclusivité). La problématique peut alors être : « Si l'école demeure une voie privilégiée de la réussite, peut-on réussir sa vie sans elle, et à quelles conditions ? »

\begin{popfigure}
\begin{tikzpicture}[font=\small]
\foreach \y/\w/\txt/\col in {0/10.5/{Sujet : lire et reformuler}/popblue,
 -1.3/8.5/{Mots-clés : souligner les termes importants}/poporange,
 -2.6/6.5/{Définitions : préciser le sens de chaque terme}/poppurple,
 -3.9/4.7/{Tension : repérer l'opposition entre les termes}/popgreen,
 -5.2/2.9/{\textbf{Problématique}}/popgold}{
 \node[draw=\col, fill=\col!25, text width=\w cm, align=center, rounded corners, minimum height=0.95cm] at (0,\y) {\txt};
}
\draw[-{Stealth}, very thick, popmuted] (0,-0.55) -- (0,-0.85);
\draw[-{Stealth}, very thick, popmuted] (0,-1.85) -- (0,-2.15);
\draw[-{Stealth}, very thick, popmuted] (0,-3.15) -- (0,-3.45);
\draw[-{Stealth}, very thick, popmuted] (0,-4.45) -- (0,-4.75);
\end{tikzpicture}
\legende{Schéma en entonnoir de l'analyse du sujet : on part du sujet (large) et on aboutit à la problématique (étroite et précise).}
\end{popfigure}

\courssub{Étude de cas guidée : un sujet de Baccalauréat technique}

\begin{exoresolu}[Analyse complète du sujet « L'école est-elle la seule voie de la réussite ? »]
Appliquons la méthode en cinq étapes à ce sujet classique du Baccalauréat technique camerounais.
\tcblower
\textbf{Étape 1 — Lire.} Trois lectures silencieuses. Reformulation : « Peut-on réussir sa vie uniquement grâce à l'école, ou existe-t-il d'autres chemins ? »

\textbf{Étape 2 — Souligner.} Mots-clés : \emph{école}, \emph{seule}, \emph{voie}, \emph{réussite}. Consigne : question directe (« est-elle... ? ») : sujet de réflexion invitant à prendre position.

\textbf{Étape 3 — Définir.}
\begin{itemize}
\item \emph{École} : institution de formation dispensant savoirs, diplômes et socialisation.
\item \emph{Seule} : adverbe restrictif, cœur du sujet ; il exclut toute autre possibilité.
\item \emph{Voie} : chemin, moyen d'accès.
\item \emph{Réussite} : accomplissement de soi (social, professionnel, financier, moral).
\end{itemize}

\textbf{Étape 4 — Relier.} Le sujet affirme implicitement que l'école est une voie de réussite, mais interroge son exclusivité. Tension : école \emph{contre} autres voies (formation professionnelle, entrepreneuriat, talent, expérience).

\textbf{Étape 5 — Questionner.} Problématique : « L'école est-elle réellement le seul chemin qui mène à la réussite, ou d'autres voies peuvent-elles y conduire ? »

\textbf{Axes du plan détaillé (dialectique) :}
\begin{itemize}
\item \textbf{I.} L'école, voie privilégiée de la réussite : elle transmet les savoirs (argument de logique ; exemple : les concours administratifs exigent des diplômes) ; elle donne des repères et des réseaux (argument d'expérience).
\item \textbf{II.} Mais l'école n'est pas la seule voie : la formation technique et l'apprentissage (exemple : les artisans et techniciens autodidactes prospères) ; l'entrepreneuriat et le talent (exemple : des artistes et commerçants camerounais reconnus).
\item \textbf{III.} Synthèse : la réussite dépend surtout de la volonté, du travail et de l'adaptation ; l'école reste un atout majeur, mais pas une condition absolue.
\end{itemize}
\end{exoresolu}

\courssub{Les erreurs à éviter}

\begin{attention}[Le mot-clé négligé]
L'erreur la plus grave est de répondre à côté parce qu'un mot-clé a été négligé. Dans notre sujet, le candidat qui oublie « \emph{seule} » traitera « les avantages de l'école » : c'est un hors-sujet partiel, lourdement sanctionné. Chaque mot du sujet compte, surtout les petits mots (\emph{seule, toujours, vraiment, forcément}).
\end{attention}

\begin{attention}[Paraphrase et problématique trop large]
Deux autres écueils fréquents :
\begin{itemize}
\item \textbf{La paraphrase du sujet} : répéter le sujet avec d'autres mots dans l'introduction sans l'analyser. L'introduction doit \emph{expliquer} les termes, pas les recopier.
\item \textbf{La problématique trop large} : « Parlons de l'école » n'est pas une problématique. La question doit être précise, issue de la tension entre les termes du sujet, et traitable en deux ou trois parties.
\end{itemize}
\end{attention}

\begin{savaistu}[Le hors-sujet, première cause d'échec]
Les rapports du jury du Baccalauréat au Cameroun le rappellent chaque session : dans les copies faibles de dissertation, le hors-sujet et le sujet mal compris arrivent en tête des causes d'échec, devant les fautes de langue. On estime couramment qu'environ une copie faible sur trois doit sa note insuffisante à une mauvaise analyse du sujet. Autrement dit : dix minutes bien employées à analyser le sujet valent mieux qu'une heure de rédaction précipitée. Les correcteurs répètent la même consigne : « Lisez le sujet, relisez le sujet, puis seulement écrivez. »
\end{savaistu}

\begin{exercice}[À toi de jouer]
Analyse le sujet suivant en appliquant les cinq étapes de la méthode : « Le travail est-il le seul facteur de l'épanouissement de l'homme ? » Tu souligneras les mots-clés, définiras chacun d'eux, dégageras les présupposés, puis formuleras une problématique et proposeras les axes d'un plan dialectique (thèse, antithèse, synthèse).
\end{exercice}

\begin{corrige}[Exercice]
\textbf{Mots-clés :} \emph{travail}, \emph{seul}, \emph{facteur}, \emph{épanouissement}, \emph{homme}.
\textbf{Définitions :} travail = activité rémunérée ou effort productif ; seul = exclusivité (mot restrictif, cœur du sujet) ; facteur = élément qui contribue à un résultat ; épanouissement = plein développement de la personne (bonheur, accomplissement) ; homme = l'être humain en général.
\textbf{Présupposés :} le travail contribue à l'épanouissement ; d'autres facteurs pourraient exister (famille, loisirs, foi, amitié, santé).
\textbf{Problématique :} « Le travail suffit-il à épanouir l'homme, ou d'autres dimensions de la vie y contribuent-elles tout autant ? »
\textbf{Axes possibles (plan dialectique) :} 
\textbf{I. (Thèse)} Le travail, facteur essentiel d'épanouissement (dignité, revenu, utilité sociale). 
\textbf{II. (Antithèse)} Mais d'autres facteurs complètent le travail (affection, loisirs, engagement spirituel). 
\textbf{III. (Synthèse)} L'épanouissement naît de l'équilibre entre travail et vie personnelle.
\end{corrige}

\begin{aretenir}[L'essentiel de la section]
\begin{itemize}
\item Le texte argumentatif = thèse + arguments + exemples + connecteurs, avec une visée persuasive.
\item Un bon argument = une idée + une justification + un exemple précis ; on distingue les arguments d'autorité, d'expérience, de logique et de valeur.
\item L'analyse du sujet suit cinq étapes : lire, souligner, définir, relier, questionner.
\item Il existe trois types de sujets : réflexion, commentaire, discussion ; la consigne (« Pensez-vous », « Commentez », « Discutez ») indique le plan attendu.
\item La problématique est la question centrale née de la tension entre les termes du sujet ; elle guide tout le devoir.
\item Erreurs fatales : le hors-sujet (mot-clé négligé), la paraphrase, la problématique trop large.
\end{itemize}
\end{aretenir}

\coursec{Recherche des idées et élaboration du plan détaillé}

\courssub{De l'analyse du sujet à la génération d'idées : transition et rappel}

Après avoir disséqué le sujet, identifié ses termes clés, dégagé la tension qui l'habite et formulé une problématique précise (voir section précédente « Le texte argumentatif et l'analyse du sujet »), nous voici au cœur de l'atelier de la dissertation : la recherche des idées et la construction du plan détaillé. C'est l'étape où la réflexion devient matière, où l'intuition se structure en argumentation rigoureuse. Rappelons-le : au Probatoire comme au Baccalauréat technique, le correcteur lit d'abord l'introduction et le plan apparent ; ils conditionnent la première impression et souvent la note finale. Un plan clair, équilibré, progressif, c'est déjà la moitié de la réussite.

\begin{savaistu}[L'impression du correcteur]
Au Baccalauréat technique, le correcteur dispose de quelques minutes par copie. Il lit d'abord l'introduction (amorce, sujet, problématique, annonce du plan) et repère immédiatement la structure annoncée. Si le plan est lisible, cohérent et adapté au sujet, la copie part avec un capital de confiance. Si le plan est confus, déséquilibré ou inadapté (ex. plan thématique annoncé pour un sujet dialectique), la lecture suivante se fait dans la suspicion. Soigner son plan détaillé, c'est soigner sa première impression.
\end{savaistu}

\courssub{Techniques de recherche des idées : mobiliser son capital culturel et contextuel}

La page blanche fait peur. Pourtant, les idées ne tombent pas du ciel : elles sont déjà en vous, dans vos lectures, votre environnement, votre expérience. La recherche d'idées (ou \cle{inventio} chez les Anciens) est un travail actif de mobilisation et de sélection.

\begin{methode}[Quatre réservoirs pour alimenter sa réflexion]
\textbf{1. Le brainstorming (remue-méninges) pur.} Notez sans censure, pendant 5 à 7 minutes, tout ce qui vous vient à l'esprit à partir des mots-clés du sujet et de la problématique : mots, expressions, noms d'auteurs, titres d'œuvres, faits d'actualité, souvenirs personnels, proverbes. Ne triez pas encore. L'objectif : quantité et liberté.

\textbf{2. La mobilisation des lectures au programme.} \emph{Le Vieux nègre et la médaille} (Ferdinand Oyono) et \emph{Le Lion et la perle} (Wole Soyinka) sont des réservoirs d'arguments littéraires, historiques, sociologiques. Exemple : pour un sujet sur « Tradition et modernité en Afrique », \emph{Le Vieux nègre et la médaille} offre la figure de Meka (tradition trahie, collaboration coloniale, illusion des médailles) ; \emph{Le Lion et la perle} offre Baroka (ruse traditionnelle face à la modernité), Lakunle (modernité caricaturale), Sidi (enjeu de la beauté et du pouvoir).

\textbf{3. L'actualité camerounaise et africaine.} Politique (décentralisation, bilinguisme, jeunesse, emploi), société (urbain/rural, éducation, santé, réseaux sociaux), culture (patrimoine, langues nationales, musique, littérature). Un sujet sur « L'école face aux défis du numérique » trouve des exemples concrets dans la rentrée scolaire 2023-2024 au Cameroun, le programme « École numérique », la fracture numérique ville/campagne.

\textbf{4. L'expérience personnelle et l'observation vécue.} Votre stage en entreprise, votre vie de quartier, vos responsabilités familiales, vos engagements associatifs. Un sujet sur « Le travail et la dignité » s'enrichit de votre expérience de job d'été ou du métier de vos parents. L'argument d'autorité (« Moi, j'ai vu... ») a sa place s'il est présenté comme illustration, non comme preuve universelle.
\end{methode}

\begin{exemplebox}[Brainstorming appliqué : sujet « La jeunesse africaine face à l'émigration »]
\emph{Problématique :} L'émigration est-elle une fatalité ou un choix lucide pour la jeunesse africaine ?
\begin{itemize}
\item \textbf{Lectures :} \emph{Le Vieux nègre et la médaille} (Meka, figure de l'aliénation coloniale, non de l'émigration mais du déracinement intérieur) ; \emph{Le Lion et la perle} (Lakunle veut « moderniser » le village, mais reste ; Baroka gère le pouvoir local).
\item \textbf{Actualité :} Tragédies en Méditerranée, « bush falling » au Cameroun, politiques de retour (programme Diaspora), remises de fonds (transferts financiers), brain drain vs brain gain.
\item \textbf{Expérience :} Cousin parti au Qatar, amie infirmière en Allemagne, frère revenu au pays créer une start-up agro, voisins dont le fils a disparu en Libye.
\item \textbf{Concepts :} Poussée/attraction (push/pull), capital humain, remittances, double absence (Sayad), citoyenneté transnationale.
\end{itemize}
\end{exemplebox}

\courssub{Trier, regrouper, hiérarchiser : de la masse à la structure}

Une fois le réservoir rempli, il faut filtrer. Toutes les idées ne sont pas bonnes à prendre.

\begin{propriete}[Critères de sélection des idées]
Une idée est \textbf{pertinente} si elle répond directement à la problématique (pas au sujet seul). Elle est \textbf{opérante} si elle peut être argumentée (explication + exemple). Elle est \textbf{non redondante} si elle n'apporte pas la même chose qu'une autre idée déjà retenue (même angle, même exemple). Elle est \textbf{représentative} si elle couvre un aspect significatif du débat (pas un détail anecdotique).
\end{propriete}

\begin{methode}[Tri et regroupement en trois temps]
\textbf{1. Élimination.} Barrez les idées hors-sujet, les répétitions, les jugements de valeur non argumentables (« C'est bien », « C'est mal »), les exemples trop personnels sans portée générale.

\textbf{2. Regroupement thématique.} Collez des étiquettes (post-it, surlignage) : idées qui parlent de \emph{causes}, idées qui parlent de \emph{conséquences}, idées qui parlent d'\emph{acteurs}, idées qui parlent de \emph{solutions}, idées \emph{pour}, idées \emph{contre}. Vous verrez émerger 3 à 5 grands pôles.

\textbf{3. Hiérarchisation et ordonnancement.} Dans chaque pôle, choisissez l'argument fort (le plus général, le plus structurant) comme \cle{idée directrice} de sous-partie. Les autres deviennent arguments secondaires ou exemples. Numérotez dans l'ordre logique : chronologique, causal, du plus simple au plus complexe, du particulier au général, ou selon la dialectique (thèse / antithèse / synthèse).
\end{methode}

\begin{attention}[Piège fréquent : confondre exemple et argument]
Un exemple (fait précis, citation, statistique) \textbf{ne} remplace \textbf{pas} un argument (raisonnement général). Dans le plan détaillé, chaque sous-partie doit porter une \textbf{idée directrice formulée en phrase nominale ou verbale}, soutenue par \textbf{un ou deux arguments explicatifs} et \textbf{illustrée par un exemple concret}. Écrire seulement « Exemple : mon cousin au Qatar » dans le plan est insuffisant ; il faut écrire « L'émigration comme stratégie familiale de survie économique (argument) — Ex. : les transferts de fonds du cousin X qui financent la scolarité de ses neveux à Yaoundé ».
\end{attention}

\courssub{Les types de plans : choisir l'architecture adaptée au sujet}

Le plan n'est pas un moule unique. Il doit épouser la logique de la problématique. Trois grandes architectures dominent en dissertation au lycée technique.

\begin{definition}[Plan dialectique (thèse — antithèse — synthèse)]
Structure classique pour les sujets \textbf{polémiques, à tension binaire} (« Faut-il... ? », « Est-il vrai que... ? », « Pour ou contre... ? »).
\begin{itemize}
\item \textbf{Partie I (Thèse) :} Les arguments \textbf{pour} la proposition, les évidences, le sens commun, les avantages.
\item \textbf{Partie II (Antithèse) :} Les arguments \textbf{contre}, les limites, les contre-exemples, les risques, la complexité.
\item \textbf{Partie III (Synthèse) :} Dépassement de l'opposition : nuances, conditions, réconciliation partielle, nouvelle perspective, ouverture. \emph{Attention : la synthèse n'est pas un résumé, c'est une construction nouvelle.}
\end{itemize}
\end{definition}

\begin{definition}[Plan thématique (ou par aspects)]
Structure adaptée aux sujets \textbf{ouverts, d'analyse, d'inventaire} (« Quels sont les enjeux de... ? », « Analysez les formes de... », « Étudie la place de... »).
\begin{itemize}
\item Chaque partie explore une \textbf{dimension distincte} du sujet (économique, sociale, culturelle, politique, psychologique, historique...).
\item Pas d'opposition binaire, mais une \textbf{couverture exhaustive et ordonnée} du champ problématique.
\item Exemple : « Les défis de l'insertion professionnelle des jeunes diplômés au Cameroun » → I. Déficit de formation (adéquation formation/emploi) — II. Rigidité du marché du travail (secteur public saturé, secteur informel précaire) — III. Manque d'accompagnement (orientation, financement, mentorat).
\end{itemize}
\end{definition}

\begin{definition}[Plan analytique (constat — causes — conséquences — solutions)]
Structure idéale pour les sujets \textbf{problème-solution, diagnostic, prospective} (« Comment résoudre... ? », « Pourquoi... ? », « Quelles pistes pour... ? »).
\begin{itemize}
\item \textbf{Partie I : Constat / Diagnostic} (état des lieux chiffré, description de la situation).
\item \textbf{Partie II : Causes / Facteurs explicatifs} (structurels, conjoncturels, culturels, politiques).
\item \textbf{Partie III : Conséquences / Enjeux} (si rien ne change).
\item \textbf{Partie IV : Solutions / Pistes d'action} (court, moyen, long terme ; acteurs publics, privés, société civile).
\item \emph{Souvent réduit à trois parties en fusionnant Constat/Causes ou Conséquences/Solutions selon l'amplitude.}
\end{itemize}
\end{definition}

\begin{methode}[Comment choisir son type de plan ?]
\begin{enumerate}
\item \textbf{Lisez la problématique.} Contient-elle une opposition claire (pour/contre) ? → Dialectique. Demande-t-elle une analyse multidimensionnelle ? → Thématique. Pose-t-elle un problème à résoudre ? → Analytique.
\item \textbf{Regardez vos groupes d'idées.} S'organisent-ils naturellement en deux camps opposés ? → Dialectique. En facettes complémentaires ? → Thématique. En chaîne causale ? → Analytique.
\item \textbf{Testez l'annonce du plan.} « Nous verrons d'abord que... puis que... enfin que... » — La phrase sonne-t-elle juste ? Les transitions sont-elles logiques ?
\item \textbf{Équilibrez.} Trois parties de poids comparable (environ 300-400 mots chacune au Bac). Deux sous-parties par partie (parfois trois si dense).
\end{enumerate}
\end{methode}

\begin{popfigure}
\begin{tikzpicture}[
  block/.style={rectangle, draw=popblue, thick, fill=popblueL, rounded corners=6pt, minimum width=4.5cm, minimum height=1.8cm, align=center, font=\small},
  arrow/.style={->, >=Stealth, thick, popdark},
  trans/.style={font=\footnotesize, text=poppurple, align=center}
]
% Three main blocks
\node[block, align=center] (thesis) at (0,0){\textbf{THÈSE}\\\textit{Partie I}\\\textit{Arguments pour}\\\textit{le sens commun,}\\\textit{les avantages}};
\node[block, align=center] (antithesis) at (5.5,0){\textbf{ANTITHÈSE}\\\textit{Partie II}\\\textit{Arguments contre,}\\\textit{limites, complexité}};
\node[block, align=center] (synthesis) at (11,0){\textbf{SYNTHÈSE}\\\textit{Partie III}\\\textit{Dépassement,}\\\textit{nuances, ouverture}};
% Arrows between blocks
\draw[arrow] (thesis.east) -- node[trans, above] {Transition :\newline « Cependant / Pourtant »} (antithesis.west);
\draw[arrow] (antithesis.east) -- node[trans, above] {Transition :\newline « Au-delà / En définitive »} (synthesis.west);
% Internal structure hints
\node[below=0.3cm of thesis, font=\tiny, text=popmuted] {2 sous-parties max};
\node[below=0.3cm of antithesis, font=\tiny, text=popmuted] {2 sous-parties max};
\node[below=0.3cm of synthesis, font=\tiny, text=popmuted] {2 sous-parties max};
% Vertical flow hint
\draw[arrow, popgreen] (thesis.south) -- ++(0,-0.8) -| node[trans, near start, left] {Idées directrices} ++(-2.5,-1.5);
\draw[arrow, popgreen] (antithesis.south) -- ++(0,-0.8) -| node[trans, near start, right] {Arguments + Exemples} ++(2.5,-1.5);
\draw[arrow, popgreen] (synthesis.south) -- ++(0,-0.8) -- node[trans, below] {Nouvelle thèse} ++(0,-1.5);
\end{tikzpicture}
\legende{Architecture du plan dialectique : trois blocs articulés par des transitions logiques. Chaque partie contient 2 sous-parties (idée directrice + argumentation + exemple).}
\end{popfigure}

\courssub{Construire le plan détaillé : de l'architecture à la fiche de travail}

Le \cle{plan détaillé} n'est pas le plan apparent (celui qu'on annonce dans l'introduction : « Nous étudierons d'abord..., puis..., enfin... »). C'est sa \textbf{version développée, privée, de travail}, où chaque sous-partie est rédigée en phrase complète (idée directrice), accompagnée de ses arguments explicatifs et de ses exemples précis. C'est la \cle{feuille de route} qui garantit qu'on ne s'égare pas en rédaction.

\begin{definition}[Plan détaillé]
Document de travail (souvent au brouillon, parfois remis au professeur) où la dissertation est entièrement \textbf{pré-rédigée en structure} : introduction complète (4 temps), développement avec pour chaque sous-partie (idée directrice, arguments, exemples, connecteurs logiques), transitions inter-parties rédigées, conclusion (bilan + ouverture). Il permet de vérifier la cohérence, l'équilibre, la progression avant l'écriture finale.
\end{definition}

\begin{popfigure}
\begin{tikzpicture}[
  cell/.style={rectangle, draw=popline, minimum height=0.9cm, align=left, font=\small},
  header/.style={rectangle, draw=popblue, fill=popblueL, minimum height=1cm, align=center, font=\small\bfseries, text=popdark},
  title/.style={font=\bfseries, text=popdark}
]
% Column widths
\def\colA{2.2cm}
\def\colB{3.5cm}
\def\colC{4.5cm}
\def\colD{4.5cm}
% Header row
\node[header, minimum width=\colA] (hA) at (0,0) {Partie};
\node[header, minimum width=\colB, align=center] (hB) at (\colA,0){Idée directrice\\(sous-partie)};
\node[header, minimum width=\colC, align=center] (hC) at (\colA+\colB,0){Arguments\\explicatifs};
\node[header, minimum width=\colD, align=center] (hD) at (\colA+\colB+\colC,0){Exemples\\précis (références)};
% Row 1: Intro
\node[cell, minimum width=\colA, fill=popgreenL!30] (r1A) at (0,-1) {\textbf{INTRO}};
\node[cell, minimum width=\colB] (r1B) at (\colA,-1) {Amorce : citation / fait marquant};
\node[cell, minimum width=\colC] (r1C) at (\colA+\colB,-1) {Présentation du sujet + définition termes clés};
\node[cell, minimum width=\colD] (r1D) at (\colA+\colB+\colC,-1) {Problématique + Annonce du plan};
% Row 2: Partie I
\node[cell, minimum width=\colA, fill=popblueL!30] (r2A) at (0,-2) {\textbf{I. Thèse}};
\node[cell, minimum width=\colB] (r2B) at (\colA,-2) {A. L'émigration comme réponse à la pauvreté};
\node[cell, minimum width=\colC] (r2C) at (\colA+\colB,-2) {Poussée économique : chômage, salaires bas, absence de perspectives};
\node[cell, minimum width=\colD] (r2D) at (\colA+\colB+\colC,-2) {Statistiques INS 2023 : 68\% jeunes au chômage ; Témoignage cousin au Qatar};
\node[cell, minimum width=\colA] (r3A) at (0,-3) {};
\node[cell, minimum width=\colB] (r3B) at (\colA,-3) {B. L'émigration comme devoir familial};
\node[cell, minimum width=\colC] (r3C) at (\colA+\colB,-3) {Solidarité lignagère, envoi de fonds (remittances), prestige social};
\node[cell, minimum width=\colD] (r3D) at (\colA+\colB+\colC,-3) {Étude Banque Mondiale 2022 : transferts = 3\% PIB Cameroun ; Roman « Les larmes de l'émigré »};
% Transition I-II
\node[cell, minimum width=\colA+\colB+\colC+\colD, fill=poporangeL!30, font=\itshape\small] (trans1) at (\colA+\colB+\colC+\colD/2,-3.7) {Transition I $\to$ II : « Cependant, cette fuite en avant masque des réalités dramatiques et des coûts humains insoutenables. »};
% Row 4: Partie II
\node[cell, minimum width=\colA, fill=poppurpleL!30] (r4A) at (0,-4.4) {\textbf{II. Antithèse}};
\node[cell, minimum width=\colB] (r4B) at (\colA,-4.4) {A. Les périls de la route et l'exploitation};
\node[cell, minimum width=\colC] (r4C) at (\colA+\colB,-4.4) {Traite, esclavage moderne en Libye, noyades en Méditerranée, travail forcé};
\node[cell, minimum width=\colD] (r4D) at (\colA+\colB+\colC,-4.4) {Rapport OIM 2023 : 12 000 migrants bloqués au Niger ; Témoignage rescapé dans « Jeune Afrique »};
\node[cell, minimum width=\colA] (r5A) at (0,-5.4) {};
\node[cell, minimum width=\colB] (r5B) at (\colA,-5.4) {B. Le coût pour le pays d'origine (brain drain)};
\node[cell, minimum width=\colC] (r5C) at (\colA+\colB,-5.4) {Perte de capital humain formé, vieillissement secteur santé/éducation, dépendance aux transferts};
\node[cell, minimum width=\colD] (r5D) at (\colA+\colB+\colC,-5.4) {Médecins camerounais en France : 1 pour 5 formés ; Fermeture de centres de santé ruraux};
% Transition II-III
\node[cell, minimum width=\colA+\colB+\colC+\colD, fill=poporangeL!30, font=\itshape\small] (trans2) at (\colA+\colB+\colC+\colD/2,-6.1) {Transition II $\to$ III : « Dès lors, l'enjeu n'est plus de partir ou rester, mais de transformer le départ en atout pour le territoire. »};
% Row 6: Partie III
\node[cell, minimum width=\colA, fill=poppinkL!30] (r6A) at (0,-6.8) {\textbf{III. Synthèse}};
\node[cell, minimum width=\colB] (r6B) at (\colA,-6.8) {A. Politiques de retour et diaspora constructive};
\node[cell, minimum width=\colC] (r6C) at (\colA+\colB,-6.8) {Programmes « Retour au pays », incitations fiscales, réseaux diaspora, transfert de compétences};
\node[cell, minimum width=\colD] (r6D) at (\colA+\colB+\colC,-6.8) {Ex. : startup « AgriTech » fondée par diplômé revenu de France ; Programme MINJEC 2024};
\node[cell, minimum width=\colA] (r7A) at (0,-7.8) {};
\node[cell, minimum width=\colB] (r7B) at (\colA,-7.8) {B. Développement endogène comme alternative crédible};
\node[cell, minimum width=\colC] (r7C) at (\colA+\colB,-7.8) {Investissement local, formation adaptée, économie verte, valorisation terroirs};
\node[cell, minimum width=\colD] (r7D) at (\colA+\colB+\colC,-7.8) {Modèle « Songhaï » (Bénin) ; Coopératives jeunes agriculteurs Nord-Cameroun};
% Row 8: Conclusion
\node[cell, minimum width=\colA, fill=popgoldL!30] (r8A) at (0,-8.8) {\textbf{CONCL.}};
\node[cell, minimum width=\colB] (r8B) at (\colA,-8.8) {Bilan : ni fatalité, ni solution unique};
\node[cell, minimum width=\colC] (r8C) at (\colA+\colB,-8.8) {L'émigration révèle l'urgence de politiques de jeunesse intégrées};
\node[cell, minimum width=\colD] (r8D) at (\colA+\colB+\colC,-8.8) {Ouverture : « Et si la vraie patrie, c'était l'avenir qu'on construit ensemble ? » (Soyinka)};
% Grid lines
\draw[popline] (0,0.5) -- (\colA+\colB+\colC+\colD,0.5);
\draw[popline] (0,-9.3) -- (\colA+\colB+\colC+\colD,-9.3);
\foreach \x in {0,\colA,\colA+\colB,\colA+\colB+\colC,\colA+\colB+\colC+\colD} {
  \draw[popline] (\x,0.5) -- (\x,-9.3);
}
\foreach \y in {0,-1,-2,-3,-3.7,-4.4,-5.4,-6.1,-6.8,-7.8,-8.8,-9.3} {
  \draw[popline] (0,\y) -- (\colA+\colB+\colC+\colD,\y);
}
\end{tikzpicture}
\legende{Modèle de fiche de plan détaillé (sujet : « La jeunesse africaine face à l'émigration »). Colonnes : Partie / Idée directrice / Arguments / Exemples. Transitions inter-parties rédigées. Couleurs : intro/concl (vert/or), thèse (bleu), antithèse (violet), synthèse (rose), transitions (orange).}
\end{popfigure}

\courssub{L'introduction en quatre temps : la porte d'entrée de la copie}

L'introduction se rédige \textbf{en dernier} (quand le plan détaillé est figé), mais elle se conçoit \textbf{en premier} dans la logique. Elle suit un rituel immuable en quatre temps.

\begin{methode}[Introduction en quatre temps (le « entonnoir »)]
\textbf{Temps 1 — Amorce (accroche).} Une phrase qui \textbf{capte l'attention} et ancre le sujet dans le concret ou l'universel. Trois types efficaces :
\begin{itemize}
\item \emph{Citation d'auteur} (africain de préférence) : « L'exil est une blessure qui ne cicatrise jamais » (Ken Bugul).
\item \emph{Fait d'actualité / statistique choc} : « En 2023, plus de 3 000 migrants ont péri en Méditerranée selon l'OIM. »
\item \emph{Proverbe / adage camerounais} : « Celui qui ne sait pas d'où il vient ne sait pas où il va. »
\end{itemize}
\emph{Évitez :} « De tout temps, les hommes ont migré » (banal), « Je vais vous parler de... » (scolaire).

\textbf{Temps 2 — Présentation du sujet / Définition des termes.} Reprenez les mots-clés du sujet, définissez-les brièvement (sémantique), contextualisez (spatio-temporel : Cameroun, Afrique, aujourd'hui). Montrez que vous avez compris l'enjeu.

\textbf{Temps 3 — Problématique.} La question \textbf{unique, précise, tensionnelle} qui guide toute la dissertation. Elle ne se termine pas par un point d'interrogation rhétorique multiple. Une seule question, avec un « comment », « pourquoi », « dans quelle mesure », « faut-il ».

\textbf{Temps 4 — Annonce du plan.} Phrase unique, structurée : « Pour y répondre, nous analyserons d'abord [Partie I], puis nous examinerons [Partie II], avant de montrer enfin [Partie III]. » Les verbes varient : analyser, étudier, montrer, démontrer, mettre en lumière, interroger. \textbf{Le plan annoncé doit correspondre exactement au plan détaillé.}
\end{methode}

\begin{exoresolu}[Introduction type — Sujet : « La jeunesse africaine face à l'émigration : fatalité ou choix ? »]
\textbf{Amorce :} « Partir, c'est mourir un peu », écrivait Alphonse de Lamartine ; pour des milliers de jeunes Africains aujourd'hui, partir, c'est risquer de mourir tout court.

\textbf{Présentation :} L'émigration juvénile, phénomène massif au Cameroun comme dans le Golfe de Guinée, touche une génération née dans les années 1990-2000, confrontée au chômage de masse (près de 70\% selon l'INS), à la précarité de l'emploi informel et à l'attrait d'un « Eldorado » européen ou moyen-oriental relayé par les réseaux sociaux.

\textbf{Problématique :} Dans quelle mesure l'émigration constitue-t-elle une réponse contrainte aux défaillances structurelles, plutôt qu'un choix libre et éclairé de la jeunesse africaine ?

\textbf{Annonce du plan :} Nous montrerons d'abord que l'émigration apparaît comme une stratégie de survie économique et familiale (I), avant d'analyser les coûts humains dramatiques et le prélèvement de compétences qu'elle entraîne (II), pour enfin esquisser les conditions d'une circulation des talents profitable aux deux rives (III).
\tcblower
\textbf{Analyse :} Amorce littéraire + choc réel. Présentation contextuelle (chiffre INS, génération). Problématique unique, tension « contrainte / choix ». Annonce plan dialectique (I = thèse, II = antithèse, III = synthèse) avec verbes variés.
\end{exoresolu}

\courssub{Le développement : parties, sous-parties, arguments, exemples, transitions}

Chaque partie du développement (I, II, III) contient \textbf{deux sous-parties} (A et B), parfois trois si la densité l'exige (mais attention au temps). Chaque sous-partie suit un micro-schéma invariant.

\begin{propriete}[Anatomie d'une sous-partie (paragraphe développé)]
\begin{enumerate}
\item \textbf{Phrase d'appel (idée directrice) :} Affirme le point principal en une phrase nominale ou verbale claire. Ex. : « L'émigration répond d'abord à une contrainte économique structurelle. »
\item \textbf{Argumentation (explication) :} 2-3 phrases qui \textbf{développent le raisonnement} : causes, mécanismes, logiques, concepts. C'est le « pourquoi » et le « comment ». Utilisez les connecteurs logiques (\emph{donc, par conséquent, en effet, notamment, ainsi}).
\item \textbf{Exemple(s) précis(s) :} 1 à 2 illustrations concrètes (chiffre daté, fait d'actualité nommé, référence littéraire précise, expérience vécue cadrée). L'exemple \textbf{valide} l'argument, ne le remplace pas.
\item \textbf{Mini-transition (facultatif) :} Phrase de liaison vers la sous-partie suivante : « Mais cette logique économique occulte une dimension tout aussi cruciale : le devoir de solidarité familiale. »
\end{enumerate}
\end{propriete}

\begin{attention}[Erreur fatale : annoncer un plan dialectique et rédiger un plan thématique]
C'est l'erreur n°1 au Bac. Le sujet appelle une tension (Pour/Contre), l'annonce dit « Nous verrons les avantages, puis les inconvénients, puis la synthèse ». Mais dans le développement, l'élève écrit : I. Causes économiques — II. Causes sociales — III. Solutions politiques. \textbf{Incohérence totale.} Le correcteur sanctionne : hors-sujet partiel, note plafonnée. \textbf{Vérifiez l'adéquation Plan annoncé / Plan exécuté} sur votre fiche détaillée avant de recopier.
\end{attention}

\courssub{Les transitions : l'articulation qui fait tenir l'édifice}

Une transition n'est pas un simple « Passons maintenant à... ». C'est une \textbf{phrase charnière} qui :
\begin{itemize}
\item \textbf{Rappelle} l'idée maîtresse de la partie achevée (résumé en une clause).
\item \textbf{Annonce} l'idée maîtresse de la partie qui commence (anticipation).
\item \textbf{Justifie} le passage : opposition (mais, cependant, pourtant), approfondissement (plus encore, au-delà), conséquence (par conséquent, dès lors), nuance (toutefois, néanmoins).
\end{itemize}

\begin{exemplebox}[Transitions modèles selon l'architecture]
\textbf{Dialectique I $\to$ II (opposition) :} « Si l'émigration apparaît comme une bouée de sauvetage économique pour les familles (rappel I), elle révèle pourtant un revers tragique : l'exposition aux réseaux criminels et la saignée des compétences nationales (annonce II). »

\textbf{Dialectique II $\to$ III (dépassement) :} « Au-delà du constat d'échec et de perte (rappel II), se dessine une voie médiane : transformer la diaspora en levier de développement par le retour volontaire et l'engagement à distance (annonce III). »

\textbf{Thématique I $\to$ II (complémentarité) :} « Après avoir analysé les freins structurels du marché du travail (rappel I), il convient d'examiner l'inadéquation des formations proposées, second pilier du chômage des diplômés (annonce II). »

\textbf{Analytique Causes $\to$ Solutions (logique) :} « Puisque les causes sont identifiées — précarité, désadéquation, manque d'accompagnement (rappel I-II) —, les réponses doivent s'articuler sur ces trois leviers simultanément (annonce III). »
\end{exemplebox}

\courssub{La conclusion : bilan et ouverture maîtrisée}

La conclusion n'est pas une répétition de l'introduction. Elle a deux temps distincts.

\begin{methode}[Conclusion en deux temps]
\textbf{Temps 1 — Bilan (synthèse répondant à la problématique).} Reprenez la problématique, donnez la \textbf{réponse nuancée} qui émerge de votre démonstration. Pas de « En conclusion, on peut dire que... ». Écrivez : « L'émigration n'est ni une fatalité pure, ni un choix pleinement libre : elle est une réponse contrainte aux défaillances structurelles, qui peut devenir opportunité si des politiques publiques audacieuses transforment le départ en circulation des compétences. »

\textbf{Temps 2 — Ouverture (perspective, horizon).} Une phrase qui \textbf{élargit} sans partir en hors-sujet. Trois pistes :
\begin{itemize}
\item \emph{Extension spatiale/temporelle :} « La question se pose aussi pour la jeunesse européenne face au vieillissement démographique. »
\item \emph{Renvoi littéraire/philosophique :} « Comme l'écrit Wole Soyinka dans \emph{Le Lion et la perle}, "le vieux monde et le nouveau doivent danser ensemble" : l'avenir de l'Afrique se jouera dans cette chorégraphie. »
\item \emph{Question prospective :} « Reste à savoir si la ZLECAf et la libre circulation des personnes en Afrique sauront transformer l'émigration subie en mobilité choisie. »
\end{itemize}
\emph{Interdit :} « J'espère que... », « À nous de... », « Le gouvernement doit... » (injonction morale hors champ argumentatif).
\end{methode}

\courssub{TP guidé : construction collective d'un plan détaillé au tableau}

\begin{experience}[TP : « De la problématique au plan détaillé complet » — 50 min]
\textbf{Matériel :} Tableau (4 zones : Brainstorming / Tri / Plan apparent / Plan détaillé), feutres 4 couleurs, post-it, sujet imprimé.
\textbf{Sujet :} « L'école technique peut-elle répondre aux défis de l'insertion professionnelle au Cameroun ? »

\textbf{Phase 1 — Brainstorming collectif (10 min).} Élèves au tableau, professeur scribe. Tous les idées acceptées. Résultat attendu : chômage, inadéquation formation/emploi, manque d'équipements, stigmatisation voie technique, partenariats entreprises, stages, entrepreneuriat, réussite pays asiatiques, rôle parents, orientation.

\textbf{Phase 2 — Tri et regroupement (10 min).} Professeur guide : « Quelle idée va avec laquelle ? » → 3 pôles émergent : (A) Atouts de la voie technique (insertion rapide, pratique) — (B) Obstacles structurels (équipements, image, stage) — (C) Leviers de transformation (partenariats, entrepreneuriat, orientation). \textbf{Choix du plan :} Dialectique (I. Atouts / II. Limites / III. Pistes de refonte) OU Analytique (I. Constat / II. Causes / III. Solutions). Vote à main levée : plan dialectique retenu.

\textbf{Phase 3 — Rédaction du plan détaillé (20 min).} Par groupes de 4, chaque groupe remplit une colonne du tableau (Partie / Idée directrice / Argument / Exemple) pour une sous-partie assignée. Mise en commun au tableau, correction collective, harmonisation des formulations, rédaction des deux transitions I$\to$II et II$\to$III. Professeur valide la fiche finale photographiée et diffusée.

\textbf{Phase 4 — Auto-évaluation (5 min).} Grille d'évaluation par les pairs (voir boîte méthode ci-dessous).
\end{experience}

\begin{methode}[Grille d'évaluation par les pairs (auto-correction guidée)]
Chaque groupe échange sa fiche de plan détaillé avec un autre groupe. Notation sur 4 critères (0-5 pts chacun) :
\begin{center}
\begin{tabularx}{\linewidth}{|l|X|}\hline
\textbf{Critère} & \textbf{Indicateurs de réussite} \\ \hline
\cle{Pertinence} & Chaque idée directrice répond \textbf{directement} à la problématique. Pas de hors-sujet. \\ \hline
\cle{Cohérence} & Logique interne respectée (dialectique : thèse/antithèse/synthèse bien distinctes). Arguments \textbf{soutiennent} l'idée directrice. Exemples \textbf{illustrent} l'argument. \\ \hline
\cle{Progression} & Ordre logique : du plus évident au plus complexe, du général au particulier, ou thèse $\to$ antithèse $\to$ synthèse. Transitions \textbf{rédigées} et justifiées. \\ \hline
\cle{Équilibre} & 3 parties de poids comparable. 2 sous-parties par partie (max 3). Introduction complète (4 temps). Conclusion (bilan + ouverture). \\ \hline
\end{tabularx}
\end{center}
\textbf{Total /20.} Note $<$ 12 : remédiation obligatoire (retravail en binôme avec fiche corrigée professeur). Note $\ge$ 12 : passage à la rédaction complète en devoir maison.
\end{methode}

\courssub{Confrontation, amélioration, remédiation : le cycle de l'excellence}

La dissertation ne s'apprend qu'en \textbf{écrivant, confrontant, corrigeant, réécrivant}. Le plan détaillé est l'outil pivot de ce cycle.

\begin{experience}[Séance de confrontation et remédiation — 1h]
\textbf{Dispositif :} Chaque élève a produit un plan détaillé sur un sujet de Bac blanc (ex. « La technologie améliore-t-elle la condition humaine ? »). Anonymat : numéros au lieu des noms.
\begin{enumerate}
\item \textbf{Lecture croisée (15 min).} Chaque élève reçoit 2 plans anonymes. Il les évalue avec la grille (méthode ci-dessus), note les forces/faiblesses sur post-it.
\item \textbf{Mise en commun (20 min).} Professeur projette 3 plans types (bon / moyen / faible) — anonymisés. Analyse collective : « Pourquoi ce plan est-il noté 16 ? » (cohérence, exemples datés, transitions rédigées). « Pourquoi celui-ci 8 ? » (incohérence plan annoncé/réalisé, exemples vagues, pas de synthèse).
\item \textbf{Remédiation ciblée (15 min).} Groupes de besoin :
\begin{itemize}
\item Groupe A (problème pertinence) : Re-travail sur l'analyse de sujet / problématique.
\item Groupe B (problème structure) : Exercice « Reconstruis le plan à partir d'un développement brouillon ».
\item Groupe C (problème exemples) : Atelier « Cherche 3 exemples datés pour cet argument ».
\end{itemize}
\item \textbf{Réécriture individuelle (10 min).} Chaque élève corrige \textbf{son propre plan} à la lumière des retours. Version finale remise pour notation formative.
\end{enumerate}
\end{experience}

\begin{aretenir}[Check-list finale avant de lancer la rédaction]
\begin{itemize}
\item[$\square$] Problématique unique, tensionnelle, notée en haut de la fiche.
\item[$\square$] Type de plan choisi \textbf{explicitement} (dialectique / thématique / analytique) et justifié.
\item[$\square$] 3 parties équilibrées, 2 sous-parties chacune (6 idées directrices au total).
\item[$\square$] Chaque sous-partie : \textbf{Idée directrice (phrase) + Argument(s) + Exemple(s) précis}.
\item[$\square$] Transitions I$\to$II et II$\to$III \textbf{rédigées en phrases complètes}.
\item[$\square$] Introduction : 4 temps (amorce, présentation, problématique, annonce plan) — annonce \textbf{identique} au plan détaillé.
\item[$\square$] Conclusion : Bilan (réponse à la problématique) + Ouverture (littéraire / prospective / élargie).
\item[$\square$] Relecture : cohérence terminologique (mots-clés du sujet repris), connecteurs logiques variés, pas de répétitions.
\end{itemize}
\end{aretenir}

\begin{exoresolu}[Plan détaillé complet — Sujet Bac Tech 2023 : « Le numérique : chance ou menace pour la jeunesse camerounaise ? »]
\textbf{Problématique :} Le numérique est-il un levier d'émancipation socio-économique pour la jeunesse camerounaise ou un facteur d'aliénation et de précarisation accrue ?

\textbf{Plan dialectique choisi.}

\textbf{INTRODUCTION}
Amorce : « Le numérique n'est pas un outil, c'est un environnement » (Michel Serres) — au Cameroun, 12 millions d'internautes en 2023 (ART), dont 60\% ont moins de 25 ans.
Présentation : Pénétration mobile, Facebook/WhatsApp omniprésents, e-gouvernement, fintech, mais aussi cybercriminalité (« feyman »), désinformation, addiction.
Problématique : (ci-dessus)
Annonce : Nous analyserons d'abord les potentialités d'inclusion et d'innovation (I), puis les risques d'exclusion et de dérive (II), pour enfin dégager les conditions d'un numérique souverain et éthique (III).

\textbf{I. LE NUMÉRIQUE COMME LEVIER D'INCLUSION ET D'INNOVATION (Thèse)}
A. Accès à l'information et formation autonome
— Argument : Démocratisation du savoir, MOOC, YouTube éducatif, contournement manuels obsolètes.
— Ex. : Plateforme « Kola Academy » (camerounaise), 50 000 inscrits 2023 ; lycéens ruraux préparent le Bac via smartphone.
B. Entrepreneuriat et insertion par la tech
— Argument : Faible barrière à l'entrée, marché local à digitaliser, fintech, e-commerce, agritech.
— Ex. : Startup « GiftedMom » (santé maternelle par SMS), « Agro-Hub » (mise en relation producteurs/acheteurs) ; programme « Cameroon Digital Boost » MINPOSTEL.

\textbf{II. LE NUMÉRIQUE COMME FACTEUR D'ALIÉNATION ET DE PRÉCARISATION (Antithèse)}
A. Économie de l'attention et cybercriminalité juvénile
— Argument : Modèles économiques basés sur la captation (algorithmes), dérive vers « feyman » (arnaque), perte temps scolaire.
— Ex. : Étude UNESCO 2022 : 4h/jour moyen écrans lycéens Yaoundé ; arrestations « brouteurs » Douala 2023 (police judiciaire).
B. Fracture numérique et précarisation du travail de clic
— Argument : Inégalité d'accès (électricité, connexion, équipement), travail invisible (modération contenu, data labeling) sous-payé, sans protection sociale.
— Ex. : Jeunes « annotateurs » pour IA payés 50 FCFA/heure (enquête « Jeune Afrique » 2024) ; zones blanches Est/Adamaoua.

\textbf{III. VERS UN NUMÉRIQUE SOUVERAIN ET ÉTHIQUE : CONDITIONS D'UNE CHANCE RÉELLE (Synthèse)}
A. Politiques publiques : infrastructure, régulation, éducation aux médias
— Argument : Investissement backbone fibre, taxe numérique équitable, programme « École numérique » réel, éducation critique (fake news, données perso).
— Ex. : Loi 2023 sur la cybersécurité ; formation 5000 enseignants TICE (MINEDUB/MINESEC) ; modèle Rwanda « Smart Africa ».
B. Jeunesse comme actrice : communautés open source, communs numériques, éthique par le design
— Argument : Sortir de la simple consommation, contribuer au code, aux données, à la gouvernance (multipartisme numérique).
— Ex. : Communautés « Python Cameroon », « Django Girls Yaoundé », hackathons « Civic Tech » (budget participatif numérique Douala).

\textbf{CONCLUSION}
Bilan : Le numérique n'est ni chance ni menace par essence : il est un \emph{amplificateur} des dynamiques sociales. Sans régulation publique forte et appropriation critique par la jeunesse, il creuse les inégalités ; avec elles, il devient levier de souveraineté.
Ouverture : « L'outil façonne l'homme, mais l'homme doit façonner l'outil » — le défi camerounais est de former une « génération code » qui ne soit pas seulement une « génération data » pour les géants de la Silicon Valley.
\tcblower
\textbf{Commentaire du correcteur (simulé) :} Plan dialectique parfaitement assumé. Idées directrices nettes, arguments étayés, exemples locaux datés et variés (chiffres, startups, lois, enquêtes). Transitions I$\to$II (opposition) et II$\to$III (dépassement) rédigées. Introduction/conclusion aux standards. Note estimée : 17/20.
\end{exoresolu}

\courssub{Synthèse de la leçon : du plan à la copie — le fil conducteur}

\begin{definition}[Fil conducteur de la dissertation technique]
La dissertation n'est pas une suite de paragraphes : c'est une \textbf{architecture argumentative} dont le plan détaillé est le \textbf{plan de construction}. Chaque pierre (idée directrice) repose sur du ciment (argumentation) et porte une marque distinctive (exemple). Les transitions sont les poutres qui lient les étages. L'introduction est la façade, la conclusion le toit qui ouvre sur le ciel. Respecter cette architecture, c'est garantir la solidité de l'édifice — et la note au Probatoire.
\end{definition}

\begin{attention}[Derniers pièges à éviter le jour J]
\begin{itemize}
\item \textbf{Rédiger le plan détaillé \emph{après} la dissertation} (inversion du processus) : le plan doit guider l'écriture, pas la justifier a posteriori.
\item \textbf{Changer de plan en cours de route} : si une idée nouvelle surgit, l'intégrer \textbf{dans le plan détaillé} avant de l'écrire, sinon la cohérence se brise.
\item \textbf{Négliger les transitions} : elles valent 4 à 6 points sur 20 (cohérence, articulation). Une transition manquante = rupture logique = perte de points.
\item \textbf{Confondre ouverture et conclusion morale} : « Il faut que l'État fasse... » n'est pas une ouverture. L'ouverture est intellectuelle, pas prescriptive.
\end{itemize}
\end{attention}

\begin{savaistu}[L'examen pratique : le « devoir de dissertation » au Probatoire technique]
Durée : 4 heures. Sujet unique (souvent en lien avec les thèmes du programme : jeunesse, travail, tradition/modernité, technologie, citoyenneté). Notation : 20 points répartis typiquement — Introduction (4) / Développement : fond (8) + forme (4) / Conclusion (4). Le plan détaillé n'est pas noté directement mais \textbf{conditionne} les 16 points du développement et de la conclusion. Les correcteurs rapportent que les copies « planifiées » (plan détaillé visible au brouillon) gagnent en moyenne 3 à 4 points sur les copies « improvisées ». Investir 30-40 minutes au brouillon pour un plan détaillé solide est le meilleur retour sur investissement temps/points de l'épreuve.
\end{savaistu}

\coursec{Compte rendu, remédiation et entraînement au Bac}

Dans la section précédente, nous avons appris à rechercher des idées et à élaborer un plan détaillé de dissertation. Mais un travail de production ne s'achève jamais à la première version : il faut encore \cle{rendre compte} de ce qui a été fait, \cle{corriger} ses erreurs et s'\cle{entraîner} en conditions réelles. C'est l'objet de cette dernière section, qui clôt l'unité et vous prépare directement à l'épreuve écrite du Probatoire et du Baccalauréat technique.

\courssub{Le compte rendu : définition et objectifs}

\begin{definition}[Le compte rendu]
Le \cle{compte rendu} est un texte bref qui présente de manière \textbf{objective} le contenu d'une œuvre, d'une réunion, d'un cours ou d'une activité, puis qui en propose une appréciation. En réception des textes, le compte rendu de lecture permet de prouver qu'on a lu et compris une œuvre, sans la recopier ni la résumer intégralement.
\end{definition}

L'adjectif \emph{objectif} est essentiel : dans un compte rendu, on ne dit pas d'abord « j'ai aimé » ou « je n'ai pas aimé ». On commence par \textbf{présenter les faits} (l'auteur, l'œuvre, le contenu), et le jugement personnel n'intervient qu'à la fin, de manière argumentée. Le compte rendu est donc un exercice qui mobilise à la fois la \cle{contraction de texte} (pour la partie résumé) et l'\cle{argumentation} (pour l'appréciation).

\begin{methode}[Rédiger un compte rendu de lecture en trois parties]
\begin{enumerate}
\item \textbf{La présentation} : indiquer le titre de l'œuvre, le nom de l'auteur, le genre littéraire (roman, conte, pièce de théâtre), la date de publication et le contexte (pays, époque). Exemple : « \emph{Le Lion et la perle} est une pièce de théâtre du dramaturge nigérian Wole Soyinka, prix Nobel de littérature 1986, publiée en 1963. »
\item \textbf{Le résumé} : raconter l'essentiel de l'intrigue en quelques phrases, au présent de narration, sans détail inutile et sans citer de longs passages. On répond aux questions : qui ? quoi ? où ? comment se termine l'histoire ?
\item \textbf{L'appréciation} : donner son jugement argumenté sur l'œuvre (intérêt du sujet, force des personnages, qualité de la langue, leçon à retenir), en s'appuyant sur des éléments précis du texte.
\end{enumerate}
\end{methode}

\begin{exemplebox}[Compte rendu de lecture : \emph{Le Lion et la perle} de Wole Soyinka]
\textbf{Présentation} : \emph{Le Lion et la perle} est une pièce de théâtre de l'écrivain nigérian Wole Soyinka, parue en 1963. Elle met en scène un village africain confronté à la modernité.

\textbf{Résumé} : Sidi, la plus belle jeune fille du village d'Ilujinle — « la perle » —, est courtisée par le jeune instituteur Lakunle, partisan des coutumes modernes, qui refuse de payer la dot. Mais Baroka, le vieux chef du village — « le lion » —, use de ruse pour séduire Sidi : il fait courir le bruit de son impuissance, attire la jeune fille dans son piège et l'épouse finalement.

\textbf{Appréciation} : cette pièce séduit par son humour, la vivacité des dialogues et le débat qu'elle installe entre tradition et modernité. La ruse de Baroka montre que l'expérience et l'intelligence peuvent triompher de la jeunesse et des discours. C'est une œuvre courte, accessible, qui invite à réfléchir sur les transformations des sociétés africaines.
\end{exemplebox}

\begin{exercice}[Restitution écrite individuelle]
Rédigez individuellement, en vingt lignes maximum, le compte rendu de votre lecture du \emph{Lion et la perle} (ou du \emph{Vieux nègre et la médaille} de Ferdinand Oyono), en respectant la structure en trois parties : présentation, résumé, appréciation.
\end{exercice}

\begin{corrige}[Exercice : restitution écrite]
Critères de correction : la présentation mentionne l'auteur, le titre, le genre et la date (4 points) ; le résumé est bref, fidèle, rédigé au présent de narration, sans recopie (8 points) ; l'appréciation est un jugement personnel justifié par au moins deux arguments tirés de l'œuvre (6 points) ; la langue est correcte (2 points). Toute partie manquante entraîne la perte des points correspondants.
\end{corrige}

\courssub{La remédiation : identifier et corriger ses erreurs}

\begin{definition}[La remédiation]
La \cle{remédiation} est l'ensemble des activités de correction et de renforcement qui suivent une évaluation : elle consiste à \textbf{identifier ses erreurs récurrentes}, à en comprendre la cause, puis à s'exercer de nouveau pour les éliminer. En dissertation, la remédiation porte sur trois points sensibles : l'\textbf{analyse du sujet}, le \textbf{lexique} et la construction du \textbf{plan détaillé}.
\end{definition}

La remédiation n'est pas une punition : c'est la démarche de tout sportif ou de tout artisan qui progresse. Après chaque dissertation notée, il faut classer ses erreurs pour savoir où travailler.

\begin{center}
\begin{tabularx}{\linewidth}{|l|X|X|}\hline
\rowcolor{popblueL} \textbf{Erreur fréquente} & \textbf{Cause probable} & \textbf{Remédiation conseillée} \\ \hline
Hors sujet & Sujet lu trop vite, mots-clés non définis & Refaire l'analyse du sujet : souligner les mots-clés, définir chaque terme, repérer les limites du sujet \\ \hline
Pas de problématique & Confusion entre sujet et problématique & Transformer le sujet en question : « qu'est-ce qui est en débat ici ? » \\ \hline
Plan déséquilibré ou non détaillé & Idées non classées avant la rédaction & Regrouper les idées en deux ou trois parties de poids comparable ; préciser arguments et exemples pour chaque sous-partie \\ \hline
Exemples vagues & Culture insuffisante, pas de fiches de lecture & Constituer un carnet d'exemples (œuvres lues, faits d'actualité, proverbes) \\ \hline
Lexique pauvre ou répétitif & Synonymes et champs lexicaux non mobilisés & Utiliser la synonymie, l'antonymie, l'hypéronymie pour varier le vocabulaire \\ \hline
Ambiguïtés référentielles & Pronoms sans référent clair & À la relecture, vérifier que chaque « il », « elle », « cela » renvoie à un référent identifiable \\ \hline
\end{tabularx}
\end{center}

\begin{attention}[Ne négligez jamais la relecture]
L'erreur la plus coûteuse au Bac est de rendre la copie sans la relire. Or la relecture permet de corriger les \textbf{fautes d'orthographe et de grammaire}, mais aussi les \textbf{ambiguïtés référentielles} : un pronom dont le référent est incertain (« il » renvoyant à deux personnages possibles) rend la phrase incompréhensible et fait perdre des points de cohérence. Réservez toujours les dix dernières minutes à une relecture complète, phrase par phrase.
\end{attention}

\courssub{Du plan simple au plan détaillé : rappel méthodologique}

\begin{methode}[Construire un plan détaillé en quatre étapes]
\begin{enumerate}
\item \textbf{Regrouper les idées} par affinité thématique (idées qui vont ensemble).
\item \textbf{Ordonner les groupes} selon la logique du type de plan (dialectique, thématique, analytique) : thèse / antithèse / synthèse OU causes / conséquences / solutions, etc.
\item \textbf{Nommer chaque grande partie} par un titre court, clair, qui annonce l'argument principal.
\item \textbf{Détailler chaque partie} en deux ou trois sous-parties, chacune portant \textbf{un argument précis} appuyé d'\textbf{un exemple concret} (œuvre étudiée, fait de société, proverbe, statistique).
\end{enumerate}
\end{methode}

\begin{exemplebox}[Plan simple vs plan détaillé — Sujet : « La tradition est un obstacle au progrès. »]
\textbf{Plan simple (insuffisant au Bac)} :
\begin{itemize}
\item I. La tradition freine le progrès.
\item II. La tradition aide le progrès.
\item III. Il faut concilier les deux.
\end{itemize}

\textbf{Plan détaillé (attendu)} :
\begin{itemize}
\item \textbf{I. La tradition peut constituer un frein au développement}
  \begin{itemize}
  \item A. Certaines coutumes bloquent l'accès à l'éducation (ex. : mariages précoces interrompent la scolarité des filles au Nord-Cameroun).
  \item B. Le refus des techniques nouvelles maintient la pauvreté (ex. : Baroka dans \emph{Le Lion et la perle} s'oppose d'abord à la route et à l'école).
  \end{itemize}
\item \textbf{II. La tradition est aussi un levier pour un progrès enraciné}
  \begin{itemize}
  \item A. Elle transmet des valeurs stabilisatrices : solidarité, respect des aînés (ex. : la palabre africaine comme mode de règlement des conflits).
  \item B. Les savoirs endogènes inspirent des innovations (ex. : pharmacopée traditionnelle et recherche médicale moderne).
  \end{itemize}
\item \textbf{III. Vers un tri critique : garder ce qui élève, réformer ce qui opprime}
  \begin{itemize}
  \item A. Critère de sélection : la dignité humaine et le bien commun.
  \item B. Exemple : réforme des successions coutumières pour protéger les veuves et les filles.
  \end{itemize}
\end{itemize}
\end{exemplebox}

\courssub{Grille d'auto-évaluation}

Avant de rendre tout plan ou toute dissertation, interrogez-vous honnêtement à l'aide de la grille suivante. Chaque case cochée est un point gagné ; chaque case vide désigne un point de remédiation.

\begin{popfigure}
\begin{tikzpicture}[scale=1]
\draw[thick,popblue] (0,0) rectangle (12,5.2);
\foreach \y in {0.65,1.3,1.95,2.6,3.25,3.9,4.55} {\draw[popline] (0,\y) -- (12,\y);}
\draw[popline] (1.1,0) -- (1.1,5.2);
\fill[popblueL] (0,4.55) rectangle (12,5.2);
\node[font=\bfseries] at (6.5,4.87) {Grille d'auto-évaluation de ma dissertation};
\foreach \y/\txt in {3.9/{Ai-je défini tous les mots-clés du sujet ?},3.25/{Ai-je formulé une problématique claire (une question) ?},2.6/{Mon plan comporte-t-il 2 ou 3 parties équilibrées ?},1.95/{Chaque partie a-t-elle des arguments et des exemples précis ?},1.3/{Mon lexique est-il varié (synonymes, champs lexicaux) ?},0.65/{Ai-je prévu une introduction et une conclusion ?}} {
\draw[thick,popdark] (0.25,\y-0.18) rectangle (0.65,\y+0.22);
\node[anchor=west,font=\small] at (1.3,\y) {\txt};
}
\end{tikzpicture}
\legende{Grille d'auto-évaluation à cocher avant de rendre la copie : six critères essentiels, de l'analyse du sujet à la conclusion.}
\end{popfigure}

\courssub{Entraînement au Bac : sujet complet en conditions d'examen}

\begin{exoresolu}[Sujet d'entraînement en temps limité]
\textbf{Énoncé} : « La tradition est un obstacle au progrès. » Discutez cette opinion en vous appuyant sur des arguments et des exemples précis. (Temps conseillé : 50 minutes pour l'analyse et le plan détaillé.)
\tcblower
\textbf{Solution détaillée.}

\emph{1. Analyse du sujet (environ 20 minutes).} Mots-clés : \textbf{tradition} (ensemble des coutumes, croyances et pratiques transmises de génération en génération), \textbf{obstacle} (ce qui empêche d'avancer), \textbf{progrès} (amélioration des conditions de vie, avancées techniques et sociales). Le verbe « discutez » exige un plan dialectique : il faut examiner la thèse, puis la nuancer ou la combattre.

\emph{2. Problématique.} La tradition freine-t-elle nécessairement le développement des sociétés, ou peut-elle au contraire accompagner et enrichir le progrès ?

\emph{3. Plan détaillé (environ 30 minutes).}
\begin{itemize}
\item \textbf{Thèse} : la tradition peut être un obstacle au progrès. Argument 1 : certaines coutumes bloquent l'éducation (mariages précoces qui interrompent la scolarité des filles). Argument 2 : le refus des techniques nouvelles maintient la pauvreté. Exemple : dans \emph{Le Lion et la perle}, Baroka s'oppose d'abord à la modernité incarnée par Lakunle.
\item \textbf{Antithèse} : la tradition est aussi une richesse et un guide. Argument 1 : elle transmet des valeurs (solidarité, respect des anciens) qui stabilisent la société. Argument 2 : les savoirs traditionnels (médecine, agriculture) peuvent inspirer le progrès. Exemple : la palabre africaine comme mode de règlement des conflits.
\item \textbf{Synthèse} : il faut trier dans la tradition — conserver ce qui élève l'homme, abandonner ce qui l'opprime — pour construire un progrès enraciné dans la culture.
\end{itemize}
\end{exoresolu}

\courssub{Gérer son temps le jour de l'épreuve}

\begin{exemplebox}[Barème indicatif au Bac : le plan pèse lourd]
Dans les grilles de correction usuelles de la dissertation, la \textbf{pertinence du plan} (respect du type de plan attendu, équilibre des parties, niveau de détail) et la \textbf{cohérence} (enchaînement logique, clarté référentielle) représentent à elles deux une part majeure de la note — souvent près de la moitié des points —, devant la langue et la présentation. Un plan solide mais une langue moyenne rapporte donc plus qu'une belle langue sur un plan confus.
\end{exemplebox}

Pour une épreuve de quatre heures, la répartition conseillée du temps est la suivante : analyse du sujet (20 minutes), élaboration du plan détaillé au brouillon (30 minutes), rédaction de l'introduction et du développement (environ 2 h 30), rédaction de la conclusion (20 minutes), relecture (20 minutes).

\begin{popfigure}
\begin{tikzpicture}[scale=1.1]
% Camembert : analyse 20min, plan 30min, redaction 150min, conclusion 20min, relecture 20min (total 240)
\fill[popblue] (0,0) -- (0:2) arc (0:30:2) -- cycle;
\fill[poporange] (0,0) -- (30:2) arc (30:75:2) -- cycle;
\fill[popgreen] (0,0) -- (75:2) arc (75:300:2) -- cycle;
\fill[poppurple] (0,0) -- (300:2) arc (300:330:2) -- cycle;
\fill[popgold] (0,0) -- (330:2) arc (330:360:2) -- cycle;
\node[font=\small\bfseries] at (15:2.7) {Analyse : 20 min};
\node[font=\small\bfseries] at (52:2.7) {Plan : 30 min};
\node[font=\small\bfseries,align=center] at (187:2.9) {Rédaction : 2 h 30};
\node[font=\small\bfseries,align=center] at (315:2.8) {Conclusion : 20 min};
\node[font=\small\bfseries] at (345:2.75) {Relecture : 20 min};
\end{tikzpicture}
\legende{Répartition conseillée du temps pour une épreuve de dissertation de 4 heures : ne jamais sacrifier l'analyse, le plan ni la relecture.}
\end{popfigure}

\begin{aretenir}[L'essentiel de la section]
\begin{itemize}
\item Le \cle{compte rendu} de lecture comporte trois parties : présentation, résumé, appréciation argumentée ; il doit rester objectif.
\item La \cle{remédiation} consiste à repérer ses erreurs récurrentes (analyse du sujet, problématique, plan, lexique, références) et à s'exercer pour les corriger.
\item Le \cle{plan détaillé} se distingue du plan simple par la précision des arguments et des exemples pour chaque sous-partie.
\item La grille d'auto-évaluation permet de vérifier, avant de rendre la copie, que la problématique existe, que le plan est équilibré et détaillé, et que les exemples sont précis.
\item Le jour du Bac, gérez votre temps : 20 minutes d'analyse, 30 minutes de plan, l'essentiel pour la rédaction, et toujours 20 minutes de relecture finale.
\end{itemize}
\end{aretenir}

\begin{savaistu}[La dissertation au Baccalauréat technique camerounais]
À l'épreuve de français du Probatoire et du Baccalauréat technique, la dissertation porte le plus souvent sur des sujets de réflexion liés à la vie sociale : l'éducation, la tradition et la modernité, le travail, la jeunesse, les nouvelles technologies. Les correcteurs valorisent les copies où les exemples sont tirés de la littérature africaine étudiée en classe (Oyono, Soyinka, Beti...) et de la réalité camerounaise : c'est la preuve que l'élève sait relier l'école à la vie.
\end{savaistu}

\newpage

\coursec{Activité d'intégration}

\coursec{Leçon 07 : Dissertation — Produire le plan détaillé d'une dissertation}

\courssub{Objectifs de l'activité}

À l'issue de cet atelier d'intégration, chaque élève doit être capable de :

\begin{itemize}
\item analyser un sujet de dissertation (dégager les mots-clés, identifier le contenu latent, déterminer le type de sujet) ;
\item formuler une problématique claire à partir du sujet ;
\item mobiliser les relations lexicales (synonymie, antonymie, hypéronymie/hyponymie) pour construire le champ lexical du sujet ;
\item rechercher et sélectionner des idées et des exemples pertinents (œuvres au programme, réalité camerounaise) ;
\item élaborer un plan détaillé complet (introduction, développement en trois parties, conclusion) ;
\item confronter, évaluer et améliorer une production à l'aide d'une grille, puis rédiger un compte rendu justifiant ses choix.
\end{itemize}

\begin{aretenir}[Compétence visée]
Appliquer les notions de l'unité à une situation concrète d'examen et justifier une démarche de production d'écrit, conformément aux savoir-faire du programme de Terminale technique (Probatoire / Baccalauréat technique).
\end{aretenir}

\courssub{Situation déclenchante}

Le lycée technique de Bafoussam organise, chaque mois de mars, une « Semaine des lettres et des métiers ». Cette année, le club de lecture, que préside Aïcha, élève de Terminale Gestion, a été invité par le proviseur à animer le grand débat de clôture devant les classes de Première et de Terminale. Le thème retenu s'inspire de la lecture de \emph{Le Lion et la perle} de Wole Soyinka et du \emph{Vieux nègre et la médaille} de Ferdinand Oyono, étudiés en classe : l'affrontement entre les coutumes villageoises et les modes de vie modernes.

Pour préparer ce débat, le professeur de français remet à chaque groupe de quatre élèves le sujet suivant, tiré d'une annale du Baccalauréat technique :

\begin{exemplebox}[Sujet de dissertation proposé]
\textbf{« Faut-il préserver les traditions ou s'ouvrir à la modernité ? »}

Vous développerez votre point de vue sur cette question en vous appuyant sur des arguments et des exemples précis tirés de vos lectures, notamment des œuvres au programme, et de votre expérience de la société camerounaise.
\end{exemplebox}

Le club dispose de données recueillies lors d'un mini-sondage réalisé auprès de 120 élèves du lycée : 55 élèves déclarent que les traditions « freinent le progrès », 40 estiment qu'elles « protègent notre identité », et 25 pensent qu'« il faut concilier les deux ». Ces chiffres pourront nourrir l'argumentation.

\begin{popfigure}
\begin{tikzpicture}
\begin{axis}[
    ybar,
    bar width=1.2cm,
    width=\linewidth,
    height=7cm,
    symbolic x coords={Freinent le progrès, Protègent l'identité, Concilier les deux},
    xtick=data,
    x tick label style={rotate=30, anchor=east, font=\small},
    ylabel={Nombre d'élèves},
    ymin=0, ymax=65,
    ymajorgrids=true,
    grid style=dashed,
    nodes near coords,
    nodes near coords align={vertical},
    every node near coord/.append style={font=\small, popdark},
    title style={font=\normalsize\bfseries, text=popdark},
    title={Résultats du mini-sondage (120 élèves)}
]
\addplot[fill=popblue!70, draw=popblue] coordinates {(Freinent le progrès,55) (Protègent l'identité,40) (Concilier les deux,25)};
\end{axis}
\end{tikzpicture}
\legende{Figure 1 : Répartition des avis des élèves sur le rapport tradition/modernité (sondage interne, 120 répondants).}
\end{popfigure}

Chaque groupe doit produire, sur une fiche-modèle fournie, un \cle{plan détaillé} complet de la dissertation : introduction (amorce, présentation du sujet, problématique, annonce du plan), trois parties développées avec arguments et exemples, et conclusion. Les fiches seront ensuite affichées au mur de la salle, évaluées par les pairs à l'aide de la grille du club, améliorées, puis chaque groupe rédigera un compte rendu d'une quinzaine de lignes justifiant ses choix.

\courssub{Consignes}

Travail par groupes de quatre. Répondez dans l'ordre aux tâches suivantes.

\begin{enumerate}
\item \textbf{Analyse du sujet.} Soulignez les mots-clés du sujet. Précisez le sens de chacun (que recouvrent « traditions » et « modernité » ?). Identifiez le contenu latent : que suppose la question posée ? De quel type de sujet s'agit-il (sujet de réflexion, de discussion, de débat d'idées) ?
\item \textbf{Problématique.} Formulez en une ou deux phrases interrogatives la problématique qui orientera toute la dissertation.
\item \textbf{Champ lexical.} Construisez un tableau à trois colonnes donnant, pour « tradition » et « modernité » : deux synonymes ou expressions voisines, deux antonymes, un hyperonyme et deux hyponymes.
\item \textbf{Recherche des idées.} Par remue-méninges, notez au moins six idées (trois pour la préservation des traditions, trois pour l'ouverture à la modernité) et associez à chacune un exemple précis : personnages ou situations de \emph{Le Lion et la perle} (Baroka, Sidi, Lakunle) et du \emph{Vieux nègre et la médaille} (Meka), ou faits de la vie camerounaise (sondage du lycée, langues locales, nouvelles technologies, rites traditionnels).
\item \textbf{Choix du plan.} Justifiez le choix d'un plan dialectique (thèse -- antithèse -- synthèse) pour ce sujet.
\item \textbf{Plan détaillé.} Rédigez sur la fiche-modèle le plan détaillé complet : introduction rédigée en quelques lignes, trois parties avec pour chacune un titre, deux arguments et deux exemples, conclusion avec ouverture.
\item \textbf{Confrontation.} Affichez votre fiche, évaluez celle d'un autre groupe à l'aide de la grille (pertinence de la problématique sur 4 points, cohérence du plan sur 4 points, qualité des exemples sur 4 points, correction de la langue sur 4 points, présentation sur 4 points), puis améliorez votre propre production.
\item \textbf{Compte rendu.} Rédigez un compte rendu d'environ quinze lignes présentant vos choix de plan et les modifications apportées après la confrontation.
\end{enumerate}

\courssub{Ressources à mobiliser}

\begin{definition}[Le référent et ses substituts]
Dans un texte, le \cle{référent} est l'être, l'objet ou le concept désigné. Ses \cle{substituts} (pronoms, synonymes, périphrases, ellipses) évitent la répétition et assurent la cohésion. Exemple : \emph{Meka} (référent) $\rightarrow$ \emph{le vieux nègre, le récipiendaire, il, ce dernier} (substituts). Maîtriser cette chaîne référentielle est indispensable pour fluidifier l'argumentation dans la dissertation.
\end{definition}

\begin{definition}[Contenus manifestes et contenus latents]
Le \cle{contenu manifeste} est ce qui est explicitement énoncé dans l'énoncé (ici : « Faut-il préserver les traditions ou s'ouvrir à la modernité ? »). Le \cle{contenu latent} est ce qui est sous-entendu, présupposé : l'opposition binaire, l'exclusion mutuelle suggérée par la conjonction « ou », l'idée qu'on ne peut pas faire les deux à la fois. Le dégager permet de formuler une problématique qui dépasse l'alternative simpliste.
\end{definition}

\begin{aretenir}[Rappels utiles de la leçon]
\begin{itemize}
\item \textbf{Analyse du sujet} : repérer les mots porteurs de sens, dégager le contenu manifeste et le contenu latent.
\item \textbf{Problématique} : une question centrale, née de la tension entre les termes du sujet, à laquelle tout le devoir répond.
\item \textbf{Relations lexicales} : la \cle{synonymie} (mots de sens proche : coutume/tradition), l'\cle{antonymie} (mots de sens contraire : traditionnel/moderne), l'\cle{hypéronymie} (terme général : « valeur » est hyperonyme de « solidarité ») et l'hyponymie (terme spécifique).
\item \textbf{Texte argumentatif} : une thèse défendue par des arguments étayés d'exemples, organisée selon des connecteurs logiques (d'abord, en effet, cependant, par conséquent).
\item \textbf{Plan dialectique} : I. Thèse (il faut préserver les traditions) ; II. Antithèse (mais il faut s'ouvrir à la modernité) ; III. Synthèse (concilier les deux en dépassant l'opposition).
\item \textbf{Structure du plan détaillé} : introduction (amorce, sujet, problématique, annonce du plan), développement (parties, arguments, exemples), conclusion (bilan, réponse à la problématique, ouverture).
\end{itemize}
\end{aretenir}

\courssub{Démarche et corrigé commenté}

\begin{exoresolu}[Production attendue : du sujet au plan détaillé]
Énoncé : à partir du sujet « Faut-il préserver les traditions ou s'ouvrir à la modernité ? », produire l'analyse du sujet, la problématique, le champ lexical, la liste des idées et le plan détaillé complet, puis le compte rendu de groupe.

\tcblower

\textbf{Étape 1 : analyse du sujet.} Les mots-clés sont « préserver », « traditions », « s'ouvrir », « modernité ». Les traditions désignent les coutumes, les rites, les langues et les valeurs transmises par les ancêtres ; la modernité renvoie aux nouvelles technologies, à l'éducation occidentale, aux modes de vie contemporains. Le contenu latent du sujet est l'idée que tradition et modernité s'excluraient mutuellement : le « ou » invite à choisir. Il s'agit d'un sujet de débat d'idées appelant un plan dialectique.

\textbf{Étape 2 : problématique.} « Les traditions et la modernité sont-elles irréconciliables, ou la société camerounaise peut-elle préserver son identité tout en s'ouvrant au progrès ? »

\textbf{Étape 3 : champ lexical.}

\begin{center}
\begin{tabularx}{\linewidth}{|l|X|X|}\hline
\rowcolor{popblueL} \textbf{Relation} & \textbf{Tradition} & \textbf{Modernité} \\ \hline
Synonymes & coutume, héritage & progrès, innovation \\ \hline
Antonymes & modernité, nouveauté & tradition, archaïsme \\ \hline
Hyperonyme & valeur culturelle & évolution sociale \\ \hline
Hyponymes & dot, rite initiatique, langue locale & Internet, téléphonie mobile, ville moderne \\ \hline
\end{tabularx}
\end{center}

\textbf{Étape 4 : recherche des idées et des exemples.}
\begin{itemize}
\item \textbf{Pour la préservation des traditions :} elles fondent l'identité et la dignité (exemple : Meka, dans \emph{Le Vieux nègre et la médaille}, reste attaché à sa terre et à sa case malgré l'humiliation) ; elles assurent la cohésion sociale et la sagesse collective (exemple : la sagesse de Baroka, gardien des coutumes d'Ilujinle) ; le sondage le confirme : 40 élèves sur 120, soit environ 33 \%, estiment qu'elles « protègent notre identité ».
\item \textbf{Pour l'ouverture à la modernité :} elle apporte le progrès matériel, sanitaire et éducatif (exemple : les projets d'école et de routes de Lakunle pour Ilujinle ; la téléphonie mobile dans les villages camerounais) ; elle permet de questionner ou de dépasser certaines pratiques contestables (exemple : le refus du paiement de la dot par Lakunle) ; 55 élèves sur 120, soit environ 46 \%, jugent que les traditions « freinent le progrès ».
\item \textbf{Pour la synthèse (concilier les deux) :} la modernité sans racines est fragile et séductrice mais vide (exemple : l'échec du discours moderniste de Lakunle face à Sidi, qui choisit Baroka) ; la tradition intelligente sait s'approprier les outils du moderne (exemple : Baroka utilise la machine à timbrer et la photographie ; les radios communautaires camerounaises diffusant en langues locales) ; 25 élèves sur 120, soit environ 21 \%, souhaitent « concilier les deux ».
\end{itemize}

\textbf{Étape 5 : justification du plan.} Le sujet oppose deux termes par « ou » ; le plan dialectique (thèse, antithèse, synthèse) permet d'examiner les deux positions avant de les dépasser dans une vision nuancée : c'est le plan le plus adapté à un sujet de débat d'idées.

\textbf{Étape 6 : plan détaillé.}

\emph{Introduction.} Amorce : le Cameroun, comme toute l'Afrique, vit sous la double influence de ses racines ancestrales et du monde moderne. Sujet amené et posé. Problématique (étape 2). Annonce du plan : nous verrons d'abord qu'il faut préserver les traditions, ensuite qu'une ouverture à la modernité est nécessaire, enfin qu'il convient de concilier les deux dans une dynamique d'appropriation critique.

\emph{I. Il faut préserver les traditions, socle de l'identité et de la cohésion.}
\begin{itemize}
\item Argument 1 : Elles fondent l'identité collective et la dignité personnelle. Exemple : Meka (\emph{Le Vieux nègre et la médaille}) dont l'attachement à la terre et à la case résiste à l'aliénation coloniale.
\item Argument 2 : Elles garantissent la cohésion sociale par des valeurs partagées (solidarité, respect des aînés). Exemple : Baroka, le « Lion », incarne la sagesse coutumière à Ilujinle ; parallèlement, les rites initiatiques et les langues locales (bassa, bamileke, fulfulde, etc.) structurent le vivre-ensemble au Cameroun.
\end{itemize}

\emph{II. Mais l'ouverture à la modernité est indispensable pour le développement.}
\begin{itemize}
\item Argument 1 : Elle apporte le progrès technique, médical et éducatif. Exemple : les projets d'école et de dispensaire de Lakunle (\emph{Le Lion et la perle}) ; aujourd'hui, la téléphonie mobile et l'accès à Internet transforment l'économie rurale camerounaise (mobile banking, information agricole).
\item Argument 2 : Elle permet de réformer des pratiques devenues injustes ou obsolètes. Exemple : le refus de la dot par Lakunle, perçue comme une marchandisation de la femme ; le sondage du lycée : 46 \% des élèves jugent les traditions freinantes pour l'émancipation.
\end{itemize}

\emph{III. Il faut concilier tradition et modernité : une appropriation critique.}
\begin{itemize}
\item Argument 1 : La modernité déracinée échoue face à l'ancrage culturel. Exemple : l'échec de Lakunle, dont le modernisme superficiel ne séduit pas Sidi, attachée à sa valeur symbolique.
\item Argument 2 : La tradition vivante intègre les apports modernes sans s'y dissoudre. Exemple : Baroka utilise la machine à timbrer et la photo pour asseoir son pouvoir ; au Cameroun, les chefferies traditionnelles gèrent des coopératives agricoles modernes et les médias communautaires émettent en langues nationales.
\end{itemize}

\emph{Conclusion.} Bilan : ni repli identitaire, ni reniement mimétique. Réponse à la problématique : préserver l'essentiel des traditions (valeurs, langues, solidarité) tout en adoptant les apports utiles de la modernité (science, technique, droits humains). Ouverture : quelle place l'école camerounaise doit-elle faire aux langues nationales et aux savoirs endogènes dans les programmes officiels pour opérer cette synthèse ?

\textbf{Étape 7 : confrontation (auto-évaluation par les pairs).} À l'affichage, la grille d'évaluation (problématique 4 pts, cohérence du plan 4 pts, exemples 4 pts, langue 4 pts, présentation 4 pts) révèle par exemple qu'un groupe voisin a mieux hiérarchisé ses arguments (du plus général au plus spécifique) ou a utilisé des connecteurs plus variés. Le groupe déplace alors l'exemple du sondage dans l'argument qu'il illustre le mieux (II-2), reformule un titre de partie trop vague (ex. « La modernité » $\rightarrow$ « L'ouverture à la modernité, condition du développement ») et vérifie l'accord des participes passés.

\textbf{Étape 8 : compte rendu (modèle — environ 15 lignes).}
« Notre groupe a choisi le plan dialectique (thèse/antithèse/synthèse) car le sujet pose une alternative exclusive ("ou") qu'il faut dépasser. Nous avons d'abord placé l'exemple du choix de Sidi dans la première partie (identité), mais l'évaluation par les pairs a montré qu'il illustrait mieux la synthèse (rejet d'une modernité vide, acceptation d'un traditionnel modernisé) : nous l'avons déplacé en III-1. Nous avons également précisé notre problématique, initialement trop générale ("Tradition et modernité sont-elles compatibles ?"), en l'ancrant explicitement dans le contexte camerounais. Enfin, nous avons enrichi le champ lexical (hyperonymes/hyponymes) pour éviter les répétitions dans les titres de parties. »
\end{exoresolu}

\begin{attention}[Pièges fréquents à éviter]
\begin{itemize}
\item Confondre le plan détaillé avec la rédaction complète : le plan détaillé présente des arguments et des exemples formulés, mais pas des paragraphes rédigés.
\item Annoncer un plan dialectique puis produire seulement deux parties : la synthèse (troisième partie) est obligatoire pour dépasser l'opposition.
\item Citer des exemples sans les relier explicitement à l'argument : chaque exemple doit \emph{prouver} quelque chose (liaison argument/exemple).
\item Oublier de répondre explicitement à la problématique dans la conclusion (le « bilan » doit y répondre).
\item Négliger la correction linguistique (accords, syntaxe, connecteurs) même dans un plan : la grille d'évaluation y attribue 4 points.
\end{itemize}
\end{attention}

\courssub{Bilan de l'activité}

Cet atelier a permis de mettre en évidence trois acquis essentiels, directement transférables à l'épreuve du Probatoire ou du Baccalauréat technique. D'abord, la qualité d'une dissertation se joue avant la rédaction : une analyse rigoureuse du sujet, qui distingue contenu manifeste et contenu latent, conditionne une problématique pertinente et un plan cohérent. Ensuite, les outils de lexicologie (synonymie, antonymie, hypéronymie/hyponymie) et de cohésion textuelle (référent/substituts) ne sont pas de simples exercices scolaires : ils enrichissent le vocabulaire de l'argumentation, évitent les répétitions et assurent la fluidité du discours. Enfin, la confrontation des plans et l'évaluation par les pairs montrent qu'un plan détaillé n'est jamais définitif : il s'améliore par la justification, la remédiation et la prise en compte du regard d'autrui. Chaque élève dispose désormais d'une démarche complète, balisée et reproductible, de l'analyse du sujet jusqu'au compte rendu réflexif.

\newpage

\coursec{Méthodes et exercices résolus}

\coursec{Leçon 07 : Dissertation – Produire le plan détaillé d'une dissertation}

\courssub{Méthode 1 : Analyser un sujet de dissertation et formuler la problématique}

\begin{methode}[Analyser un sujet et dégager la problématique]
L'analyse du sujet est la première étape de la dissertation : elle conditionne tout le reste du travail. Voici la démarche en étapes.
\begin{enumerate}
\item \textbf{Lire le sujet plusieurs fois} et souligner les \cle{mots-clés} (les termes porteurs de sens).
\item \textbf{Définir chaque mot-clé} : donner son sens précis, ses synonymes, ses contraires. Attention aux mots polysémiques (qui ont plusieurs sens).
\item \textbf{Délimiter le sujet} : repérer ce qui est demandé (le champ du sujet) et ce qui est exclu (hors sujet).
\item \textbf{Dégager la tension} : un sujet de dissertation oppose souvent deux idées (par exemple : tradition / modernité, individu / société). C'est cette opposition qui fait problème.
\item \textbf{Formuler la problématique} : transformer la tension en une \cle{question problématique}, c'est-à-dire une question ouverte qui appelle une réflexion argumentée (commençant souvent par « Dans quelle mesure... », « Faut-il... », « Comment... »).
\item \textbf{Vérifier} : la problématique doit reprendre tous les mots-clés du sujet et permettre au moins deux réponses opposées.
\end{enumerate}
\end{methode}

\begin{exoresolu}[Analyse d'un sujet et formulation de la problématique]
\textbf{Énoncé (type Probatoire / Baccalauréat technique) :} « La littérature africaine doit-elle se contenter de divertir le lecteur ou doit-elle l'engager et l'éduquer ? » Analysez ce sujet et formulez la problématique.
\tcblower
\textbf{Solution détaillée.}

\textbf{Étape 1 : repérage des mots-clés.} Les termes porteurs de sens sont : \emph{littérature africaine}, \emph{divertir}, \emph{engager}, \emph{éduquer}.

\textbf{Étape 2 : définition des mots-clés.}
\begin{itemize}
\item \emph{Littérature africaine} : ensemble des œuvres écrites par des auteurs africains (romans, poésies, théâtre), souvent en lien avec l'histoire et les réalités du continent.
\item \emph{Divertir} : amuser, distraire, procurer du plaisir au lecteur.
\item \emph{Engager} : prendre position, pousser le lecteur à réfléchir et à agir sur les problèmes de sa société.
\item \emph{Éduquer} : instruire, former l'esprit et la conscience du lecteur.
\end{itemize}

\textbf{Étape 3 : délimitation.} Le sujet porte sur la \cle{fonction} de la littérature africaine, pas sur son histoire ni sur ses auteurs. Parler uniquement de la littérature européenne serait hors sujet.

\textbf{Étape 4 : tension du sujet.} Le sujet oppose deux conceptions : la littérature-plaisir (divertir) et la littérature-utilité (engager, éduquer). Le « ou » suggère un choix, mais on peut montrer que les deux fonctions se combinent.

\textbf{Étape 5 : problématique.} « Dans quelle mesure la littérature africaine peut-elle à la fois divertir le lecteur et remplir une mission d'engagement et d'éducation ? »

\textbf{Étape 6 : vérification.} La question reprend tous les mots-clés et permet deux réponses opposées (divertissement seul / engagement nécessaire), donc elle est bien problématique.

\textbf{Conclusion :} la problématique retenue est : « Dans quelle mesure la littérature africaine peut-elle à la fois divertir et éduquer le lecteur ? »
\end{exoresolu}

\courssub{Méthode 2 : Rechercher et organiser les idées}

\begin{methode}[Mobiliser ses idées et ses exemples]
\begin{enumerate}
\item \textbf{Faire un brainstorming} : noter au brouillon, sans trier, toutes les idées, exemples, œuvres, citations liés au sujet.
\item \textbf{Classer les idées} en deux ou trois colonnes selon la tension du sujet (pour / contre ; causes / conséquences ; etc.).
\item \textbf{Associer à chaque idée un exemple précis} : une œuvre étudiée en classe (par exemple \emph{Le Vieux nègre et la médaille} de Ferdinand Oyono), un fait de la vie courante, une situation camerounaise ou africaine.
\item \textbf{Éliminer} les idées hors sujet, répétitives ou impossibles à illustrer.
\item \textbf{Ordonner} : aller du plus simple au plus complexe, du plus faible au plus fort argument.
\item \textbf{Vérifier l'équilibre} : chaque grande partie doit contenir au moins deux arguments étayés par des exemples.
\end{enumerate}
\end{methode}

\begin{exoresolu}[Recherche et classement des idées]
\textbf{Énoncé :} Sujet : « L'école est-elle le seul lieu d'éducation de la jeunesse ? » Faites la recherche des idées et présentez le classement obtenu.
\tcblower
\textbf{Solution détaillée.}

\textbf{Étape 1 : brainstorming.} Idées notées : l'école donne des savoirs ; la famille éduque ; les médias et les réseaux sociaux influencent ; la rue et le quartier ; l'église ou la mosquée ; les amis ; les traditions et les contes ; les diplômes ; les valeurs morales ; le sport ; les maquis et les mauvaises fréquentations.

\textbf{Étape 2 : classement en colonnes.}
\begin{center}
\begin{tabularx}{\linewidth}{|X|X|}
\hline
\rowcolor{popblueL} \textbf{L'école, lieu central d'éducation} & \textbf{Les autres lieux d'éducation} \\
\hline
Elle transmet les savoirs et les diplômes & La famille transmet les valeurs morales et la langue \\
\hline
Elle apprend la vie en société (règles, ponctualité) & Les traditions (contes, initiations) forment l'identité \\
\hline
Elle ouvre l'esprit aux sciences et au monde & Les médias et réseaux sociaux informent (parfois mal) \\
\hline
\end{tabularx}
\end{center}

\textbf{Étape 3 : exemples associés.} Famille : l'enfant apprend le respect des aînés au village ; traditions : les contes du soir transmettent la sagesse ; école : l'apprentissage du calcul et du français au lycée technique.

\textbf{Étape 4 : élimination.} « Les maquis » est écarté : idée négative difficile à exploiter dans une dissertation scolaire.

\textbf{Étape 5 : ordonnancement.} On ira du rôle reconnu de l'école vers les autres instances éducatives, puis vers la complémentarité.

\textbf{Conclusion :} les idées sont classées en deux colonnes équilibrées, chacune illustrée d'exemples précis, prêtes pour l'élaboration du plan.
\end{exoresolu}

\courssub{Méthode 3 : Choisir le type de plan adapté}

\begin{methode}[Les trois grands types de plans]
\begin{enumerate}
\item \textbf{Le plan dialectique} (thèse / antithèse / synthèse) : on présente d'abord une position (I), puis la position opposée (II), puis un dépassement qui concilie les deux (III). Il convient aux sujets où deux opinions s'affrontent (« faut-il... ? », « pour ou contre... ? »).
\item \textbf{Le plan thématique} : on examine les différents aspects d'une notion (I : aspect 1 ; II : aspect 2 ; III : aspect 3). Il convient aux sujets descriptifs ou aux questions ouvertes sur un thème.
\item \textbf{Le plan analytique} (constat / causes / conséquences ou solutions) : on décrit un fait (I), on en cherche les causes (II), on en tire les conséquences ou les solutions (III). Il convient aux sujets portant sur un phénomène de société.
\end{enumerate}
\textbf{Réflexe :} le choix du plan dépend de la formulation du sujet et de la problématique, jamais du hasard.
\end{methode}

\begin{exoresolu}[Choisir le bon plan]
\textbf{Énoncé :} Pour chacun des sujets suivants, indiquez le type de plan le plus adapté et justifiez :
a) « Faut-il interdire les téléphones portables au lycée technique ? »
b) « Les fonctions du chef traditionnel dans le Cameroun d'aujourd'hui. »
c) « L'exode rural dans les pays africains. »
\tcblower
\textbf{Solution détaillée.}

\textbf{a) « Faut-il interdire les téléphones portables au lycée technique ? »} La formulation « Faut-il... ? » appelle deux réponses opposées (oui / non). On choisit donc le \cle{plan dialectique} : I. Les arguments pour l'interdiction (thèse) ; II. Les arguments contre (antithèse) ; III. Une solution équilibrée, par exemple un usage encadré (synthèse).

\textbf{b) « Les fonctions du chef traditionnel... »} Le sujet invite à inventorier les différents rôles d'une institution. On choisit le \cle{plan thématique} : I. Le chef, garant des traditions et de la cohésion sociale ; II. Le chef, auxiliaire de l'administration ; III. Les limites et les défis de son rôle aujourd'hui.

\textbf{c) « L'exode rural... »} Le sujet désigne un phénomène de société à expliquer. On choisit le \cle{plan analytique} : I. Le constat (les jeunes quittent massivement les villages) ; II. Les causes (pauvreté, manque d'infrastructures, attrait de la ville) ; III. Les conséquences et les solutions (développement des campagnes).

\textbf{Conclusion :} sujet a) $\rightarrow$ plan dialectique ; sujet b) $\rightarrow$ plan thématique ; sujet c) $\rightarrow$ plan analytique. Le choix du plan découle toujours de la formulation du sujet.
\end{exoresolu}

\courssub{Méthode 4 : Rédiger le plan détaillé}

\begin{methode}[Construire un plan détaillé complet]
Un \cle{plan détaillé} n'est pas une simple liste de titres : c'est la maquette de la dissertation entière.
\begin{enumerate}
\item \textbf{Introduction} : noter l'amorce (accroche), la présentation du sujet, la problématique et l'annonce du plan.
\item \textbf{Développement} : pour chaque grande partie (I, II, III), rédiger un titre complet (une phrase ou une expression claire), puis indiquer les sous-parties (A, B) avec, pour chacune, l'argument et l'exemple précis qui l'illustre.
\item \textbf{Transitions} : prévoir une phrase de liaison entre chaque grande partie.
\item \textbf{Conclusion} : noter le bilan (réponse à la problématique) et l'ouverture (élargissement vers une question voisine).
\item \textbf{Vérifier la progression} : les parties doivent se succéder logiquement, sans répétition, et répondre ensemble à la problématique.
\end{enumerate}
\end{methode}

\begin{exoresolu}[Production d'un plan détaillé]
\textbf{Énoncé (type Probatoire / Baccalauréat technique) :} Sujet : « La lecture est-elle encore utile à l'heure des réseaux sociaux ? » Rédigez le plan détaillé d'une dissertation sur ce sujet.
\tcblower
\textbf{Solution détaillée.}

\textbf{Analyse préalable :} mots-clés : \emph{lecture}, \emph{utile}, \emph{réseaux sociaux}. Tension : lecture (ancien média) contre réseaux sociaux (nouveaux médias). Problématique : « Dans quelle mesure la lecture conserve-t-elle son utilité face à l'essor des réseaux sociaux ? » Plan dialectique.

\textbf{Introduction :} amorce (les jeunes passent des heures sur leurs téléphones) ; présentation du sujet ; problématique ; annonce du plan.

\textbf{I. La lecture semble perdre son utilité face aux réseaux sociaux (thèse)}
\begin{itemize}
\item A. Les réseaux sociaux informent plus vite que les livres. \emph{Exemple :} un élève apprend les actualités en quelques secondes sur son téléphone.
\item B. Les réseaux divertissent davantage et captent le temps de lecture. \emph{Exemple :} beaucoup de jeunes préfèrent les vidéos courtes aux romans.
\end{itemize}
\emph{Transition :} Pourtant, la rapidité ne signifie pas la profondeur.

\textbf{II. La lecture garde des vertus irremplaçables (antithèse)}
\begin{itemize}
\item A. Elle développe l'esprit critique et le vocabulaire. \emph{Exemple :} la lecture de romans comme \emph{Le Vieux nègre et la médaille} enrichit la langue et fait réfléchir sur la société coloniale.
\item B. Elle forme la conscience et l'imagination, contrairement aux informations superficielles des réseaux. \emph{Exemple :} un lecteur de contes africains découvre la sagesse de sa culture.
\end{itemize}
\emph{Transition :} Faut-il alors opposer lecture et réseaux sociaux ?

\textbf{III. Lecture et réseaux sociaux peuvent se compléter (synthèse)}
\begin{itemize}
\item A. Les réseaux peuvent promouvoir la lecture. \emph{Exemple :} des clubs de lecture en ligne recommandent des livres.
\item B. L'essentiel est d'éduquer les jeunes à un usage équilibré des deux. \emph{Exemple :} un temps d'écran limité et un temps de lecture quotidien.
\end{itemize}

\textbf{Conclusion :} bilan (la lecture reste utile car elle forme l'esprit, à condition de cohabiter avec les réseaux) ; ouverture (comment l'école peut-elle redonner le goût de lire ?).
\end{exoresolu}

\courssub{Méthode 5 : Utiliser le lexique de l'argumentation (synonymie, antonymie, hypéronymie)}

\begin{methode}[Enrichir et nuancer son expression]
\begin{enumerate}
\item \textbf{La synonymie} : utiliser des mots de sens proche pour éviter les répétitions (ex. \emph{instruire / éduquer / former}). Attention : les synonymes sont rarement parfaits ; vérifier le sens exact.
\item \textbf{L'antonymie} : opposer des contraires pour construire la tension argumentative (ex. \emph{tradition / modernité}, \emph{ignorance / savoir}). C'est un outil précieux pour bâtir thèse et antithèse.
\item \textbf{L'hypéronymie et l'hyponymie} : l'hypéronyme est le terme général (ex. \emph{écrit}), l'hyponyme le terme particulier (ex. \emph{roman, poème, conte}). Passer du général au particulier permet d'illustrer ; passer du particulier au général permet de conclure.
\item \textbf{Réflexe de rédaction} : dans chaque paragraphe, varier le vocabulaire (synonymes), marquer les oppositions (antonymes) et articuler général et particulier (hypéronymes / hyponymes).
\end{enumerate}
\end{methode}

\begin{exoresolu}[Lexique de l'argumentation]
\textbf{Énoncé :} a) Donnez deux synonymes de « former ». b) Donnez l'antonyme de « ignorance » et de « liberté ». c) Classez les mots suivants du plus général au plus particulier : \emph{roman, œuvre, littérature}. d) Expliquez comment ces outils servent la dissertation.
\tcblower
\textbf{Solution détaillée.}

\textbf{a) Synonymes de « former » :} \emph{instruire}, \emph{éduquer}. Ces mots ont un sens proche mais pas identique : « instruire » insiste sur les savoirs, « éduquer » sur les valeurs et les comportements.

\textbf{b) Antonymes :} l'antonyme de « ignorance » est \emph{savoir} (ou \emph{connaissance}) ; l'antonyme de « liberté » est \emph{servitude} (ou \emph{esclavage}, \emph{oppression}).

\textbf{c) Classement du général au particulier :} \emph{littérature} (hypéronyme le plus général) $\rightarrow$ \emph{œuvre} $\rightarrow$ \emph{roman} (hyponyme le plus particulier). En effet, la littérature regroupe toutes les œuvres, et le roman est un type d'œuvre parmi d'autres (poème, pièce de théâtre, conte).

\textbf{d) Utilité pour la dissertation :}
\begin{itemize}
\item Les synonymes évitent les répétitions et montrent la richesse du vocabulaire, ce qui valorise la copie.
\item Les antonymes permettent de construire l'opposition thèse / antithèse : par exemple opposer « l'ignorance qui enchaîne » au « savoir qui libère ».
\item L'hypéronymie permet d'organiser les exemples : on annonce le terme général (« les œuvres engagées ») puis on cite les hyponymes précis (\emph{Le Vieux nègre et la médaille}, \emph{Une si longue lettre}).
\end{itemize}

\textbf{Conclusion :} maîtriser synonymie, antonymie et hypéronymie, c'est disposer des outils lexicaux indispensables pour argumenter avec précision et sans répétition.
\end{exoresolu}

\courssub{Méthode 6 : Confronter et améliorer sa production (compte rendu et remédiation)}

\begin{methode}[Évaluer et corriger un plan détaillé]
\begin{enumerate}
\item \textbf{Relire le plan} en posant quatre questions : le sujet est-il respecté ? La problématique est-elle claire ? Les parties répondent-elles à la problématique ? Chaque argument est-il illustré ?
\item \textbf{Confronter} sa production à celle d'un camarade : comparer les choix de plan, les arguments et les exemples.
\item \textbf{Repérer les défauts} : hors sujet, parties déséquilibrées, arguments sans exemples, répétitions, titres vagues.
\item \textbf{Corriger} : reformuler les titres, ajouter les exemples manquants, supprimer les passages hors sujet, rééquilibrer les parties.
\item \textbf{Rédiger le compte rendu} : présenter objectivement ce qui a été retenu, ce qui a été modifié et pourquoi.
\end{enumerate}
\end{methode}

\begin{exoresolu}[Remédiation d'un plan défaillant]
\textbf{Énoncé :} Voici le plan d'un élève sur le sujet « Faut-il préserver les langues locales au Cameroun ? » : I. Les langues locales sont nombreuses. II. Le français est la langue de l'école. III. J'aime ma langue. Relevez les défauts de ce plan et proposez un plan corrigé.
\tcblower
\textbf{Solution détaillée.}

\textbf{Relevé des défauts :}
\begin{enumerate}
\item \textbf{Absence de problématique} : le plan n'annonce aucune tension ; le sujet est une question (« faut-il... ? ») qui exige un débat.
\item \textbf{Titres vagues et non argumentatifs} : « Les langues locales sont nombreuses » est un constat, pas un argument ; « J'aime ma langue » est un avis personnel sans justification.
\item \textbf{Pas de type de plan identifiable} : ni dialectique, ni thématique, ni analytique ; les parties ne s'enchaînent pas logiquement.
\item \textbf{Aucun exemple annoncé} : rien n'indique comment les arguments seront illustrés.
\end{enumerate}

\textbf{Plan corrigé (dialectique).} Problématique : « Dans quelle mesure la préservation des langues locales est-elle nécessaire au Cameroun ? »

\textbf{I. Les langues locales semblent menacées et peu utiles (thèse)}
\begin{itemize}
\item A. Le français domine l'école et l'administration. \emph{Exemple :} les examens officiels se passent en français ou en anglais.
\item B. Certaines langues locales perdent leurs locuteurs. \emph{Exemple :} des jeunes citadins ne parlent plus la langue de leurs grands-parents.
\end{itemize}

\textbf{II. Pourtant, les langues locales sont indispensables (antithèse)}
\begin{itemize}
\item A. Elles portent l'identité et la culture. \emph{Exemple :} les contes et proverbes en ewondo, en duala ou en fulfuldé transmettent la sagesse des ancêtres.
\item B. Elles favorisent la communication dans la vie quotidienne. \emph{Exemple :} au marché et au village, la langue locale unit les générations.
\end{itemize}

\textbf{III. Il faut concilier langues officielles et langues locales (synthèse)}
\begin{itemize}
\item A. L'école peut intégrer les langues locales. \emph{Exemple :} l'enseignement des langues nationales prévu au Cameroun.
\item B. Chacun peut valoriser sa langue tout en maîtrisant le français. \emph{Exemple :} un élève bilingue est un atout pour son pays.
\end{itemize}

\textbf{Conclusion :} le plan corrigé repose sur une problématique claire, un plan dialectique équilibré et des exemples précis : il est désormais prêt à être rédigé.
\end{exoresolu}

\begin{attention}[Les 5 erreurs qui coûtent des points au Probatoire / Baccalauréat technique]
\begin{enumerate}
\item \textbf{Traiter à côté du sujet (hors sujet)} : ne pas définir les mots-clés avant de rédiger conduit à répondre à une autre question. Réflexe : souligner et définir chaque terme du sujet au brouillon.
\item \textbf{Oublier la problématique} : une introduction sans question problématique est une introduction ratée. La problématique doit reprendre tous les mots-clés et annoncer le plan.
\item \textbf{Confondre plan détaillé et rédaction} : le plan détaillé se compose de titres, d'arguments et d'exemples notés, pas de paragraphes rédigés. Inversement, une simple liste de mots sans arguments ni exemples n'est pas un plan détaillé.
\item \textbf{Affirmer sans illustrer} : chaque argument doit être suivi d'un exemple précis (œuvre étudiée, fait de société, situation camerounaise). Un argument sans exemple est une affirmation gratuite, donc non notée.
\item \textbf{Déséquilibrer les parties ou se contredire} : les grandes parties doivent avoir une importance comparable et la conclusion doit découler du développement. Une conclusion qui dit le contraire du corps du devoir révèle un plan mal construit : il faut alors refaire le plan, pas maquiller la fin.
\end{enumerate}
\end{attention}

\newpage

\coursec{Exercices}

\courssub{Je vérifie mes connaissances}

\begin{exercice}[QCM : les notions clés de l'unité]
Pour chaque question, entoure la (ou les) bonne(s) réponse(s).

\textbf{1.} Le \cle{paratexte} d'une œuvre comprend :
\begin{enumerate}[label=\alph*.]
\item le titre de l'œuvre ;
\item le nom de l'auteur et la couverture ;
\item la quatrième de couverture ;
\item le dernier chapitre du roman.
\end{enumerate}

\textbf{2.} Dans la phrase « Meka reçut sa médaille. \emph{Il} la contempla longuement », le pronom « il » est :
\begin{enumerate}[label=\alph*.]
\item un référent ;
\item un substitut du référent « Meka » ;
\item un champ lexical ;
\item un connecteur logique.
\end{enumerate}

\textbf{3.} Un \cle{hypéronyme} est :
\begin{enumerate}[label=\alph*.]
\item un mot de sens contraire ;
\item un mot de sens voisin ;
\item un terme générique qui englobe d'autres termes ;
\item un mot plus spécifique qu'un autre.
\end{enumerate}

\textbf{4.} Le plan le plus adapté à un sujet de réflexion du type « Faut-il... ? » est :
\begin{enumerate}[label=\alph*.]
\item le plan thématique ;
\item le plan dialectique (thèse / antithèse / synthèse) ;
\item le plan analytique (constat / causes / conséquences) ;
\item le récit chronologique.
\end{enumerate}
\end{exercice}

\begin{exercice}[Vrai ou faux ? Justifie ta réponse]
Réponds par VRAI ou FAUX, puis justifie en une ou deux phrases.
\begin{enumerate}
\item La problématique d'une dissertation se formule généralement sous la forme d'une ou plusieurs questions.
\item Deux synonymes ont toujours exactement le même sens et sont interchangeables dans tous les contextes.
\item Le contenu latent d'un énoncé est ce qui est dit explicitement, mot à mot.
\item Dans un plan détaillé, chaque partie doit être subdivisée en sous-parties avec des arguments et des exemples précis.
\item Dans \emph{Le Vieux nègre et la médaille} de Ferdinand Oyono, la médaille récompense réellement le mérite de Meka.
\end{enumerate}
\end{exercice}

\begin{exercice}[Définitions]
Définis en deux ou trois phrases chacune des notions suivantes, en donnant pour chacune un exemple :
\begin{enumerate}
\item le \cle{référent} et le \cle{substitut} ;
\item la \cle{synonymie} et l'\cle{antonymie} ;
\item l'\cle{hypéronymie} et l'\cle{hyponymie} ;
\item le \cle{contenu manifeste} et le \cle{contenu latent} ;
\item le \cle{plan détaillé} d'une dissertation.
\end{enumerate}
\end{exercice}

\begin{exercice}[Repérage rapide dans le paratexte]
On donne les éléments suivants, relevés sur un livre : « \emph{Le Vieux nègre et la médaille} — Ferdinand OYONO — Éditions Présence Africaine — Roman — 1956 ».
\begin{enumerate}
\item Identifie le titre, l'auteur, l'éditeur, le genre et la date de publication.
\item Quelles informations ce paratexte te donne-t-il avant même la lecture du texte ?
\item Quelles hypothèses de lecture peux-tu formuler à partir du seul titre ?
\end{enumerate}
\end{exercice}

\courssub{Je m'entraîne}

\begin{exercice}[Chaînes de référence]
Dans l'extrait suivant, inspiré de \emph{Le Vieux nègre et la médaille} : « Meka avait donné ses terres à la mission. Le vieil homme croyait que ce geste lui vaudrait la reconnaissance des Blancs. Le jour de la cérémonie, il attendit des heures sous le soleil. » :
\begin{enumerate}
\item Relève tous les substituts qui renvoient au référent « Meka ».
\item Réécris le passage en remplaçant « Meka » par « le vieux nègre » à chaque occurrence : quel effet produit cette répétition ?
\item Propose deux autres substituts possibles pour éviter les répétitions.
\end{enumerate}
\end{exercice}

\begin{exercice}[Synonymie et antonymie]
On donne les dix mots suivants : \emph{injustice, équité, bonheur, malheur, progrès, régression, liberté, servitude, générosité, équité}.
\begin{enumerate}
\item Classe ces mots en couples de synonymes.
\item Forme ensuite des couples d'antonymes.
\item Explique pourquoi « injustice » et « servitude » ne sont pas des synonymes parfaits.
\end{enumerate}
\end{exercice}

\begin{exercice}[Arbre hypéronymie / hyponymie]
À partir du mot « injustice », construis un arbre lexical :
\begin{enumerate}
\item Trouve un hypéronyme (terme plus général) et deux hyponymes (types d'injustice, par exemple : injustice sociale, injustice judiciaire).
\item Pour chaque hyponyme, propose un exemple concret tiré de la vie courante au Cameroun.
\item Utilise ce lexique pour rédiger une phrase argumentative sur le thème de l'injustice.
\end{enumerate}
\end{exercice}

\begin{exercice}[Contenus manifestes et contenus latents]
Voici trois affirmations inspirées de situations de \emph{Le Lion et la perle} de Wole Soyinka :
\begin{enumerate}
\item « Lakunle refuse de payer la dot : il dit défendre la modernité. »
\item « Baroka, le vieux chef, accueille l'étranger avec un large sourire. »
\item « Sidi pose fièrement devant le photographe de la ville. »
\end{enumerate}
Pour chaque affirmation :
\begin{enumerate}[label=\alph*.]
\item dégage le contenu manifeste (ce qui est dit) ;
\item dégage le contenu latent (ce qui est sous-entendu : ironie, critique, intention cachée).
\end{enumerate}
\end{exercice}

\courssub{Je me prépare au Bac}

\begin{exercice}[Sujet type Baccalauréat technique : dissertation]
\textbf{Sujet :} « Le progrès technique suffit-il à faire le bonheur de l'homme ? »

\begin{enumerate}
\item Analyse le sujet : dégage le ou les mots-clés, définis-les et repère la question posée.
\item Délimite le sujet : quels domaines peux-tu mobiliser (santé, communication, travail, agriculture au Cameroun...) ?
\item Formule la problématique sous forme de deux ou trois questions ordonnées.
\item Recherche les idées : note au moins trois arguments pour la thèse (le progrès technique rend heureux) et trois arguments pour l'antithèse (ses limites ou ses dangers), avec un exemple précis pour chaque argument.
\item Choisis le type de plan adapté et justifie ton choix en cinq lignes.
\item Rédige le plan détaillé complet : introduction (amorce, présentation du sujet, problématique, annonce du plan), trois parties subdivisées en sous-parties avec arguments et exemples, conclusion (bilan et ouverture).
\end{enumerate}
\end{exercice}

\begin{exercice}[Sujet type Probatoire : choix et justification d'un plan]
\textbf{Sujet :} « L'argent fait-il le bonheur ? »

Trois élèves proposent les plans suivants :
\begin{itemize}
\item \textbf{Élève A} — I. L'argent rend heureux. II. L'argent ne rend pas heureux. III. L'argent contribue au bonheur sans le garantir.
\item \textbf{Élève B} — I. Définition de l'argent. II. Histoire de la monnaie. III. Les monnaies modernes.
\item \textbf{Élève C} — I. Constat : l'argent occupe une place centrale dans nos sociétés. II. Causes : il permet de satisfaire les besoins essentiels. III. Conséquences : il ne suffit pas à combler les besoins affectifs et moraux.
\end{itemize}

\begin{enumerate}
\item Identifie le type de plan de chaque élève (dialectique, thématique, analytique).
\item Quel plan est hors sujet ? Justifie.
\item Entre les deux plans restants, lequel te semble le plus adapté à ce sujet de réflexion ? Justifie ton choix en cinq lignes.
\item Rédige la problématique du sujet, puis le plan détaillé correspondant au plan que tu as retenu (parties, sous-parties, arguments, exemples).
\end{enumerate}
\end{exercice}

\begin{exercice}[Remédiation : corriger un plan défaillant]
Voici le plan détaillé rédigé par un élève pour le sujet : « La lecture est-elle encore utile à l'ère du numérique ? »

\textbf{Plan de l'élève :}
\begin{itemize}
\item I. La lecture est très ancienne : elle existe depuis l'invention de l'écriture.
\item II. Les réseaux sociaux sont dangereux pour la jeunesse camerounaise.
\item III. Conclusion : la lecture est utile.
\end{itemize}

\begin{enumerate}
\item Releve au moins quatre erreurs dans ce plan (déséquilibre des parties, hors-sujet partiel, absence de problématique, sous-parties manquantes, conclusion réduite à une phrase).
\item Explique en quoi la partie II constitue un hors-sujet partiel : quelle question faudrait-il se poser pour la réintégrer au sujet ?
\item Formule une problématique correcte pour ce sujet.
\item Propose un plan détaillé corrigé : introduction complète, deux ou trois parties équilibrées avec sous-parties, arguments et exemples (tu pourras évoquer l'école, le téléphone portable, les bibliothèques), conclusion avec ouverture.
\end{enumerate}
\end{exercice}

\newpage

\coursec{Corrigés des exercices}

\begin{corrige}[Exercice 1 — QCM : les notions clés de l'unité]

\textbf{1.} Bonnes réponses : \textbf{a, b et c}. Le paratexte regroupe tout ce qui entoure le texte sans en faire partie : titre, nom de l'auteur, couverture, quatrième de couverture, préface, genre annoncé. La réponse \textbf{d} est fausse : le dernier chapitre appartient au texte lui-même, pas au paratexte.

\textbf{2.} Bonne réponse : \textbf{b}. Le pronom « il » reprend et remplace le nom « Meka » : c'est un \cle{substitut} (reprise anaphorique) du \cle{référent} « Meka ». Ce n'est ni un champ lexical (ensemble de mots liés à un même thème), ni un connecteur logique (mot de liaison comme « donc », « cependant »).

\textbf{3.} Bonne réponse : \textbf{c}. Un hypéronyme est un terme générique qui englobe des termes plus spécifiques (ses hyponymes). Exemple : « fruit » est l'hypéronyme de « mangue ». La réponse a définit l'antonyme, la b le synonyme, la d l'hyponyme.

\textbf{4.} Bonne réponse : \textbf{b}. Un sujet interrogatif de type « Faut-il... ? » appelle une confrontation d'opinions opposées : le plan \cle{dialectique} (thèse / antithèse / synthèse) est le plus adapté. Le plan analytique convient plutôt à un sujet de constat ou d'explication.

\textbf{Points de vigilance :} ne pas confondre paratexte et texte ; ne pas confondre hypéronyme (terme général) et hyponyme (terme spécifique) — piège classique au Bac.
\end{corrige}

\begin{corrige}[Exercice 2 — Vrai ou faux ? Justifie ta réponse]

\begin{enumerate}
\item \textbf{VRAI.} La problématique est la question centrale que le sujet pose ; elle se formule le plus souvent par une ou plusieurs questions ordonnées qui guideront la réflexion (ex. : « Le progrès technique rend-il toujours heureux ? N'a-t-il pas des limites ? »).

\item \textbf{FAUX.} Les synonymes ont des sens voisins mais rarement identiques : ils diffèrent par la nuance, le registre de langue ou l'emploi. « Demeure » et « maison » sont synonymes, mais « demeure » est soutenu ; on ne peut pas les intervertir dans tous les contextes.

\item \textbf{FAUX.} C'est le \cle{contenu manifeste} qui est dit explicitement. Le \cle{contenu latent} est ce qui est sous-entendu, suggéré, implicite (ironie, intention cachée, critique voilée).

\item \textbf{VRAI.} Un plan détaillé n'est pas une simple liste de titres : chaque partie est subdivisée en sous-parties, et chaque sous-partie comporte un argument développé et un exemple précis. C'est ce qui distingue le plan détaillé du plan simple.

\item \textbf{FAUX.} Dans le roman d'Oyono, la médaille ne récompense pas un mérite réel : elle est un instrument de la politique coloniale, destinée à remercier Meka d'avoir cédé ses terres à la mission et à donner bonne conscience aux colonisateurs. L'ironie du roman repose sur cette fausse récompense, suivie de l'humiliation de Meka.
\end{enumerate}
\end{corrige}

\begin{corrige}[Exercice 3 — Définitions]

\begin{enumerate}
\item \textbf{Référent et substitut.} Le \cle{référent} est l'être, l'objet ou la notion dont on parle dans un texte, désigné par une première mention. Le \cle{substitut} est tout mot ou groupe de mots qui reprend ce référent pour éviter la répétition : pronom, nom synonyme, périphrase. \emph{Exemple :} « Meka reçut sa médaille. Le vieil homme la contempla » — « Meka » est le référent, « le vieil homme » et « la » sont des substituts.

\item \textbf{Synonymie et antonymie.} La \cle{synonymie} est la relation de sens voisin entre deux mots (avec des nuances de registre ou d'intensité). L'\cle{antonymie} est la relation de sens contraire. \emph{Exemples :} « courageux » et « vaillant » sont synonymes ; « courageux » et « lâche » sont antonymes.

\item \textbf{Hypéronymie et hyponymie.} L'\cle{hypéronyme} est le terme générique qui englobe plusieurs termes spécifiques ; l'\cle{hyponyme} est chacun de ces termes spécifiques. \emph{Exemple :} « animal » est l'hypéronyme de « lion » et « chèvre », qui sont ses hyponymes.

\item \textbf{Contenu manifeste et contenu latent.} Le \cle{contenu manifeste} est ce que l'énoncé dit explicitement, à la lettre. Le \cle{contenu latent} est ce qu'il suggère sans le dire : sous-entendu, ironie, intention cachée. \emph{Exemple :} « Tu es enfin arrivé, il est seulement 18 heures... » dit manifestement une heure d'arrivée, mais sous-entend (latent) un reproche sur le retard.

\item \textbf{Plan détaillé.} Le \cle{plan détaillé} d'une dissertation est l'organisation complète et écrite du devoir : introduction (amorce, sujet, problématique, annonce du plan), développement en deux ou trois parties subdivisées en sous-parties avec arguments et exemples, conclusion (bilan et ouverture). Il constitue l'architecture du devoir avant la rédaction.
\end{enumerate}
\end{corrige}

\begin{corrige}[Exercice 4 — Repérage rapide dans le paratexte]

\begin{enumerate}
\item Identification des éléments :
\begin{center}
\begin{tabularx}{\linewidth}{|l|X|}\hline
\rowcolor{popblueL} \textbf{Élément} & \textbf{Information relevée} \\ \hline
Titre & \emph{Le Vieux nègre et la médaille} \\ \hline
Auteur & Ferdinand Oyono \\ \hline
Éditeur & Éditions Présence Africaine \\ \hline
Genre & Roman \\ \hline
Date de publication & 1956 \\ \hline
\end{tabularx}
\end{center}

\item Informations données avant la lecture : on sait qu'il s'agit d'une œuvre de fiction narrative (roman), écrite par un auteur camerounais, publiée par une maison d'édition engagée dans la valorisation des cultures africaines, à une époque (1956) où le Cameroun était encore sous domination coloniale. On peut donc s'attendre à un texte en lien avec la colonisation.

\item Hypothèses de lecture à partir du titre : le roman mettra en scène un vieil homme noir (« le vieux nègre ») ; la « médaille » évoque une récompense, une distinction honorifique ; l'association des deux termes laisse supposer une récompense accordée à un Africain, peut-être par l'administration coloniale, et invite à s'interroger sur sa valeur réelle (mérite sincère ou piège ?). Le titre annonce donc une possible ironie.
\end{enumerate}
\end{corrige}

\begin{corrige}[Exercice 5 — Chaînes de référence]

\begin{enumerate}
\item Substituts renvoyant au référent « Meka » : « \emph{ses} terres » (adjectif possessif), « \emph{Le vieil homme} » (périphrase nominale), « \emph{lui} vaudrait » (pronom complément), « \emph{il} attendit » (pronom sujet). Toute la chaîne de reprises maintient la cohérence du texte sans répéter le nom propre.

\item Réécriture : « Le vieux nègre avait donné ses terres à la mission. Le vieux nègre croyait que ce geste lui vaudrait la reconnaissance des Blancs. Le jour de la cérémonie, le vieux nègre attendit des heures sous le soleil. » \textbf{Effet produit :} la répétition alourdit le texte, crée une monotonie et une insistance appuyée ; elle peut produire un effet d'insistance volontaire (dénomination méprisante répétée, qui souligne le regard colonial), mais dans un devoir elle est le signe d'un vocabulaire pauvre.

\item Deux autres substituts possibles : « le vieillard », « le père Meka », « le bénéficiaire de la médaille », « le protagoniste ». Exemple : « Le vieillard croyait que ce geste lui vaudrait la reconnaissance des Blancs. »
\end{enumerate}

\textbf{Points de vigilance :} au Bac, la multiplication des substituts variés et précis est un critère de qualité d'expression ; éviter les reprises vagues qui rendent le référent ambigu.
\end{corrige}

\begin{corrige}[Exercice 6 — Synonymie et antonymie]

\begin{enumerate}
\item Couples de synonymes : dans la liste proposée, on peut rapprocher \emph{injustice} et \emph{iniquité} (sens voisin : absence d'équité), \emph{bonheur} et \emph{félicité}, \emph{progrès} et \emph{avancement}. Attention : deux synonymes ne sont jamais parfaitement interchangeables ; ils diffèrent par la nuance ou le registre (ainsi « félicité » est plus soutenu que « bonheur »).

\item Couples d'antonymes :
\begin{itemize}
\item \emph{injustice} $\leftrightarrow$ \emph{équité} (ou justice) ;
\item \emph{bonheur} $\leftrightarrow$ \emph{malheur} ;
\item \emph{progrès} $\leftrightarrow$ \emph{régression} ;
\item \emph{liberté} $\leftrightarrow$ \emph{servitude} ;
\item \emph{générosité} $\leftrightarrow$ \emph{avarice}.
\end{itemize}

\item « Injustice » et « servitude » ne sont pas des synonymes parfaits : l'\emph{injustice} désigne l'absence d'équité, le fait de violer le droit ; la \emph{servitude} désigne l'état de dépendance, d'asservissement d'une personne. Les deux notions se recoupent parfois (la servitude est une forme d'injustice), mais une injustice peut exister sans servitude (un procès inéquitable, par exemple). Ce sont des termes liés par le champ lexical du mal social, non des synonymes.
\end{enumerate}

\textbf{Points de vigilance :} au Bac, employer un synonyme pour éviter une répétition suppose de vérifier que la nuance convient au contexte ; un antonyme doit appartenir à la même classe grammaticale (nom contre nom, adjectif contre adjectif).
\end{corrige}

\begin{corrige}[Exercice 7 — Arbre hypéronymie / hyponymie]

\begin{enumerate}
\item Arbre lexical :
\begin{itemize}
\item \textbf{Hypéronyme} (terme général) : le \emph{mal social} (ou « le vice », « le tort »).
\item \textbf{Terme} : l'\emph{injustice}.
\item \textbf{Hyponymes} (types d'injustice) : l'\emph{injustice sociale}, l'\emph{injustice judiciaire} (on peut ajouter l'injustice économique, l'injustice scolaire).
\end{itemize}

\item Exemples concrets au Cameroun :
\begin{itemize}
\item \emph{Injustice sociale} : deux travailleurs accomplissant le même travail ne reçoivent pas le même salaire ; l'accès inégal aux soins de santé entre la ville et le village.
\item \emph{Injustice judiciaire} : un innocent condamné faute de moyens pour payer un avocat ; un puissant échappant à toute poursuite.
\end{itemize}

\item Phrase argumentative : « L'injustice, sous toutes ses formes — sociale, judiciaire ou économique — demeure l'un des maux sociaux les plus graves, car elle détruit la confiance des citoyens envers leurs institutions ; c'est pourquoi l'État camerounais doit garantir à chacun un procès équitable et un salaire juste. »
\end{enumerate}
\end{corrige}

\begin{corrige}[Exercice 8 — Contenus manifestes et contenus latents]

\begin{enumerate}
\item \textbf{Lakunle et la dot.}
\begin{itemize}
\item \emph{Contenu manifeste :} Lakunle refuse de payer la dot et justifie son refus par la défense de la modernité.
\item \emph{Contenu latent :} sa « modernité » peut cacher un prétexte égoïste (économiser l'argent de la dot) ou un mépris maladroit des traditions de son propre village ; l'énoncé invite à douter de la sincérité du personnage.
\end{itemize}

\item \textbf{Baroka et l'étranger.}
\begin{itemize}
\item \emph{Contenu manifeste :} le vieux chef accueille l'étranger en souriant.
\item \emph{Contenu latent :} le « large sourire » peut dissimuler la ruse : Baroka est un personnage calculateur qui ménage ses apparences pour mieux manipuler autrui ; l'hospitalité affichée cache une méfiance ou un dessein.
\end{itemize}

\item \textbf{Sidi et le photographe.}
\begin{itemize}
\item \emph{Contenu manifeste :} Sidi pose avec fierté devant le photographe venu de la ville.
\item \emph{Contenu latent :} cette fierté révèle la vanité du personnage et la séduction exercée par la modernité (la photographie, la ville) sur la jeunesse du village ; elle annonce le conflit entre tradition et modernité qui traverse la pièce.
\end{itemize}
\end{enumerate}

\textbf{Points de vigilance :} distinguer clairement ce qui est écrit (manifeste) de ce qui est interprété (latent) ; au commentaire comme à la dissertation, le contenu latent doit être justifié par des indices du texte, non inventé.
\end{corrige}

\begin{corrige}[Exercice 9 — Sujet type Baccalauréat technique : « Le progrès technique suffit-il à faire le bonheur de l'homme ? »]

\begin{enumerate}
\item \textbf{Analyse du sujet.} Mots-clés : « progrès technique » (ensemble des innovations scientifiques et technologiques qui transforment la vie : médecine, transports, téléphone, machines agricoles) ; « suffire » (être la condition unique et complète) ; « bonheur de l'homme » (état de satisfaction durable, bien-être matériel et moral). La question posée : le progrès technique est-il, à lui seul, la source du bonheur ?

\item \textbf{Délimitation.} Domaines mobilisables : la santé (vaccins, hôpitaux), la communication (téléphone portable, Internet), le travail (machines, gain de temps), l'agriculture au Cameroun (tracteurs, semences améliorées), mais aussi les dangers (chômage, dépendance aux écrans, pollution, isolement). Le sujet porte sur l'homme en général, pas seulement sur la technique.

\item \textbf{Problématique.} Le progrès technique améliore-t-il réellement les conditions de vie de l'homme ? Ne comporte-t-il pas des limites et des dangers qui empêchent le bonheur complet ? Dès lors, quelles autres conditions (morales, affectives, sociales) sont nécessaires au bonheur ?

\item \textbf{Recherche des idées.}
\begin{itemize}
\item \emph{Thèse (le progrès rend heureux) :} 1) il soulage la peine et prolonge la vie (ex. : les vaccins et les hôpitaux sauvent des vies au Cameroun) ; 2) il rapproche les hommes (ex. : le téléphone portable permet à un étudiant de Yaoundé de parler à sa famille au village) ; 3) il libère du temps et augmente les richesses (ex. : les machines agricoles améliorent les récoltes).
\item \emph{Antithèse (limites et dangers) :} 1) il crée la dépendance et l'isolement (ex. : les jeunes enfermés dans les réseaux sociaux négligent les relations réelles) ; 2) il détruit des emplois et creuse les inégalités (ex. : les machines remplacent les ouvriers) ; 3) il menace la santé et la nature (ex. : pollution, accidents de la route liés à la vitesse).
\end{itemize}

\item \textbf{Choix du plan.} Le sujet est une question fermée (« suffit-il... ? ») qui appelle une réponse nuancée : le \cle{plan dialectique} (thèse / antithèse / synthèse) est le plus adapté, car il permet d'examiner les bienfaits, puis les limites, puis de dépasser l'opposition en montrant que le bonheur exige davantage que la technique.

\item \textbf{Plan détaillé complet.}

\textbf{Introduction :} \emph{Amorce :} depuis l'invention de la machine à vapeur jusqu'au téléphone portable, la technique transforme la vie humaine. \emph{Présentation du sujet :} beaucoup voient dans le progrès technique la clé du bonheur. \emph{Problématique :} le progrès technique suffit-il à faire le bonheur de l'homme ? \emph{Annonce du plan :} nous verrons d'abord que le progrès technique améliore la vie humaine, puis qu'il comporte des limites et des dangers, enfin que le bonheur exige des valeurs que la technique ne peut offrir.

\textbf{I. Le progrès technique contribue au bonheur de l'homme.}
\begin{itemize}
\item \emph{a) Il préserve la santé et prolonge la vie.} Argument : la médecine moderne vainc des maladies autrefois mortelles. Exemple : les campagnes de vaccination au Cameroun.
\item \emph{b) Il facilite la communication et le travail.} Argument : il rapproche les hommes et allège les peines. Exemple : le téléphone portable, les machines agricoles.
\end{itemize}

\textbf{II. Mais le progrès technique ne suffit pas et comporte des dangers.}
\begin{itemize}
\item \emph{a) Il crée de nouvelles servitudes.} Argument : dépendance aux écrans, isolement. Exemple : des jeunes qui délaissent les études pour les réseaux sociaux.
\item \emph{b) Il engendre inégalités et destructions.} Argument : chômage technique, pollution. Exemple : les ouvriers remplacés par les machines ; les accidents de la circulation.
\end{itemize}

\textbf{III. Le bonheur repose sur des valeurs que la technique ne peut donner.}
\begin{itemize}
\item \emph{a) Le bonheur est d'abord affectif et moral.} Argument : l'amour, l'amitié, la paix intérieure ne s'achètent pas. Exemple : un malade riche et seul reste malheureux.
\item \emph{b) Le progrès n'est heureux que s'il est maîtrisé et partagé.} Argument : la technique doit servir l'homme et non l'asservir. Exemple : l'usage raisonné d'Internet pour étudier.
\end{itemize}

\textbf{Conclusion :} \emph{Bilan :} le progrès technique est une condition du bien-être, mais non du bonheur complet, qui exige des valeurs humaines. \emph{Ouverture :} reste à savoir comment éduquer les jeunes à un usage sage des nouvelles technologies.
\end{enumerate}

\textbf{Barème indicatif (sur 20) :} analyse du sujet et problématique : 4 pts ; recherche des idées (arguments + exemples) : 6 pts ; choix et justification du plan : 3 pts ; plan détaillé (structure, équilibre, sous-parties) : 5 pts ; correction de la langue : 2 pts.

\textbf{Points de vigilance :} la problématique doit reprendre les mots-clés du sujet ; chaque argument doit être suivi d'un exemple précis ; les trois parties doivent être équilibrées ; ne jamais annoncer un plan dialectique puis rédiger un plan thématique.
\end{corrige}

\begin{corrige}[Exercice 10 — Sujet type Probatoire : « L'argent fait-il le bonheur ? »]

\begin{enumerate}
\item \textbf{Types de plans.} Élève A : plan \cle{dialectique} (thèse / antithèse / synthèse). Élève B : plan \cle{thématique} (inventaire d'aspects autour d'une notion). Élève C : plan \cle{analytique} (constat / causes / conséquences).

\item \textbf{Plan hors sujet.} Le plan de l'élève B est hors sujet : il traite de l'argent en soi (définition, histoire, monnaies modernes) mais ne répond jamais à la question posée, qui porte sur le rapport entre l'argent et le \emph{bonheur}. C'est un exposé sur la monnaie, pas une dissertation sur le sujet.

\item \textbf{Plan le plus adapté.} Le plan dialectique de l'élève A est le plus adapté : le sujet est une question fermée (« fait-il... ? ») qui appelle une confrontation entre l'opinion favorable (l'argent procure confort et sécurité) et l'opinion contraire (l'argent ne donne ni amour ni santé), avant un dépassement (l'argent est un moyen, non une fin). Le plan analytique de C est acceptable mais moins naturel ici, car le sujet n'invite pas d'abord à un constat de fait mais à un débat d'idées.

\item \textbf{Problématique et plan détaillé (plan dialectique retenu).}

\emph{Problématique :} L'argent procure-t-il le bonheur en satisfaisant nos besoins ? Ne révèle-t-il pas aussi ses limites devant les biens qui ne s'achètent pas ? Comment alors faut-il concevoir le rôle de l'argent dans une vie heureuse ?

\textbf{I. L'argent semble faire le bonheur.}
\begin{itemize}
\item \emph{a) Il satisfait les besoins essentiels.} Argument : sans argent, on ne peut se nourrir, se loger, se soigner. Exemple : une famille qui peut payer les soins d'un enfant malade.
\item \emph{b) Il donne sécurité et considération sociale.} Argument : il protège de la misère et ouvre les portes de la réussite. Exemple : l'entrepreneur prospère respecté dans sa communauté.
\end{itemize}

\textbf{II. Pourtant, l'argent ne fait pas le bonheur.}
\begin{itemize}
\item \emph{a) Il n'achète ni l'amour, ni la santé, ni la paix du cœur.} Argument : les biens affectifs et moraux échappent à l'argent. Exemple : un homme riche abandonné par tous à l'hôpital.
\item \emph{b) Il peut même détruire le bonheur.} Argument : la soif de l'argent engendre corruption, crimes, familles déchirées. Exemple : des héritages qui divisent les familles camerounaises.
\end{itemize}

\textbf{III. L'argent contribue au bonheur sans le garantir.}
\begin{itemize}
\item \emph{a) L'argent est un moyen, non une fin.} Argument : il rend service lorsqu'il sert des projets humains. Exemple : financer les études de ses enfants.
\item \emph{b) Le bonheur naît de l'équilibre entre le matériel et le moral.} Argument : il faut joindre au confort les valeurs d'amour, de travail et de partage. Exemple : une famille modeste mais unie et solidaire.
\end{itemize}
\end{enumerate}

\textbf{Barème indicatif :} identification des plans : 3 pts ; repérage du hors-sujet justifié : 3 pts ; choix justifié (5 lignes) : 4 pts ; problématique : 3 pts ; plan détaillé : 7 pts.

\textbf{Points de vigilance :} un plan thématique n'est pas toujours faux, mais il devient hors sujet s'il oublie la question posée ; toujours vérifier que chaque partie répond à la problématique ; la synthèse ne doit pas être une simple répétition de la thèse.
\end{corrige}

\begin{corrige}[Exercice 11 — Remédiation : corriger un plan défaillant]

\begin{enumerate}
\item \textbf{Erreurs relevées (au moins quatre) :}
\begin{itemize}
\item absence d'introduction et de problématique : le plan commence directement par la partie I ;
\item déséquilibre des parties : la partie I est historique et hors proportion, la partie II est un hors-sujet partiel ;
\item absence de sous-parties : aucune partie n'est subdivisée en arguments et exemples ;
\item partie I peu argumentative : elle constate l'ancienneté de la lecture sans répondre à la question de son utilité actuelle ;
\item conclusion réduite à une phrase, sans bilan ni ouverture ;
\item le plan ne comporte que deux parties de développement, sans confrontation d'idées.
\end{itemize}

\item \textbf{Le hors-sujet partiel de la partie II.} « Les réseaux sociaux sont dangereux pour la jeunesse » traite des réseaux sociaux en général, non de la \emph{lecture} à l'ère du numérique. Pour la réintégrer au sujet, il faudrait se demander : « la concurrence des réseaux sociaux et des écrans détourne-t-elle les jeunes de la lecture ? » — autrement dit, lier le numérique à la pratique de lire.

\item \textbf{Problématique correcte.} À l'ère du numérique, la lecture conserve-t-elle son utilité pour l'élève et le citoyen ? Les écrans ne la rendent-ils pas secondaire ? Quelle place faut-il alors lui accorder dans la formation de l'homme ?

\item \textbf{Plan détaillé corrigé.}

\textbf{Introduction :} \emph{Amorce :} le téléphone portable et Internet occupent désormais une grande part du temps des jeunes. \emph{Présentation du sujet :} certains annoncent la fin de la lecture, jugée dépassée. \emph{Problématique :} la lecture est-elle encore utile à l'ère du numérique ? \emph{Annonce du plan :} nous montrerons d'abord que le numérique concurrence fortement la lecture, puis que la lecture demeure irremplaçable, enfin qu'il faut concilier lecture et outils numériques.

\textbf{I. À l'ère du numérique, la lecture semble menacée.}
\begin{itemize}
\item \emph{a) Les écrans captent le temps des jeunes.} Argument : réseaux sociaux, jeux et vidéos remplacent le livre. Exemple : des élèves qui passent leurs soirées sur le téléphone au lieu de lire.
\item \emph{b) L'information rapide semble suffire.} Argument : on croit tout savoir par un résumé en ligne. Exemple : les élèves qui se contentent de fiches trouvées sur Internet.
\end{itemize}

\textbf{II. Pourtant, la lecture reste irremplaçable.}
\begin{itemize}
\item \emph{a) Elle forme l'esprit et la langue.} Argument : lire enrichit le vocabulaire, l'orthographe et la réflexion, indispensables à la réussite scolaire. Exemple : les bons élèves en dissertation sont d'abord de bons lecteurs.
\item \emph{b) Elle développe l'esprit critique et l'imagination.} Argument : le livre apprend à analyser en profondeur, contrairement aux informations fragmentées des écrans. Exemple : la lecture de romans comme \emph{Le Vieux nègre et la médaille} fait comprendre l'histoire coloniale du Cameroun.
\end{itemize}

\textbf{III. Il faut concilier lecture et numérique.}
\begin{itemize}
\item \emph{a) Le numérique peut servir la lecture.} Argument : livres numériques, bibliothèques en ligne donnent accès à de nombreux textes. Exemple : télécharger gratuitement des classiques sur son téléphone.
\item \emph{b) L'école et la famille doivent entretenir le goût de lire.} Argument : bibliothèques scolaires, clubs de lecture, temps de lecture quotidien. Exemple : les clubs de lecture dans les lycées de Yaoundé et de Douala.
\end{itemize}

\textbf{Conclusion :} \emph{Bilan :} menacée par les écrans, la lecture demeure indispensable à la formation de l'esprit ; il faut donc l'allier au numérique au lieu de l'y sacrifier. \emph{Ouverture :} comment encourager la lecture dans les familles camerounaises où le livre reste coûteux ?
\end{enumerate}

\textbf{Points de vigilance :} un plan défaillant se corrige d'abord par la problématique ; chaque partie doit contenir au moins deux sous-parties ; vérifier que chaque partie répond à une question de la problématique ; la conclusion exige un bilan ET une ouverture.
\end{corrige}

\newpage

\coursec{Fiche bilan}

\coursec{Leçon 07 : Dissertation – Produire le plan détaillé}

\begin{popfigure}
\begin{center}\tcbox[colback=popink,colframe=popink,coltext=white,arc=9pt,boxsep=4pt,left=6pt,right=6pt]{\bfseries\small Dissertation : plan détaillé}\end{center}
\vspace{-2pt}
\begin{tcbraster}[raster columns=3,raster equal height=rows,raster column skip=6pt,raster row skip=6pt,enhanced,blankest]
\begin{tcolorbox}[enhanced,colback=white,colframe=poporange,boxrule=0.9pt,arc=3pt,title={Analyser le sujet},fonttitle=\bfseries\scriptsize,coltitle=white,colbacktitle=poporange,left=4pt,right=4pt,top=2pt,bottom=2pt,boxsep=1pt,toptitle=1.5pt,bottomtitle=1.5pt,halign=left]
\begin{itemize}[nosep,leftmargin=*,itemsep=1pt,font=\scriptsize]
\item Mots-clés
\item Problématique
\end{itemize}
\end{tcolorbox}
\begin{tcolorbox}[enhanced,colback=white,colframe=popgreen,boxrule=0.9pt,arc=3pt,title={Chercher les idées},fonttitle=\bfseries\scriptsize,coltitle=white,colbacktitle=popgreen,left=4pt,right=4pt,top=2pt,bottom=2pt,boxsep=1pt,toptitle=1.5pt,bottomtitle=1.5pt,halign=left]
\begin{itemize}[nosep,leftmargin=*,itemsep=1pt,font=\scriptsize]
\item Brainstorming
\item Classer, trier
\end{itemize}
\end{tcolorbox}
\begin{tcolorbox}[enhanced,colback=white,colframe=poppurple,boxrule=0.9pt,arc=3pt,title={Types de plan},fonttitle=\bfseries\scriptsize,coltitle=white,colbacktitle=poppurple,left=4pt,right=4pt,top=2pt,bottom=2pt,boxsep=1pt,toptitle=1.5pt,bottomtitle=1.5pt,halign=left]
\begin{itemize}[nosep,leftmargin=*,itemsep=1pt,font=\scriptsize]
\item Dialectique
\item Thématique
\item Analytique
\end{itemize}
\end{tcolorbox}
\begin{tcolorbox}[enhanced,colback=white,colframe=popgold,boxrule=0.9pt,arc=3pt,title={Lexique argumentatif},fonttitle=\bfseries\scriptsize,coltitle=white,colbacktitle=popgold,left=4pt,right=4pt,top=2pt,bottom=2pt,boxsep=1pt,toptitle=1.5pt,bottomtitle=1.5pt,halign=left]
\begin{itemize}[nosep,leftmargin=*,itemsep=1pt,font=\scriptsize]
\item Synonymie / antonymie
\item Hypéronymie / hyponymie
\end{itemize}
\end{tcolorbox}
\begin{tcolorbox}[enhanced,colback=white,colframe=poppink,boxrule=0.9pt,arc=3pt,title={Texte argumentatif},fonttitle=\bfseries\scriptsize,coltitle=white,colbacktitle=poppink,left=4pt,right=4pt,top=2pt,bottom=2pt,boxsep=1pt,toptitle=1.5pt,bottomtitle=1.5pt,halign=left]
\begin{itemize}[nosep,leftmargin=*,itemsep=1pt,font=\scriptsize]
\item Thèse, arguments, exemples
\item Manifeste / latent
\item Connecteurs logiques
\end{itemize}
\end{tcolorbox}
\begin{tcolorbox}[enhanced,colback=white,colframe=popblue,boxrule=0.9pt,arc=3pt,title={Référent et substituts},fonttitle=\bfseries\scriptsize,coltitle=white,colbacktitle=popblue,left=4pt,right=4pt,top=2pt,bottom=2pt,boxsep=1pt,toptitle=1.5pt,bottomtitle=1.5pt,halign=left]
\begin{itemize}[nosep,leftmargin=*,itemsep=1pt,font=\scriptsize]
\item Reprises nominales
\item Pronoms
\end{itemize}
\end{tcolorbox}
\end{tcbraster}

\legende{Carte mentale de la leçon : les six compétences clés pour produire le plan détaillé d'une dissertation au Baccalauréat technique.}
\end{popfigure}

\begin{bilan}
\textbf{Définitions clés}
\begin{itemize}
  \item \cle{Dissertation} : exercice écrit qui discute un sujet de manière ordonnée, objective et argumentée (introduction, développement en 2 ou 3 parties, conclusion).
  \item \cle{Problématique} : question centrale (souvent formulée en 2 ou 3 questions articulées) née de l'analyse du sujet ; elle guide tout le devoir et en assure la cohérence.
  \item \cle{Plan détaillé} : plan structuré avec titres des parties, sous-parties numérotées, arguments et exemples précis pour chaque idée directrice.
  \item \cle{Référent} : réalité désignée par un mot ; ses \cle{substituts} (pronoms personnels, relatifs, démonstratifs ; reprises nominales ; synonymes) évitent les répétitions et assurent la cohésion.
  \item \cle{Synonymie} / \cle{antonymie} : ressemblance / opposition de sens entre mots. \cle{Hypéronymie} : terme générique (\emph{fleur}) ; \cle{hyponymie} : terme spécifique (\emph{rose}, \emph{tulipe}).
  \item \cle{Contenu manifeste} : sens explicite, littéral du texte ; \cle{contenu latent} : sens implicite, sous-entendu, ironique ou idéologique.
\end{itemize}

\textbf{Les trois types de plan à connaître par c\oe ur}
\begin{center}
\begin{tabularx}{\linewidth}{|l|X|X|}
\hline
\rowcolor{popblueL} \textbf{Plan} & \textbf{Structure type} & \textbf{Usage privilégié} \\
\hline
Dialectique & Thèse / Antithèse / Synthèse & Sujets polémiques, questions ouvertes (\emph{« Faut-il... ? »}, \emph{« Est-il juste que... ? »}) \\
\hline
Thématique & Aspect 1 / Aspect 2 / Aspect 3 & Sujets descriptifs, explicatifs ou de définition (\emph{« Les causes de... »}, \emph{« Les formes de... »}) \\
\hline
Analytique & Constat / Causes / Conséquences / Solutions & Problèmes de société, sujets à dimension pratique (\emph{« Comment lutter contre... ? »}) \\
\hline
\end{tabularx}
\end{center}

\textbf{Méthode express : du sujet au plan détaillé}
\begin{enumerate}
  \item Lire le sujet 2 fois ; souligner les mots-clés et définir chacun (synonymes, champ lexical).
  \item Dégager le thème général, puis formuler la problématique (questions claires, liées au sujet).
  \item Chercher les idées par brainstorming (mots, arguments, exemples, citations), les trier et les regrouper par affinité.
  \item Choisir le type de plan adapté au sujet ; répartir les groupes d'idées en 2 ou 3 parties équilibrées.
  \item Détailler : chaque partie = sous-parties numérotées (1, 2, 3) + argument + exemple précis (fait, chiffre, citation, référence littéraire).
  \item Vérifier la cohérence : connecteurs logiques (\emph{d'abord, en effet, cependant, par conséquent}), substituts du référent, aucune idée hors sujet, progression logique.
\end{enumerate}

\textbf{Caractéristiques du texte argumentatif} : une thèse (soutenue ou rejetée), des arguments étayés par des exemples concrets, des connecteurs logiques (\emph{d'abord, en effet, cependant, donc, par ailleurs}), un lexique varié et précis (synonymes, antonymes, champs lexicaux, hypéronymes/hyponymes), distinction entre contenus manifestes et latents.
\end{bilan}

\begin{aretenir}[Checklist avant le Bac – Plan détaillé]
Avant de rendre votre plan détaillé, vérifiez chaque point :
\begin{itemize}
  \item $\square$ J'ai souligné et défini tous les mots-clés du sujet.
  \item $\square$ J'ai dégagé le thème sans déformer le sujet.
  \item $\square$ Ma problématique est formulée en questions claires et articulées.
  \item $\square$ J'ai choisi le type de plan adapté au sujet (dialectique, thématique, analytique).
  \item $\square$ Mon plan comporte 2 ou 3 parties équilibrées en volume.
  \item $\square$ Chaque partie contient des sous-parties numérotées (1, 2, 3...).
  \item $\square$ Chaque idée directrice est accompagnée d'un argument et d'un exemple précis.
  \item $\square$ J'ai évité les répétitions grâce aux substituts du référent (pronoms, reprises, synonymes).
  \item $\square$ J'ai utilisé des synonymes, antonymes et le champ lexical pertinent.
  \item $\square$ Mes idées sont reliées par des connecteurs logiques variés.
  \item $\square$ Je n'ai introduit aucune idée hors sujet.
  \item $\square$ Je me suis relu et j'ai corrigé l'orthographe, la grammaire, la ponctuation.
\end{itemize}
\end{aretenir}

\findecours

\end{document}
